% Options for packages loaded elsewhere
% Options for packages loaded elsewhere
\PassOptionsToPackage{unicode}{hyperref}
\PassOptionsToPackage{hyphens}{url}
%
\documentclass[
  a4paper,
]{scrbook}
\usepackage{xcolor}
\usepackage[margin=1in]{geometry}
\usepackage{amsmath,amssymb}
\setcounter{secnumdepth}{5}
\usepackage{iftex}
\ifPDFTeX
  \usepackage[T1]{fontenc}
  \usepackage[utf8]{inputenc}
  \usepackage{textcomp} % provide euro and other symbols
\else % if luatex or xetex
  \usepackage{unicode-math} % this also loads fontspec
  \defaultfontfeatures{Scale=MatchLowercase}
  \defaultfontfeatures[\rmfamily]{Ligatures=TeX,Scale=1}
\fi
\usepackage{lmodern}
\ifPDFTeX\else
  % xetex/luatex font selection
\fi
% Use upquote if available, for straight quotes in verbatim environments
\IfFileExists{upquote.sty}{\usepackage{upquote}}{}
\IfFileExists{microtype.sty}{% use microtype if available
  \usepackage[]{microtype}
  \UseMicrotypeSet[protrusion]{basicmath} % disable protrusion for tt fonts
}{}
\makeatletter
\@ifundefined{KOMAClassName}{% if non-KOMA class
  \IfFileExists{parskip.sty}{%
    \usepackage{parskip}
  }{% else
    \setlength{\parindent}{0pt}
    \setlength{\parskip}{6pt plus 2pt minus 1pt}}
}{% if KOMA class
  \KOMAoptions{parskip=half}}
\makeatother
% Make \paragraph and \subparagraph free-standing
\makeatletter
\ifx\paragraph\undefined\else
  \let\oldparagraph\paragraph
  \renewcommand{\paragraph}{
    \@ifstar
      \xxxParagraphStar
      \xxxParagraphNoStar
  }
  \newcommand{\xxxParagraphStar}[1]{\oldparagraph*{#1}\mbox{}}
  \newcommand{\xxxParagraphNoStar}[1]{\oldparagraph{#1}\mbox{}}
\fi
\ifx\subparagraph\undefined\else
  \let\oldsubparagraph\subparagraph
  \renewcommand{\subparagraph}{
    \@ifstar
      \xxxSubParagraphStar
      \xxxSubParagraphNoStar
  }
  \newcommand{\xxxSubParagraphStar}[1]{\oldsubparagraph*{#1}\mbox{}}
  \newcommand{\xxxSubParagraphNoStar}[1]{\oldsubparagraph{#1}\mbox{}}
\fi
\makeatother

\usepackage{color}
\usepackage{fancyvrb}
\newcommand{\VerbBar}{|}
\newcommand{\VERB}{\Verb[commandchars=\\\{\}]}
\DefineVerbatimEnvironment{Highlighting}{Verbatim}{commandchars=\\\{\}}
% Add ',fontsize=\small' for more characters per line
\usepackage{framed}
\definecolor{shadecolor}{RGB}{241,243,245}
\newenvironment{Shaded}{\begin{snugshade}}{\end{snugshade}}
\newcommand{\AlertTok}[1]{\textcolor[rgb]{0.68,0.00,0.00}{#1}}
\newcommand{\AnnotationTok}[1]{\textcolor[rgb]{0.37,0.37,0.37}{#1}}
\newcommand{\AttributeTok}[1]{\textcolor[rgb]{0.40,0.45,0.13}{#1}}
\newcommand{\BaseNTok}[1]{\textcolor[rgb]{0.68,0.00,0.00}{#1}}
\newcommand{\BuiltInTok}[1]{\textcolor[rgb]{0.00,0.23,0.31}{#1}}
\newcommand{\CharTok}[1]{\textcolor[rgb]{0.13,0.47,0.30}{#1}}
\newcommand{\CommentTok}[1]{\textcolor[rgb]{0.37,0.37,0.37}{#1}}
\newcommand{\CommentVarTok}[1]{\textcolor[rgb]{0.37,0.37,0.37}{\textit{#1}}}
\newcommand{\ConstantTok}[1]{\textcolor[rgb]{0.56,0.35,0.01}{#1}}
\newcommand{\ControlFlowTok}[1]{\textcolor[rgb]{0.00,0.23,0.31}{\textbf{#1}}}
\newcommand{\DataTypeTok}[1]{\textcolor[rgb]{0.68,0.00,0.00}{#1}}
\newcommand{\DecValTok}[1]{\textcolor[rgb]{0.68,0.00,0.00}{#1}}
\newcommand{\DocumentationTok}[1]{\textcolor[rgb]{0.37,0.37,0.37}{\textit{#1}}}
\newcommand{\ErrorTok}[1]{\textcolor[rgb]{0.68,0.00,0.00}{#1}}
\newcommand{\ExtensionTok}[1]{\textcolor[rgb]{0.00,0.23,0.31}{#1}}
\newcommand{\FloatTok}[1]{\textcolor[rgb]{0.68,0.00,0.00}{#1}}
\newcommand{\FunctionTok}[1]{\textcolor[rgb]{0.28,0.35,0.67}{#1}}
\newcommand{\ImportTok}[1]{\textcolor[rgb]{0.00,0.46,0.62}{#1}}
\newcommand{\InformationTok}[1]{\textcolor[rgb]{0.37,0.37,0.37}{#1}}
\newcommand{\KeywordTok}[1]{\textcolor[rgb]{0.00,0.23,0.31}{\textbf{#1}}}
\newcommand{\NormalTok}[1]{\textcolor[rgb]{0.00,0.23,0.31}{#1}}
\newcommand{\OperatorTok}[1]{\textcolor[rgb]{0.37,0.37,0.37}{#1}}
\newcommand{\OtherTok}[1]{\textcolor[rgb]{0.00,0.23,0.31}{#1}}
\newcommand{\PreprocessorTok}[1]{\textcolor[rgb]{0.68,0.00,0.00}{#1}}
\newcommand{\RegionMarkerTok}[1]{\textcolor[rgb]{0.00,0.23,0.31}{#1}}
\newcommand{\SpecialCharTok}[1]{\textcolor[rgb]{0.37,0.37,0.37}{#1}}
\newcommand{\SpecialStringTok}[1]{\textcolor[rgb]{0.13,0.47,0.30}{#1}}
\newcommand{\StringTok}[1]{\textcolor[rgb]{0.13,0.47,0.30}{#1}}
\newcommand{\VariableTok}[1]{\textcolor[rgb]{0.07,0.07,0.07}{#1}}
\newcommand{\VerbatimStringTok}[1]{\textcolor[rgb]{0.13,0.47,0.30}{#1}}
\newcommand{\WarningTok}[1]{\textcolor[rgb]{0.37,0.37,0.37}{\textit{#1}}}

\usepackage{longtable,booktabs,array}
\newcounter{none} % for unnumbered tables
\usepackage{calc} % for calculating minipage widths
% Correct order of tables after \paragraph or \subparagraph
\usepackage{etoolbox}
\makeatletter
\patchcmd\longtable{\par}{\if@noskipsec\mbox{}\fi\par}{}{}
\makeatother
% Allow footnotes in longtable head/foot
\IfFileExists{footnotehyper.sty}{\usepackage{footnotehyper}}{\usepackage{footnote}}
\makesavenoteenv{longtable}
\usepackage{graphicx}
\makeatletter
\newsavebox\pandoc@box
\newcommand*\pandocbounded[1]{% scales image to fit in text height/width
  \sbox\pandoc@box{#1}%
  \Gscale@div\@tempa{\textheight}{\dimexpr\ht\pandoc@box+\dp\pandoc@box\relax}%
  \Gscale@div\@tempb{\linewidth}{\wd\pandoc@box}%
  \ifdim\@tempb\p@<\@tempa\p@\let\@tempa\@tempb\fi% select the smaller of both
  \ifdim\@tempa\p@<\p@\scalebox{\@tempa}{\usebox\pandoc@box}%
  \else\usebox{\pandoc@box}%
  \fi%
}
% Set default figure placement to htbp
\def\fps@figure{htbp}
\makeatother


% definitions for citeproc citations
\NewDocumentCommand\citeproctext{}{}
\NewDocumentCommand\citeproc{mm}{%
  \begingroup\def\citeproctext{#2}\cite{#1}\endgroup}
\makeatletter
 % allow citations to break across lines
 \let\@cite@ofmt\@firstofone
 % avoid brackets around text for \cite:
 \def\@biblabel#1{}
 \def\@cite#1#2{{#1\if@tempswa , #2\fi}}
\makeatother
\newlength{\cslhangindent}
\setlength{\cslhangindent}{1.5em}
\newlength{\csllabelwidth}
\setlength{\csllabelwidth}{3em}
\newenvironment{CSLReferences}[2] % #1 hanging-indent, #2 entry-spacing
 {\begin{list}{}{%
  \setlength{\itemindent}{0pt}
  \setlength{\leftmargin}{0pt}
  \setlength{\parsep}{0pt}
  % turn on hanging indent if param 1 is 1
  \ifodd #1
   \setlength{\leftmargin}{\cslhangindent}
   \setlength{\itemindent}{-1\cslhangindent}
  \fi
  % set entry spacing
  \setlength{\itemsep}{#2\baselineskip}}}
 {\end{list}}
\usepackage{calc}
\newcommand{\CSLBlock}[1]{\hfill\break\parbox[t]{\linewidth}{\strut\ignorespaces#1\strut}}
\newcommand{\CSLLeftMargin}[1]{\parbox[t]{\csllabelwidth}{\strut#1\strut}}
\newcommand{\CSLRightInline}[1]{\parbox[t]{\linewidth - \csllabelwidth}{\strut#1\strut}}
\newcommand{\CSLIndent}[1]{\hspace{\cslhangindent}#1}



\setlength{\emergencystretch}{3em} % prevent overfull lines

\providecommand{\tightlist}{%
  \setlength{\itemsep}{0pt}\setlength{\parskip}{0pt}}



 


\makeatletter
\@ifpackageloaded{bookmark}{}{\usepackage{bookmark}}
\makeatother
\makeatletter
\@ifpackageloaded{caption}{}{\usepackage{caption}}
\AtBeginDocument{%
\ifdefined\contentsname
  \renewcommand*\contentsname{Table of contents}
\else
  \newcommand\contentsname{Table of contents}
\fi
\ifdefined\listfigurename
  \renewcommand*\listfigurename{List of Figures}
\else
  \newcommand\listfigurename{List of Figures}
\fi
\ifdefined\listtablename
  \renewcommand*\listtablename{List of Tables}
\else
  \newcommand\listtablename{List of Tables}
\fi
\ifdefined\figurename
  \renewcommand*\figurename{Figure}
\else
  \newcommand\figurename{Figure}
\fi
\ifdefined\tablename
  \renewcommand*\tablename{Table}
\else
  \newcommand\tablename{Table}
\fi
}
\@ifpackageloaded{float}{}{\usepackage{float}}
\floatstyle{ruled}
\@ifundefined{c@chapter}{\newfloat{codelisting}{h}{lop}}{\newfloat{codelisting}{h}{lop}[chapter]}
\floatname{codelisting}{Listing}
\newcommand*\listoflistings{\listof{codelisting}{List of Listings}}
\makeatother
\makeatletter
\makeatother
\makeatletter
\@ifpackageloaded{caption}{}{\usepackage{caption}}
\@ifpackageloaded{subcaption}{}{\usepackage{subcaption}}
\makeatother
\usepackage{bookmark}
\IfFileExists{xurl.sty}{\usepackage{xurl}}{} % add URL line breaks if available
\urlstyle{same}
% fallback for those not using the hyperref driver hyperxmp:
\makeatletter
\@ifundefined{xmpquote}{\newcommand{\xmpquote}[1]{#1}}{}
\makeatother
\hypersetup{
  pdftitle={Applied AI for Political Scientists},
  pdfauthor={Babak RezaeeDaryakenari},
  hidelinks,
  pdfcreator={LaTeX via pandoc}}


\title{Applied AI for Political Scientists}
\author{Babak RezaeeDaryakenari}
\date{2026-09-18}
\begin{document}
\frontmatter
\maketitle

\renewcommand*\contentsname{Table of contents}
{
\setcounter{tocdepth}{2}
\tableofcontents
}

\mainmatter
\bookmarksetup{startatroot}

\chapter*{Preface}\label{preface}
\addcontentsline{toc}{chapter}{Preface}

\markboth{Preface}{Preface}

These handouts are prepared for Applied AI for Political Scientists by
Dr.~Babak RezaeeDaryakenari
\href{mailto:s.rezaeedaryakenari@fsw.leidenuniv.nl}{\nolinkurl{s.rezaeedaryakenari@fsw.leidenuniv.nl}}.

Work through the \textbf{R Bootcamp} before the first session on 2
September. The \textbf{Weekly Handouts} accompany the sessions.

This is an ongoing project. Please do not share or use it without the
author's written permission.

\part{R Bootcamp}

\chapter*{Installing R and RStudio}\label{installing-r-and-rstudio}
\addcontentsline{toc}{chapter}{Installing R and RStudio}

\markboth{Installing R and RStudio}{Installing R and RStudio}

Two separate pieces of software are needed, and the distinction between
them causes more confusion than it should. \textbf{R} is the language:
the engine that does the computing. \textbf{RStudio} is the environment
you work in: the editor, the panes, the buttons. RStudio cannot do
anything without R behind it, which is why R goes first.

Everything here is free, and the whole process takes about fifteen
minutes on a normal connection. Do it before the first session on 2
September rather than during it.

\section*{Step 1: install R}\label{step-1-install-r}
\addcontentsline{toc}{section}{Step 1: install R}

\markright{Step 1: install R}

Go to \href{https://cran.r-project.org}{\textbf{cran.r-project.org}} and
choose the download link for your system.

\textbf{On Windows}, click \emph{Download R for Windows}, then
\emph{base}, then the download link at the top. Run the installer and
accept the defaults.

\textbf{On macOS}, click \emph{Download R for macOS} and pick the
correct \texttt{.pkg} file. There are usually two. If your Mac has an
Apple silicon chip (M1, M2, M3, or M4), take the \textbf{arm64} version.
If it has an Intel chip, take the other. To check, click the Apple menu,
then \emph{About This Mac}, and read the Chip or Processor line.

Choosing the wrong one on a Mac is the single most common installation
problem, and it produces confusing errors later rather than a clear
failure at install time.

\section*{Step 2: install RStudio}\label{step-2-install-rstudio}
\addcontentsline{toc}{section}{Step 2: install RStudio}

\markright{Step 2: install RStudio}

Go to
\href{https://posit.co/download/rstudio-desktop}{\textbf{posit.co/download/rstudio-desktop}}.
The page detects your system and offers the right installer. Take the
free Desktop version; there is no need for any paid tier or for RStudio
Server.

Install it, then open \textbf{RStudio}, not R. From here on, RStudio is
the only one of the two you will launch directly.

\section*{Step 3: check that it
worked}\label{step-3-check-that-it-worked}
\addcontentsline{toc}{section}{Step 3: check that it worked}

\markright{Step 3: check that it worked}

You should see four panes, or two panes with tabs. Find the
\textbf{Console}, on the left, and type this at the
\texttt{\textgreater{}} prompt, then press Enter.

\begin{Shaded}
\begin{Highlighting}[]
\DecValTok{2} \SpecialCharTok{+} \DecValTok{2}
\end{Highlighting}
\end{Shaded}

If \texttt{{[}1{]}\ 4} appears, both pieces are installed and talking to
each other.

Then check the version:

\begin{Shaded}
\begin{Highlighting}[]
\NormalTok{R.version.string}
\end{Highlighting}
\end{Shaded}

Anything from \textbf{R 4.3.0} onward is fine for this course. If you
have an older installation lying around from a previous course, update
it now rather than debugging a version problem in Week 3.

\section*{Step 4: install the
packages}\label{step-4-install-the-packages}
\addcontentsline{toc}{section}{Step 4: install the packages}

\markright{Step 4: install the packages}

A package is a bundle of functions someone else has written and shared.
R ships with a great many built in, and everything else is installed
once and then loaded whenever needed.

Run this in the Console. It will take several minutes and produce a
large amount of scrolling output, which is normal.

\begin{Shaded}
\begin{Highlighting}[]
\FunctionTok{install.packages}\NormalTok{(}\FunctionTok{c}\NormalTok{(}\StringTok{"tidyverse"}\NormalTok{, }\StringTok{"rmarkdown"}\NormalTok{, }\StringTok{"knitr"}\NormalTok{))}
\end{Highlighting}
\end{Shaded}

Then confirm they load:

\begin{Shaded}
\begin{Highlighting}[]
\FunctionTok{library}\NormalTok{(tidyverse)}
\end{Highlighting}
\end{Shaded}

You will see a short message listing the packages that came along with
it and noting a couple of conflicts. That message is informational, not
an error.

The remaining packages, for classification, text, and language model
work, are installed in the week they are first needed. There is no
advantage to installing everything now.

\section*{Step 5: install Quarto}\label{step-5-install-quarto}
\addcontentsline{toc}{section}{Step 5: install Quarto}

\markright{Step 5: install Quarto}

Quarto renders documents that combine code, prose, and output, and your
final project is submitted in that form.

Recent versions of RStudio include it already. Check in the Console:

\begin{Shaded}
\begin{Highlighting}[]
\NormalTok{quarto}\SpecialCharTok{::}\FunctionTok{quarto\_version}\NormalTok{()}
\end{Highlighting}
\end{Shaded}

If that returns a version number, you are done. If it reports that the
package is not installed, run \texttt{install.packages("quarto")} and
try again. If it still cannot find Quarto itself, download it from
\href{https://quarto.org/docs/get-started}{\textbf{quarto.org/docs/get-started}}
and install it separately.

\section*{When installation goes
wrong}\label{when-installation-goes-wrong}
\addcontentsline{toc}{section}{When installation goes wrong}

\markright{When installation goes wrong}

Three problems account for nearly all of them.

\textbf{``Package not available'' or a compilation error on macOS.}
Almost always the arm64 and Intel mismatch from Step 1. Uninstall R,
check your chip, and install the matching version.

\textbf{Permission denied, or the installation appears to succeed but
the package will not load.} Common on university-managed laptops, where
the default library folder is not writable. R usually offers to use a
personal library instead, and you should accept. If it does not offer,
run this once and then reinstall:

\begin{Shaded}
\begin{Highlighting}[]
\FunctionTok{dir.create}\NormalTok{(}\FunctionTok{Sys.getenv}\NormalTok{(}\StringTok{"R\_LIBS\_USER"}\NormalTok{), }\AttributeTok{recursive =} \ConstantTok{TRUE}\NormalTok{)}
\end{Highlighting}
\end{Shaded}

\textbf{The download stalls or times out.} More often a network problem
than a problem with your machine. Try again, and if it persists, try a
different connection or a different CRAN mirror, which you can choose
under \textbf{Tools \textgreater{} Global Options \textgreater{}
Packages}.

If none of these apply, bring the laptop to the first session. A
five-minute look in person is faster than an evening of email, and
installation problems are common enough that nobody should feel behind
because of one.

\section*{A little configuration}\label{a-little-configuration}
\addcontentsline{toc}{section}{A little configuration}

\markright{A little configuration}

Two settings are worth changing now. Go to \textbf{Tools \textgreater{}
Global Options \textgreater{} General} and:

\begin{itemize}
\tightlist
\item
  Set \textbf{Save workspace to .RData on exit} to \textbf{Never}.
\item
  Untick \textbf{Restore .RData into workspace at startup}.
\end{itemize}

Chapter 3 explains why these matter. In short, they are what make it
possible to trust that your code actually works rather than merely
appearing to.

While you are in Global Options, \textbf{Appearance} lets you change the
theme and font size. Not important, but you will be looking at this
window for eight weeks, so it is worth a minute.

\section*{Before you move on}\label{before-you-move-on}
\addcontentsline{toc}{section}{Before you move on}

\markright{Before you move on}

You should be able to open RStudio, type \texttt{2\ +\ 2} in the
Console, get \texttt{4}, and run \texttt{library(tidyverse)} without an
error.

That is the whole requirement. The next chapter starts from there.

\chapter*{Four Ways to Run R Code}\label{four-ways-to-run-r-code}
\addcontentsline{toc}{chapter}{Four Ways to Run R Code}

\markboth{Four Ways to Run R Code}{Four Ways to Run R Code}

Within the first hour of using RStudio you will encounter at least four
different places where R code can be typed and executed: the console, a
script, an R Markdown document, and a Quarto document. Students almost
always assume that one of these is the real way to work and that the
others are variations on it, and they then spend the semester wondering
which one they were supposed to be using. The assumption is
understandable, and it is wrong. The four are not competing versions of
the same thing. Each answers a different question, and the question you
are asking should determine which one you open.

The distinction matters more in this course than it would in an
introductory statistics class. Your final project is submitted as a
document that must run from raw inputs to reported numbers on a machine
that is not yours. Ten percent of its grade depends on that property
alone. Understanding what each of these environments preserves, and what
each of them quietly discards, is therefore not a preliminary detail. It
is the foundation of everything you submit.

\section*{The console: thinking out
loud}\label{the-console-thinking-out-loud}
\addcontentsline{toc}{section}{The console: thinking out loud}

\markright{The console: thinking out loud}

The console is the panel where R is actually running. When you type
\texttt{2\ +\ 2} and press Enter, the console evaluates the expression,
prints the answer, and forgets that you asked.

\begin{Shaded}
\begin{Highlighting}[]
\DecValTok{2} \SpecialCharTok{+} \DecValTok{2}
\CommentTok{\#\textgreater{} [1] 4}
\end{Highlighting}
\end{Shaded}

That last property is the important one. The console keeps a history you
can scroll through, but it keeps no record you can hand to anyone, and
it imposes no order on what you did. If you fit a model, adjust a
variable, fit it again, and then produce a result you are pleased with,
the console cannot tell you which sequence of steps produced it. Neither
can you, a week later.

The console is therefore the right tool for questions you are asking of
yourself rather than of the world. Checking the dimensions of a dataset,
inspecting a variable, reading a help file, and testing whether a line
does what you expect are all console work. Use it constantly, and use it
freely. Nothing you do there is precious, which is precisely what makes
it useful.

\section*{The script: a record of
instructions}\label{the-script-a-record-of-instructions}
\addcontentsline{toc}{section}{The script: a record of instructions}

\markright{The script: a record of instructions}

An R script, saved with the extension \texttt{.R}, is a plain text file
containing lines of R code and nothing else. When you run a script, R
executes each line in order, exactly as if you had typed them into the
console yourself.

The difference is that the file survives. A script is a record of what
you did, in the order you did it, and it can be run again tomorrow, or
by someone else, or by you on a different machine. This is what makes it
the first genuinely reproducible unit of work, and it is why the
transition from console to script is the single most consequential habit
shift in learning R.

\begin{Shaded}
\begin{Highlighting}[]
\CommentTok{\# analysis.R}
\FunctionTok{library}\NormalTok{(tidyverse)}

\NormalTok{fl }\OtherTok{\textless{}{-}} \FunctionTok{read.csv}\NormalTok{(}\StringTok{"https://raw.githubusercontent.com/babakrezaee/AppliedAI{-}PoliSci/main/data/FearonLaitin2003.csv"}\NormalTok{)}

\NormalTok{fl\_clean }\OtherTok{\textless{}{-}}\NormalTok{ fl }\SpecialCharTok{\%\textgreater{}\%}
  \FunctionTok{filter}\NormalTok{(}\SpecialCharTok{!}\FunctionTok{is.na}\NormalTok{(onset)) }\SpecialCharTok{\%\textgreater{}\%}
  \FunctionTok{mutate}\NormalTok{(}\AttributeTok{log\_pop =} \FunctionTok{log}\NormalTok{(pop))}

\NormalTok{model }\OtherTok{\textless{}{-}} \FunctionTok{glm}\NormalTok{(onset }\SpecialCharTok{\textasciitilde{}}\NormalTok{ log\_pop }\SpecialCharTok{+}\NormalTok{ gdpen, }\AttributeTok{data =}\NormalTok{ fl\_clean, }\AttributeTok{family =} \StringTok{"binomial"}\NormalTok{)}
\FunctionTok{summary}\NormalTok{(model)}
\end{Highlighting}
\end{Shaded}

As explained in the previous section, you can check the dimension of
your datasets and summarize your variables as well:

\begin{Shaded}
\begin{Highlighting}[]
\FunctionTok{dim}\NormalTok{(fl)}
\FunctionTok{summary}\NormalTok{(fl}\SpecialCharTok{$}\NormalTok{gdpen)}
\NormalTok{?mean}
\end{Highlighting}
\end{Shaded}

What a script preserves is the instructions. What it does not preserve
is the output. Run the script and the results appear in the console,
where they are as ephemeral as anything else typed there. Run it again
in three months, after the data file has been updated or a package has
changed, and you may get different numbers with no record of what the
old ones were. Scripts are also, by design, hostile to prose. You can
write comments, and you should, but a comment is a note to yourself, not
an argument addressed to a reader.

Scripts remain the right choice for work whose product is an action
rather than a document: downloading and cleaning a dataset, defining
functions you will use elsewhere, or running a long job whose output is
a saved file. In the repository for this course you will find several
such scripts in the \texttt{R/} folder.

\section*{R Markdown: code, prose, and output in one
file}\label{r-markdown-code-prose-and-output-in-one-file}
\addcontentsline{toc}{section}{R Markdown: code, prose, and output in
one file}

\markright{R Markdown: code, prose, and output in one file}

An R Markdown document, saved as \texttt{.Rmd}, holds three things at
once: prose written in Markdown, chunks of R code, and the output those
chunks produce. When you knit the document, R runs every chunk in order,
captures the results, and weaves them into a finished report as HTML,
PDF, or Word.

\begin{Shaded}
\begin{Highlighting}[]
\CommentTok{{-}{-}{-}}
\AnnotationTok{title:}\CommentTok{ "Civil War Onset"}
\AnnotationTok{output:}\CommentTok{ html\_document}
\CommentTok{{-}{-}{-}}

\NormalTok{Fearon and Laitin (2003) report that per capita income predicts}
\NormalTok{civil war onset more reliably than measures of ethnic diversity.}
\NormalTok{The model below reproduces that pattern.}

\InformationTok{\textasciigrave{}\textasciigrave{}\textasciigrave{}\{r\}}
\NormalTok{model }\OtherTok{\textless{}{-}} \FunctionTok{glm}\NormalTok{(onset }\SpecialCharTok{\textasciitilde{}}\NormalTok{ gdpen }\SpecialCharTok{+}\NormalTok{ ethfrac, }\AttributeTok{data =}\NormalTok{ fl, }\AttributeTok{family =} \StringTok{"binomial"}\NormalTok{)}
\FunctionTok{summary}\NormalTok{(model)}
\InformationTok{\textasciigrave{}\textasciigrave{}\textasciigrave{}}

\NormalTok{The coefficient on income is negative and precisely estimated,}
\NormalTok{while the coefficient on ethnic fractionalization is not}
\NormalTok{distinguishable from zero.}
\end{Highlighting}
\end{Shaded}

Two consequences follow, and both matter for this course.

The first is that the numbers in your text cannot drift away from the
numbers your code produced, because the document generates them each
time it is knitted. A reported coefficient that was copied by hand into
a sentence will eventually become wrong, quietly, when the analysis
changes and the sentence does not. A reported coefficient that is
produced by the code cannot.

The second is subtler and catches almost everyone at least once.
Knitting starts a fresh R session that knows nothing about your
workspace. Any object you created in the console, any package you loaded
by hand, and any file path that happens to work on your machine are all
invisible to it. A document that runs perfectly line by line in your own
session will therefore fail to knit if it depends on something you did
outside it. This is not an inconvenience to be worked around. It is the
mechanism that makes the format trustworthy, because a document that
knits from a clean session is a document that will run on someone else's
computer.

The habit that follows is simple and worth adopting now: before you
submit anything, restart R (Session, then Restart R) and knit. If it
does not run from clean, it does not run.

\section*{Quarto: the same idea, more
broadly}\label{quarto-the-same-idea-more-broadly}
\addcontentsline{toc}{section}{Quarto: the same idea, more broadly}

\markright{Quarto: the same idea, more broadly}

Quarto, saved as \texttt{.qmd}, is the successor to R Markdown.
Everything described above applies to it unchanged, and if you have
understood R Markdown you already understand Quarto.

The differences are largely a matter of reach. Quarto is not tied to R,
so the same document can execute chunks written in Python or Julia,
which is directly relevant given that most of the applied AI ecosystem
lives in Python. Its options are declared in a consistent way across
output formats, and it renders books and websites as readily as single
documents. The materials you are reading now are a Quarto book.

\begin{Shaded}
\begin{Highlighting}[]
\InformationTok{\textasciigrave{}\textasciigrave{}\textasciigrave{}\{r\}}
\CommentTok{\#| label: fig{-}onset}
\CommentTok{\#| fig{-}cap: "Predicted probability of onset by income"}
\CommentTok{\#| warning: false}

\FunctionTok{ggplot}\NormalTok{(fl, }\FunctionTok{aes}\NormalTok{(}\AttributeTok{x =}\NormalTok{ gdpen, }\AttributeTok{y =}\NormalTok{ fitted)) }\SpecialCharTok{+}
  \FunctionTok{geom\_line}\NormalTok{()}
\InformationTok{\textasciigrave{}\textasciigrave{}\textasciigrave{}}
\end{Highlighting}
\end{Shaded}

For a single analysis document the practical difference between the two
is small, and either is acceptable for your final project. I use Quarto
for course materials because it handles a multi-chapter book cleanly,
and I would suggest it for new work on the general principle that
development effort is going there rather than into R Markdown.

\section*{Choosing between them}\label{choosing-between-them}
\addcontentsline{toc}{section}{Choosing between them}

\markright{Choosing between them}

{\def\LTcaptype{none} % do not increment counter
\begin{longtable}[]{@{}
  >{\raggedright\arraybackslash}p{(\linewidth - 8\tabcolsep) * \real{0.2000}}
  >{\raggedright\arraybackslash}p{(\linewidth - 8\tabcolsep) * \real{0.2000}}
  >{\raggedright\arraybackslash}p{(\linewidth - 8\tabcolsep) * \real{0.2000}}
  >{\raggedright\arraybackslash}p{(\linewidth - 8\tabcolsep) * \real{0.2000}}
  >{\raggedright\arraybackslash}p{(\linewidth - 8\tabcolsep) * \real{0.2000}}@{}}
\toprule\noalign{}
\begin{minipage}[b]{\linewidth}\raggedright
\end{minipage} & \begin{minipage}[b]{\linewidth}\raggedright
Console
\end{minipage} & \begin{minipage}[b]{\linewidth}\raggedright
Script (\texttt{.R})
\end{minipage} & \begin{minipage}[b]{\linewidth}\raggedright
R Markdown (\texttt{.Rmd})
\end{minipage} & \begin{minipage}[b]{\linewidth}\raggedright
Quarto (\texttt{.qmd})
\end{minipage} \\
\midrule\noalign{}
\endhead
\bottomrule\noalign{}
\endlastfoot
Preserves code & No & Yes & Yes & Yes \\
Preserves output & No & No & Yes & Yes \\
Holds prose & No & Comments only & Yes & Yes \\
Runs in a clean session & Not applicable & On request & Every render &
Every render \\
Languages & R & R & R, mainly & R, Python, Julia \\
\end{longtable}
}

The rule of thumb is short. Explore in the console, because it costs
nothing and you will throw the work away. Automate in a script, when the
product is an action rather than an argument. Report in Quarto or R
Markdown, when someone else needs to see both what you concluded and how
you got there.

Your final project falls squarely in the third category, which is why
the submission format is a \texttt{.qmd} or \texttt{.Rmd} file together
with the document it renders. The requirement is not a formatting
preference. A project submitted as a script and a separate write-up
cannot demonstrate that the numbers in the write-up came from the
script, and demonstrating exactly that is a substantial part of what the
project is for.

\section*{Exercise}\label{exercise}
\addcontentsline{toc}{section}{Exercise}

\markright{Exercise}

Work through this in order. It takes about fifteen minutes and it is
worth doing before the first session rather than during it.

\begin{enumerate}
\def\labelenumi{\arabic{enumi}.}
\item
  In the console, compute \texttt{mean(c(4,\ 8,\ 15,\ 16,\ 23,\ 42))}.
  Note that the answer appears and then nothing remains except a history
  entry.
\item
  Open a new script (File, New File, R Script). Write the same line,
  save the file as \texttt{first-script.R}, and run it with the Source
  button. The result appears in the console again, but the instruction
  now lives in a file you could send to someone.
\item
  Open a new Quarto document (File, New File, Quarto Document). Write
  one sentence of prose describing what you are about to compute, then a
  code chunk containing the same line, then one sentence interpreting
  the result. Render it.
\item
  Now break it deliberately. In the console, create an object with
  \texttt{x\ \textless{}-\ c(4,\ 8,\ 15,\ 16,\ 23,\ 42)}. Change the
  chunk in your Quarto document to \texttt{mean(x)} and render again.
\end{enumerate}

The render fails, and the error will say that \texttt{x} was not found.
Sit with that for a moment before fixing it, because you have just met
the most common reason a working analysis refuses to knit the night
before a deadline. The object exists in your session and does not exist
in the document. The fix is to define \texttt{x} inside the document,
which is also the reason the format can be trusted.

\chapter*{Objects, Vectors, and Data
Frames}\label{objects-vectors-and-data-frames}
\addcontentsline{toc}{chapter}{Objects, Vectors, and Data Frames}

\markboth{Objects, Vectors, and Data Frames}{Objects, Vectors, and Data
Frames}

Everything in R is an object, and almost everything you will do in this
course amounts to creating an object, inspecting it, and passing it to a
function. This chapter covers the three kinds of object you will meet
constantly: the vector, which holds a column of values; the data frame,
which holds a table of them; and the individual value, which is really
just a vector of length one.

The order matters. Data frames are what political science data actually
look like, and they are the destination. Vectors come first because a
data frame is built out of them, and because most of the errors you will
hit in Week 2 come from a vector quietly being the wrong type.

\section*{Assignment}\label{assignment}
\addcontentsline{toc}{section}{Assignment}

\markright{Assignment}

An object is created by assigning something a name. R accepts two
operators for this, and you will see both in the wild.

\begin{Shaded}
\begin{Highlighting}[]
\NormalTok{x }\OtherTok{\textless{}{-}} \DecValTok{5}
\NormalTok{y }\OtherTok{=} \DecValTok{5}
\end{Highlighting}
\end{Shaded}

Use \texttt{\textless{}-}. The two are equivalent in this context, but
\texttt{=} also means other things in R, including naming an argument
inside a function call, so keeping the two visually distinct saves
confusion later. In RStudio, \textbf{Alt and minus} on Windows or
\textbf{Option and minus} on a Mac types it for you.

To see what an object contains, type its name.

\begin{Shaded}
\begin{Highlighting}[]
\NormalTok{x}
\CommentTok{\#\textgreater{} [1] 5}
\end{Highlighting}
\end{Shaded}

To remove one, use \texttt{rm()}.

\begin{Shaded}
\begin{Highlighting}[]
\FunctionTok{rm}\NormalTok{(x)}
\end{Highlighting}
\end{Shaded}

Names are subject to a few rules. They may contain letters, numbers, a
period, and an underscore, but not spaces or symbols such as \texttt{!},
\texttt{+}, or \texttt{\#}. They cannot begin with a number. And R is
case sensitive, so \texttt{polity} and \texttt{Polity} are two different
objects, which is a more common source of ``object not found'' than
anyone expects.

\section*{Vectors}\label{vectors}
\addcontentsline{toc}{section}{Vectors}

\markright{Vectors}

A vector is an ordered collection of values of the same type. You build
one with \texttt{c()}, for combine.

\begin{Shaded}
\begin{Highlighting}[]
\NormalTok{gdp }\OtherTok{\textless{}{-}} \FunctionTok{c}\NormalTok{(}\FloatTok{2.1}\NormalTok{, }\FloatTok{4.7}\NormalTok{, }\FloatTok{1.3}\NormalTok{, }\FloatTok{8.9}\NormalTok{)}
\NormalTok{gdp}
\CommentTok{\#\textgreater{} [1] 2.1 4.7 1.3 8.9}
\end{Highlighting}
\end{Shaded}

Sequences of consecutive integers have a shorthand.

\begin{Shaded}
\begin{Highlighting}[]
\NormalTok{years }\OtherTok{\textless{}{-}} \DecValTok{1990}\SpecialCharTok{:}\DecValTok{1995}
\NormalTok{years}
\CommentTok{\#\textgreater{} [1] 1990 1991 1992 1993 1994 1995}
\end{Highlighting}
\end{Shaded}

The phrase ``of the same type'' is the part worth pausing on, and it is
where the first real surprise lives.

\section*{Types, and why they matter}\label{types-and-why-they-matter}
\addcontentsline{toc}{section}{Types, and why they matter}

\markright{Types, and why they matter}

A vector holds numbers, or text, or logical values, but never a mixture.
If you supply a mixture, R does not complain. It silently converts
everything to the most accommodating type, which is almost always text.

\begin{Shaded}
\begin{Highlighting}[]
\NormalTok{v1 }\OtherTok{\textless{}{-}} \FunctionTok{c}\NormalTok{(}\DecValTok{1}\NormalTok{, }\DecValTok{2}\NormalTok{, }\DecValTok{3}\NormalTok{)}
\NormalTok{v2 }\OtherTok{\textless{}{-}} \FunctionTok{c}\NormalTok{(}\DecValTok{1}\NormalTok{, }\DecValTok{2}\NormalTok{, }\StringTok{"a"}\NormalTok{)}

\NormalTok{v1}
\CommentTok{\#\textgreater{} [1] 1 2 3}

\NormalTok{v2}
\CommentTok{\#\textgreater{} [1] "1" "2" "a"}
\end{Highlighting}
\end{Shaded}

Look closely at \texttt{v2}. The numbers now have quotation marks around
them. They are no longer numbers, and arithmetic on them will fail.

\begin{Shaded}
\begin{Highlighting}[]
\FunctionTok{mean}\NormalTok{(v2)}
\CommentTok{\#\textgreater{} Warning: argument is not numeric or logical: returning NA}
\end{Highlighting}
\end{Shaded}

This behavior, called coercion, is the single most common reason a
calculation returns \texttt{NA} for no visible reason. A real dataset
with one stray text entry in an otherwise numeric column, for instance a
\texttt{"n/a"} typed by hand, will import as a text column and every
subsequent calculation on it will quietly fail.

The habit that protects you is to check the type whenever a result
surprises you.

\begin{Shaded}
\begin{Highlighting}[]
\FunctionTok{class}\NormalTok{(v1)}
\CommentTok{\#\textgreater{} [1] "numeric"}

\FunctionTok{class}\NormalTok{(v2)}
\CommentTok{\#\textgreater{} [1] "character"}
\end{Highlighting}
\end{Shaded}

Logical vectors hold \texttt{TRUE} and \texttt{FALSE}, and you will meet
them constantly as the result of a comparison.

\begin{Shaded}
\begin{Highlighting}[]
\NormalTok{gdp }\SpecialCharTok{\textgreater{}} \DecValTok{3}
\CommentTok{\#\textgreater{} [1] FALSE  TRUE FALSE  TRUE}
\end{Highlighting}
\end{Shaded}

\section*{Arithmetic happens element by
element}\label{arithmetic-happens-element-by-element}
\addcontentsline{toc}{section}{Arithmetic happens element by element}

\markright{Arithmetic happens element by element}

Operations on vectors apply to every element at once, which is why R
rarely needs the loops that other languages depend on.

\begin{Shaded}
\begin{Highlighting}[]
\NormalTok{a }\OtherTok{\textless{}{-}} \FunctionTok{c}\NormalTok{(}\DecValTok{1}\NormalTok{, }\DecValTok{3}\NormalTok{, }\DecValTok{5}\NormalTok{, }\DecValTok{7}\NormalTok{)}
\NormalTok{b }\OtherTok{\textless{}{-}} \FunctionTok{c}\NormalTok{(}\DecValTok{1}\NormalTok{, }\DecValTok{2}\NormalTok{, }\DecValTok{4}\NormalTok{, }\DecValTok{8}\NormalTok{)}

\NormalTok{a }\SpecialCharTok{*} \DecValTok{5}
\CommentTok{\#\textgreater{} [1]  5 15 25 35}

\NormalTok{a }\SpecialCharTok{+}\NormalTok{ b}
\CommentTok{\#\textgreater{} [1]  2  5  9 15}
\end{Highlighting}
\end{Shaded}

If the two vectors are of unequal length, R recycles the shorter one to
match the longer, rather than raising an error.

\begin{Shaded}
\begin{Highlighting}[]
\FunctionTok{c}\NormalTok{(}\DecValTok{10}\NormalTok{, }\DecValTok{20}\NormalTok{) }\SpecialCharTok{+} \FunctionTok{c}\NormalTok{(}\DecValTok{1}\NormalTok{, }\DecValTok{2}\NormalTok{, }\DecValTok{3}\NormalTok{, }\DecValTok{4}\NormalTok{)}
\CommentTok{\#\textgreater{} [1] 11 22 13 24}
\end{Highlighting}
\end{Shaded}

This is occasionally convenient and frequently dangerous. If you did not
intend the recycling, you now have a result that looks plausible and is
wrong. When two vectors should be the same length, check with
\texttt{length()} before combining them.

\section*{Getting at parts of a
vector}\label{getting-at-parts-of-a-vector}
\addcontentsline{toc}{section}{Getting at parts of a vector}

\markright{Getting at parts of a vector}

Square brackets extract elements by position, counting from one rather
than from zero.

\begin{Shaded}
\begin{Highlighting}[]
\NormalTok{gdp[}\DecValTok{2}\NormalTok{]}
\CommentTok{\#\textgreater{} [1] 4.7}

\NormalTok{gdp[}\FunctionTok{c}\NormalTok{(}\DecValTok{1}\NormalTok{, }\DecValTok{3}\NormalTok{)]}
\CommentTok{\#\textgreater{} [1] 2.1 1.3}

\NormalTok{gdp[}\SpecialCharTok{{-}}\DecValTok{1}\NormalTok{]        }\CommentTok{\# everything except the first}
\CommentTok{\#\textgreater{} [1] 4.7 1.3 8.9}
\end{Highlighting}
\end{Shaded}

You can also select with a logical condition, which is the form you will
use most.

\begin{Shaded}
\begin{Highlighting}[]
\NormalTok{gdp[gdp }\SpecialCharTok{\textgreater{}} \DecValTok{3}\NormalTok{]}
\CommentTok{\#\textgreater{} [1] 4.7 8.9}
\end{Highlighting}
\end{Shaded}

\section*{Missing values}\label{missing-values}
\addcontentsline{toc}{section}{Missing values}

\markright{Missing values}

Missing data is written \texttt{NA}, and it is contagious by design.

\begin{Shaded}
\begin{Highlighting}[]
\NormalTok{pop }\OtherTok{\textless{}{-}} \FunctionTok{c}\NormalTok{(}\DecValTok{12}\NormalTok{, }\DecValTok{45}\NormalTok{, }\ConstantTok{NA}\NormalTok{, }\DecValTok{8}\NormalTok{)}
\FunctionTok{mean}\NormalTok{(pop)}
\CommentTok{\#\textgreater{} [1] NA}
\end{Highlighting}
\end{Shaded}

R is telling you something true: the mean of a set that includes an
unknown value is unknown. If you are content to ignore the missing case,
say so explicitly.

\begin{Shaded}
\begin{Highlighting}[]
\FunctionTok{mean}\NormalTok{(pop, }\AttributeTok{na.rm =} \ConstantTok{TRUE}\NormalTok{)}
\CommentTok{\#\textgreater{} [1] 21.66667}
\end{Highlighting}
\end{Shaded}

Say so deliberately rather than reflexively. Dropping missing cases is a
substantive decision about your sample, and the final project asks you
to make such decisions visible rather than silent.

To find them:

\begin{Shaded}
\begin{Highlighting}[]
\FunctionTok{is.na}\NormalTok{(pop)}
\CommentTok{\#\textgreater{} [1] FALSE FALSE  TRUE FALSE}

\FunctionTok{sum}\NormalTok{(}\FunctionTok{is.na}\NormalTok{(pop))}
\CommentTok{\#\textgreater{} [1] 1}
\end{Highlighting}
\end{Shaded}

\section*{Data frames}\label{data-frames}
\addcontentsline{toc}{section}{Data frames}

\markright{Data frames}

A data frame is a table: several vectors of equal length standing side
by side as columns, each with a name, each with its own type. This is
what a political science dataset is.

Load one. The \texttt{read.csv()} function accepts a web address
directly, so nothing needs downloading.

\begin{Shaded}
\begin{Highlighting}[]
\NormalTok{fl }\OtherTok{\textless{}{-}} \FunctionTok{read.csv}\NormalTok{(}\StringTok{"https://raw.githubusercontent.com/babakrezaee/AppliedAI{-}PoliSci/main/data/FearonLaitin2003.csv"}\NormalTok{)}
\end{Highlighting}
\end{Shaded}

These are the data behind Fearon and Laitin's 2003 article on ethnicity,
insurgency, and civil war, one of the most cited studies in the conflict
literature. The unit of analysis is the country-year, and we will return
to it in Week 6.

\section*{Look before you calculate}\label{look-before-you-calculate}
\addcontentsline{toc}{section}{Look before you calculate}

\markright{Look before you calculate}

Before doing anything with a new dataset, find out what you have.

\begin{Shaded}
\begin{Highlighting}[]
\FunctionTok{dim}\NormalTok{(fl)          }\CommentTok{\# rows and columns}
\FunctionTok{nrow}\NormalTok{(fl)}
\FunctionTok{names}\NormalTok{(fl)        }\CommentTok{\# column names}
\FunctionTok{head}\NormalTok{(fl, }\AttributeTok{n =} \DecValTok{3}\NormalTok{)  }\CommentTok{\# first three rows}
\FunctionTok{str}\NormalTok{(fl)          }\CommentTok{\# structure: every column, its type, first few values}
\end{Highlighting}
\end{Shaded}

\texttt{str()} is the most useful of these, because it reports the type
of every column. A numeric variable that arrived as text will announce
itself here rather than three steps later.

For a single variable, \texttt{summary()} gives the distribution and,
importantly, the count of missing values.

\begin{Shaded}
\begin{Highlighting}[]
\FunctionTok{summary}\NormalTok{(fl}\SpecialCharTok{$}\NormalTok{gdpen)}
\end{Highlighting}
\end{Shaded}

The dollar sign extracts one column by name. It is the workhorse of base
R and you will use it constantly.

\section*{An example of why you look
first}\label{an-example-of-why-you-look-first}
\addcontentsline{toc}{section}{An example of why you look first}

\markright{An example of why you look first}

The outcome variable in these data is \texttt{onset}, coded one in a
country-year when a civil war begins and zero otherwise. Count its
values.

\begin{Shaded}
\begin{Highlighting}[]
\FunctionTok{table}\NormalTok{(fl}\SpecialCharTok{$}\NormalTok{onset)}
\end{Highlighting}
\end{Shaded}

You will find a handful of zeros and ones, as expected, and also a
single observation coded \texttt{4}. That is a data entry error, present
in the distributed file, and it has survived into a great many
replications.

It matters. A model fitted with \texttt{onset} treated as numeric would
take that four seriously. Left as a factor with three levels, a
classifier would try to predict a category that occurs once in six
thousand cases. Either way the result is wrong in a manner no error
message will announce.

The fix is one line.

\begin{Shaded}
\begin{Highlighting}[]
\NormalTok{fl}\SpecialCharTok{$}\NormalTok{onset[fl}\SpecialCharTok{$}\NormalTok{onset }\SpecialCharTok{\textgreater{}} \DecValTok{1}\NormalTok{] }\OtherTok{\textless{}{-}} \DecValTok{1}
\FunctionTok{table}\NormalTok{(fl}\SpecialCharTok{$}\NormalTok{onset)}
\end{Highlighting}
\end{Shaded}

The lesson generalizes beyond this dataset. Run \texttt{table()} on
every categorical variable and \texttt{summary()} on every numeric one
before you model anything. It takes two minutes and it is the cheapest
error-catching you will ever do.

\section*{Selecting rows and columns}\label{selecting-rows-and-columns}
\addcontentsline{toc}{section}{Selecting rows and columns}

\markright{Selecting rows and columns}

Data frames take two indices inside the brackets, rows first and columns
second.

\begin{Shaded}
\begin{Highlighting}[]
\NormalTok{fl[}\DecValTok{1}\NormalTok{, ]            }\CommentTok{\# first row, all columns}
\NormalTok{fl[, }\StringTok{"onset"}\NormalTok{]      }\CommentTok{\# the onset column}
\NormalTok{fl[}\DecValTok{1}\SpecialCharTok{:}\DecValTok{5}\NormalTok{, }\FunctionTok{c}\NormalTok{(}\StringTok{"country"}\NormalTok{, }\StringTok{"year"}\NormalTok{, }\StringTok{"onset"}\NormalTok{)]}
\end{Highlighting}
\end{Shaded}

To keep rows meeting a condition:

\begin{Shaded}
\begin{Highlighting}[]
\NormalTok{africa }\OtherTok{\textless{}{-}}\NormalTok{ fl[fl}\SpecialCharTok{$}\NormalTok{ssafrica }\SpecialCharTok{==} \DecValTok{1}\NormalTok{, ]}
\FunctionTok{nrow}\NormalTok{(africa)}
\end{Highlighting}
\end{Shaded}

Note the double equals sign. A single \texttt{=} assigns; a double
\texttt{==} asks whether two things are equal. Mixing them up is a rite
of passage.

\section*{The tidyverse, and the pipe}\label{the-tidyverse-and-the-pipe}
\addcontentsline{toc}{section}{The tidyverse, and the pipe}

\markright{The tidyverse, and the pipe}

Base R can do everything above, and for many tasks the tidyverse offers
a version that reads more like a sentence. Install it once, then load it
in each session.

\begin{Shaded}
\begin{Highlighting}[]
\CommentTok{\# install.packages("tidyverse")   \# once, ever}
\FunctionTok{library}\NormalTok{(tidyverse)                }\CommentTok{\# every session}
\end{Highlighting}
\end{Shaded}

The pipe operator \texttt{\%\textgreater{}\%} takes whatever is on its
left and passes it as the first argument to the function on its right.
Written this way, a chain of operations reads in the order it happens.

\begin{Shaded}
\begin{Highlighting}[]
\NormalTok{fl }\SpecialCharTok{\%\textgreater{}\%}
  \FunctionTok{filter}\NormalTok{(ssafrica }\SpecialCharTok{==} \DecValTok{1}\NormalTok{, year }\SpecialCharTok{\textgreater{}=} \DecValTok{1980}\NormalTok{) }\SpecialCharTok{\%\textgreater{}\%}
  \FunctionTok{select}\NormalTok{(country, year, onset, gdpen) }\SpecialCharTok{\%\textgreater{}\%}
  \FunctionTok{head}\NormalTok{()}
\end{Highlighting}
\end{Shaded}

Compare that with the base R equivalent, which has to be read from the
inside out:

\begin{Shaded}
\begin{Highlighting}[]
\FunctionTok{head}\NormalTok{(fl[fl}\SpecialCharTok{$}\NormalTok{ssafrica }\SpecialCharTok{==} \DecValTok{1} \SpecialCharTok{\&}\NormalTok{ fl}\SpecialCharTok{$}\NormalTok{year }\SpecialCharTok{\textgreater{}=} \DecValTok{1980}\NormalTok{, }\FunctionTok{c}\NormalTok{(}\StringTok{"country"}\NormalTok{, }\StringTok{"year"}\NormalTok{, }\StringTok{"onset"}\NormalTok{, }\StringTok{"gdpen"}\NormalTok{)])}
\end{Highlighting}
\end{Shaded}

Both are correct. The piped version is easier to modify six weeks later,
which is why the handouts use it. In RStudio, \textbf{Cmd and Shift and
M} on a Mac, or \textbf{Ctrl and Shift and M} on Windows, types the
pipe.

Three verbs will carry most of your work. \texttt{filter()} keeps rows,
\texttt{select()} keeps columns, and \texttt{mutate()} creates new ones.

\begin{Shaded}
\begin{Highlighting}[]
\NormalTok{fl }\SpecialCharTok{\%\textgreater{}\%}
  \FunctionTok{filter}\NormalTok{(}\SpecialCharTok{!}\FunctionTok{is.na}\NormalTok{(gdpen)) }\SpecialCharTok{\%\textgreater{}\%}
  \FunctionTok{mutate}\NormalTok{(}\AttributeTok{log\_pop =} \FunctionTok{log}\NormalTok{(pop)) }\SpecialCharTok{\%\textgreater{}\%}
  \FunctionTok{select}\NormalTok{(country, year, gdpen, log\_pop) }\SpecialCharTok{\%\textgreater{}\%}
  \FunctionTok{head}\NormalTok{()}
\end{Highlighting}
\end{Shaded}

\section*{Factors}\label{factors}
\addcontentsline{toc}{section}{Factors}

\markright{Factors}

A factor is R's type for a categorical variable: values drawn from a
fixed set of levels. Modeling functions expect categorical outcomes as
factors, so you will convert deliberately.

\begin{Shaded}
\begin{Highlighting}[]
\NormalTok{fl}\SpecialCharTok{$}\NormalTok{onset }\OtherTok{\textless{}{-}} \FunctionTok{factor}\NormalTok{(fl}\SpecialCharTok{$}\NormalTok{onset)}
\FunctionTok{levels}\NormalTok{(fl}\SpecialCharTok{$}\NormalTok{onset)}
\CommentTok{\#\textgreater{} [1] "0" "1"}
\end{Highlighting}
\end{Shaded}

One warning worth carrying. Converting a factor straight to numeric
returns the level codes, not the labels, so a factor with levels ``0''
and ``1'' becomes 1 and 2. If you need the underlying numbers, route
through character first.

\begin{Shaded}
\begin{Highlighting}[]
\FunctionTok{as.numeric}\NormalTok{(}\FunctionTok{as.character}\NormalTok{(fl}\SpecialCharTok{$}\NormalTok{onset))}
\end{Highlighting}
\end{Shaded}

\section*{Matrices, briefly}\label{matrices-briefly}
\addcontentsline{toc}{section}{Matrices, briefly}

\markright{Matrices, briefly}

A matrix is a rectangle of values that are all the same type, which is
the difference from a data frame. You will not build many, but they
appear inside modeling functions, and older code uses \texttt{cbind()}
and \texttt{rbind()} to assemble them.

\begin{Shaded}
\begin{Highlighting}[]
\NormalTok{m }\OtherTok{\textless{}{-}} \FunctionTok{matrix}\NormalTok{(}\FunctionTok{c}\NormalTok{(}\DecValTok{2}\NormalTok{, }\DecValTok{4}\NormalTok{, }\DecValTok{3}\NormalTok{, }\DecValTok{1}\NormalTok{, }\DecValTok{5}\NormalTok{, }\DecValTok{7}\NormalTok{), }\AttributeTok{nrow =} \DecValTok{3}\NormalTok{, }\AttributeTok{ncol =} \DecValTok{2}\NormalTok{)}
\NormalTok{m}
\FunctionTok{t}\NormalTok{(m)               }\CommentTok{\# transpose}
\FunctionTok{cbind}\NormalTok{(m, }\FunctionTok{c}\NormalTok{(}\DecValTok{7}\NormalTok{, }\DecValTok{4}\NormalTok{, }\DecValTok{2}\NormalTok{))  }\CommentTok{\# add a column}
\end{Highlighting}
\end{Shaded}

Recognizing them is enough for now. Week 4 returns to matrices when text
becomes a document-term matrix.

\section*{Getting help}\label{getting-help}
\addcontentsline{toc}{section}{Getting help}

\markright{Getting help}

Every function has a help page, reachable with a question mark.

\begin{Shaded}
\begin{Highlighting}[]
\NormalTok{?mean}
\NormalTok{?read.csv}
\end{Highlighting}
\end{Shaded}

The \textbf{Examples} section at the bottom is usually the fastest way
in, because you can run those lines directly. If you do not know the
function's name, \texttt{??regression} searches the text of all
installed help pages.

\section*{Exercise}\label{exercise-1}
\addcontentsline{toc}{section}{Exercise}

\markright{Exercise}

Roughly twenty minutes. Do it in a script rather than the console, so
you still have it afterward.

\begin{enumerate}
\def\labelenumi{\arabic{enumi}.}
\item
  Create a numeric vector of five values and a character vector of five
  country names. Confirm the type of each with \texttt{class()}.
\item
  Combine one number and one country name into a single vector with
  \texttt{c()}. Check its class. Explain in a comment what happened and
  why.
\item
  Load the Fearon and Laitin data. Report how many rows and columns it
  has.
\item
  Run \texttt{table()} on \texttt{onset}. Note the anomaly, correct it,
  and run \texttt{table()} again to confirm.
\item
  Using the pipe, keep only country-years from 1990 onward with
  \texttt{onset} equal to one, and select \texttt{country},
  \texttt{year}, and \texttt{gdpen}. How many civil war onsets does that
  leave?
\item
  Compute the mean of \texttt{gdpen} across the whole dataset. It
  returns \texttt{NA}. Work out why, then compute it in a way that
  succeeds, and add a comment saying what your version is actually the
  mean of.
\end{enumerate}

Step 6 is the one that matters. Every dataset you meet in this course
has missing values, and the difference between a mean of the observed
cases and a mean of the sample is a difference you will have to state
plainly in your final project.

\chapter*{Projects, Paths, and Reproducible
Habits}\label{projects-paths-and-reproducible-habits}
\addcontentsline{toc}{chapter}{Projects, Paths, and Reproducible Habits}

\markboth{Projects, Paths, and Reproducible Habits}{Projects, Paths, and
Reproducible Habits}

Code that runs on your machine and code that runs are not the same
thing. The gap between them is where most of the marks on your final
project's reproducibility criterion are lost, and almost all of it comes
down to a handful of habits that cost nothing once you have them.

This chapter is about those habits. None of the material is difficult.
What makes it worth a chapter is that the failures it prevents do not
appear until the moment they are most expensive, which is when someone
else opens your work.

\section*{The problem, stated plainly}\label{the-problem-stated-plainly}
\addcontentsline{toc}{section}{The problem, stated plainly}

\markright{The problem, stated plainly}

Suppose you write this at the top of an analysis.

\begin{Shaded}
\begin{Highlighting}[]
\FunctionTok{setwd}\NormalTok{(}\StringTok{"/Users/babak/Documents/AI course/week 4/analysis"}\NormalTok{)}
\FunctionTok{read.csv}\NormalTok{(}\StringTok{"civil\_war.csv"}\NormalTok{)}
\end{Highlighting}
\end{Shaded}

On your laptop, in that folder, it works. Everywhere else it fails,
because no other computer has that folder. Your instructor does not have
it, your co-author or colleague does not have it, and neither will you
after you reorganize your files in March.

\texttt{setwd()} is the most common single reason a submitted document
does not run. You will see it constantly in older code, including in
tutorials and in published replication files, and the instinct to copy
it is understandable. Do not. Nothing in this course requires it, and
using it may cost you points on the project.

The replacement is not a different function. It is a different way of
organizing work.

\section*{R Projects}\label{r-projects}
\addcontentsline{toc}{section}{R Projects}

\markright{R Projects}

An R Project is a folder that R treats as the home of a piece of work.
Creating one puts a small file with the extension \texttt{.Rproj} inside
the folder, and from then on, opening that file does three useful things
at once: it starts a fresh R session, it sets the working directory to
that folder, and it remembers which files you had open.

That second point is the whole trick. Because the working directory is
already correct, you never have to set it, and because it is set
relative to the project rather than to your machine, the same code works
on anyone else's computer.

To create one: \textbf{File \textgreater{} New Project \textgreater{}
New Directory \textgreater{} New Project}, then give the folder a name
and a location. For work you already have in a folder, use \textbf{File
\textgreater{} New Project \textgreater{} Existing Directory} instead.

From then on, open the project by double-clicking the \texttt{.Rproj}
file rather than by opening a script directly. If you open a
\texttt{.qmd} from the Finder, R starts with no project and the paths
will not resolve, which produces a ``cannot open file'' error that looks
like a problem with the file and is not.

\section*{Relative paths}\label{relative-paths}
\addcontentsline{toc}{section}{Relative paths}

\markright{Relative paths}

Inside a project, you refer to files by their location relative to the
project folder, with no leading slash and no drive letter.

\begin{Shaded}
\begin{Highlighting}[]
\NormalTok{dat }\OtherTok{\textless{}{-}} \FunctionTok{read.csv}\NormalTok{(}\StringTok{"data/civil\_war.csv"}\NormalTok{)}
\end{Highlighting}
\end{Shaded}

That line means ``the file \texttt{civil\_war.csv}, inside a folder
called \texttt{data}, inside whatever folder this project lives in.'' It
is true on every machine.

A few conventions worth following:

\begin{itemize}
\tightlist
\item
  Forward slashes, \texttt{/}, on every operating system, including
  Windows. R accepts them everywhere.
\item
  No spaces in file or folder names. Use \texttt{civil\_war.csv} rather
  than \texttt{civil\ war.csv}. Spaces are legal, but they will
  eventually bite you.
\item
  Lower case, consistently. Some systems treat \texttt{Data/} and
  \texttt{data/} as the same folder and some do not.
\end{itemize}

\section*{A layout that works}\label{a-layout-that-works}
\addcontentsline{toc}{section}{A layout that works}

\markright{A layout that works}

You do not need an elaborate structure. This is enough for the final
project and it will still make sense a year from now.

\begin{verbatim}
YourProject/
  YourProject.Rproj
  analysis.qmd          the document you submit
  data/
    raw/                exactly as downloaded, never edited
    processed/          anything your code produced
  R/
    clean-data.R        scripts that do a job
  figures/              plots your code saved
\end{verbatim}

The distinction between \texttt{data/raw/} and \texttt{data/processed/}
is the one that earns its keep. Raw files are never opened in Excel and
never edited by hand. Every change is made in code, which means every
change is visible, reversible, and explicable when a reviewer,
instructor, or colleague asks why your sample has 4,318 rows rather than
4,502.

\section*{Note on this course}\label{note-on-this-course}
\addcontentsline{toc}{section}{Note on this course}

\markright{Note on this course}

The handouts here load their data straight from a web address, so there
are no local files to find.

\begin{Shaded}
\begin{Highlighting}[]
\NormalTok{fl }\OtherTok{\textless{}{-}} \FunctionTok{read.csv}\NormalTok{(}\StringTok{"https://raw.githubusercontent.com/babakrezaee/AppliedAI{-}PoliSci/main/data/FearonLaitin2003.csv"}\NormalTok{)}
\end{Highlighting}
\end{Shaded}

That is deliberate. It means the code in every chapter runs for you
immediately, on any machine, with nothing downloaded.

Your final project is different. You will have data files of your own,
and the paths in your document will have to resolve on a machine that is
not yours. The habits in this chapter are for that, and it is suggested
to adopted for your final project.

\section*{Never save the workspace}\label{never-save-the-workspace}
\addcontentsline{toc}{section}{Never save the workspace}

\markright{Never save the workspace}

By default, R offers to save your objects when you quit and to reload
them when you return. Turn this off.

\textbf{Tools \textgreater{} Global Options \textgreater{} General},
then set \textbf{Save workspace to .RData on exit} to \textbf{Never},
and untick \textbf{Restore .RData into workspace at startup}.

The reason is that a restored workspace hides broken code. An object
created three weeks ago by a line you have since deleted is still
sitting in memory, so your analysis appears to work while depending on
something that no longer exists anywhere in your files. The failure
surfaces when you knit, or when someone else runs it, which is exactly
the wrong time.

Your objects should always be recreatable by running your code. If they
are not, the code is incomplete.

\section*{Restart, and run from clean}\label{restart-and-run-from-clean}
\addcontentsline{toc}{section}{Restart, and run from clean}

\markright{Restart, and run from clean}

The companion habit is to restart R often. \textbf{Session
\textgreater{} Restart R}, or \textbf{Cmd and Shift and F10} on a Mac,
\textbf{Ctrl and Shift and F10} on Windows.

After restarting, run your document from the top. If it completes, it
works. If it does not, you have found a dependency on something that
existed only in your session, and you have found it now rather than the
night before a deadline.

Do this before every submission, without exception. It is the entire
test, and it takes under a minute.

\section*{Packages}\label{packages}
\addcontentsline{toc}{section}{Packages}

\markright{Packages}

Load every package you use at the top of the document, in one place.

\begin{Shaded}
\begin{Highlighting}[]
\FunctionTok{library}\NormalTok{(tidyverse)}
\FunctionTok{library}\NormalTok{(caret)}
\end{Highlighting}
\end{Shaded}

\textbf{Do not put \texttt{install.packages()} in a document you
submit.} Installation is something you do once on your own machine, from
the console. Left in a document, it re-installs packages on the reader's
computer, which is slow, occasionally destructive, and never what they
wanted.

If a package is missing, the reader gets a clear error naming it, which
is the correct outcome. A short note in your README saying which
packages the project needs is a courtesy worth extending.

\section*{Seeds}\label{seeds}
\addcontentsline{toc}{section}{Seeds}

\markright{Seeds}

Anything involving randomness, which in this course means splitting data
into training and test sets, sampling, and fitting a random forest,
gives a different answer each time it runs.

\begin{Shaded}
\begin{Highlighting}[]
\FunctionTok{set.seed}\NormalTok{(}\DecValTok{7}\NormalTok{)}
\end{Highlighting}
\end{Shaded}

Put that line near the top, once, before anything random happens. The
number itself is arbitrary and does not matter. What matters is that
with it, your reported numbers reproduce exactly, and without it, a
reader who runs your code gets results that disagree with your text for
no stated reason.

\section*{When a path fails anyway}\label{when-a-path-fails-anyway}
\addcontentsline{toc}{section}{When a path fails anyway}

\markright{When a path fails anyway}

The error will look like this.

\begin{verbatim}
Warning: cannot open file 'data/civil_war.csv': No such file or directory
Error in file(file, "rt") : cannot open the connection
\end{verbatim}

Three checks, in order.

\textbf{Where does R think it is?}

\begin{Shaded}
\begin{Highlighting}[]
\FunctionTok{getwd}\NormalTok{()}
\end{Highlighting}
\end{Shaded}

If that is not your project folder, you opened the file without opening
the project. Close everything and open the \texttt{.Rproj} file.

\textbf{Does R see the file?}

\begin{Shaded}
\begin{Highlighting}[]
\FunctionTok{list.files}\NormalTok{(}\StringTok{"data"}\NormalTok{)}
\end{Highlighting}
\end{Shaded}

If the file is not listed, the problem is the name or the location, not
the code.

\textbf{Is the name exactly right?} Capital letters count. So does the
extension, and macOS hides extensions by default, so a file that looks
like \texttt{civil\_war.csv} in the Finder may actually be
\texttt{civil\_war.csv.txt}.

Notice that the fix is never \texttt{setwd()}.

\section*{Exercise}\label{exercise-2}
\addcontentsline{toc}{section}{Exercise}

\markright{Exercise}

About fifteen minutes.

\begin{enumerate}
\def\labelenumi{\arabic{enumi}.}
\item
  Create a new R Project called \texttt{ai-practice}, somewhere you will
  find it again.
\item
  Inside the project folder, create a \texttt{data} folder. In the
  RStudio Files pane, use the \textbf{New Folder} button.
\item
  Download the Fearon and Laitin file into it, using code rather than
  your browser:

\begin{Shaded}
\begin{Highlighting}[]
\FunctionTok{download.file}\NormalTok{(}
  \StringTok{"https://raw.githubusercontent.com/babakrezaee/AppliedAI{-}PoliSci/main/data/FearonLaitin2003.csv"}\NormalTok{,}
  \AttributeTok{destfile =} \StringTok{"data/FearonLaitin2003.csv"}
\NormalTok{)}
\end{Highlighting}
\end{Shaded}
\item
  Write a script that loads it with a relative path and reports its
  dimensions. Save the script as \texttt{explore.R}.
\item
  Set the workspace options as described above. Then restart R and run
  your script from the top. It should work with an empty environment.
\item
  Now break it on purpose. Change the path to
  \texttt{Data/FearonLaitin2003.csv}, with a capital D, and run it
  again. Read the error carefully, then run \texttt{getwd()} and
  \texttt{list.files("data")} and confirm what each one tells you.
\end{enumerate}

Step 6 is the useful one. That error message will find you again during
the block, and recognizing it in ten seconds rather than twenty minutes
is the only thing separating an evening's work from an evening's
frustration.

\chapter*{Reading Data into R}\label{reading-data-into-r}
\addcontentsline{toc}{chapter}{Reading Data into R}

\markboth{Reading Data into R}{Reading Data into R}

Importing data feels like the trivial step before the real work, and it
is where a surprising share of analytical errors are born. The reason is
that R almost never refuses to import a file. It makes a decision about
what your data mean, tells you nothing, and carries on. A column of
numbers arrives as text, a missing value code of \texttt{-99} is treated
as a real observation, and everything downstream is quietly wrong.

This chapter covers the formats you will actually meet and, more
importantly, the check that catches the failures.

\section*{CSV, the ordinary case}\label{csv-the-ordinary-case}
\addcontentsline{toc}{section}{CSV, the ordinary case}

\markright{CSV, the ordinary case}

A comma-separated (CSV) file is plain text with one row per line and
commas between fields. It is the format most political science data
arrive in, and R reads it out of the box.

\begin{Shaded}
\begin{Highlighting}[]
\NormalTok{fl }\OtherTok{\textless{}{-}} \FunctionTok{read.csv}\NormalTok{(}\StringTok{"data/FearonLaitin2003.csv"}\NormalTok{)}
\end{Highlighting}
\end{Shaded}

There is also a tidyverse version, \texttt{read\_csv()}, with an
underscore rather than a period.

\begin{Shaded}
\begin{Highlighting}[]
\FunctionTok{library}\NormalTok{(tidyverse)}
\NormalTok{fl }\OtherTok{\textless{}{-}} \FunctionTok{read\_csv}\NormalTok{(}\StringTok{"data/FearonLaitin2003.csv"}\NormalTok{)}
\end{Highlighting}
\end{Shaded}

Three differences matter in practice. \texttt{read\_csv()} is faster on
large files, it prints a summary of the column types it guessed, which
is useful information rather than clutter, and it returns a tibble
rather than a data frame. A tibble is a data frame with better printing
behavior, and everything you learned in the previous chapter applies to
it unchanged.

Either is acceptable in this course. The handouts use
\texttt{read.csv()} because it needs no package, so the first line of a
chapter always runs.

\section*{Straight from a web
address}\label{straight-from-a-web-address}
\addcontentsline{toc}{section}{Straight from a web address}

\markright{Straight from a web address}

Both functions accept a URL wherever they accept a filename, which means
data can be loaded with nothing downloaded.

\begin{Shaded}
\begin{Highlighting}[]
\NormalTok{fl }\OtherTok{\textless{}{-}} \FunctionTok{read.csv}\NormalTok{(}\StringTok{"https://raw.githubusercontent.com/babakrezaee/AppliedAI{-}PoliSci/main/data/FearonLaitin2003.csv"}\NormalTok{)}
\end{Highlighting}
\end{Shaded}

This is how every dataset in these handouts is loaded. It has two
advantages worth naming: the code runs identically on any machine, and
there is no file path to get wrong.

It has one disadvantage. If the address ever changes, the code stops
working, and if the file behind the address is edited, your results
change without your knowing. For the final project, download the file
once and put it on your cloud-based platform, then you have control over
your data source than relying on someone else's server.

Note the address on Github must be the \textbf{raw} one. Clicking a CSV
on GitHub shows a formatted preview page, and its address returns HTML
rather than data. The \textbf{Raw} button gives the address you want,
and it always begins \texttt{raw.githubusercontent.com}. You can also
upload your files on Dropbox, Googledrive, and Onedrive; just make sure
that you have the permission to put the files on these public
cloud-based platforms.

\section*{Check the import, every
time}\label{check-the-import-every-time}
\addcontentsline{toc}{section}{Check the import, every time}

\markright{Check the import, every time}

This is the habit the chapter exists for. Three lines, immediately after
loading anything.

\begin{Shaded}
\begin{Highlighting}[]
\FunctionTok{dim}\NormalTok{(fl)}
\FunctionTok{str}\NormalTok{(fl)}
\FunctionTok{summary}\NormalTok{(fl)}
\end{Highlighting}
\end{Shaded}

You are looking for three things.

\textbf{Did every column get the type you expected?} A numeric variable
listed as \texttt{chr} in the \texttt{str()} output has a non-numeric
value hiding in it somewhere.

\textbf{Are the row and column counts right?} If the file should have
6,610 rows and you have 6,609, the header row was probably read as data,
or the last line lacks a line break.

\textbf{Do the ranges make sense?} \texttt{summary()} reports the
minimum and maximum of each numeric column. A \texttt{polity2} score of
\texttt{-88} or a population of \texttt{-99} is a missing value code,
not an observation, and it will wreck any model that takes it seriously.

\section*{Missing value codes}\label{missing-value-codes}
\addcontentsline{toc}{section}{Missing value codes}

\markright{Missing value codes}

Older datasets frequently encode missing values as numbers, with
\texttt{-99}, \texttt{-999}, and \texttt{-88} being the common choices.
R has no way of knowing this and will treat them as real.

Tell it explicitly at import.

\begin{Shaded}
\begin{Highlighting}[]
\NormalTok{dat }\OtherTok{\textless{}{-}} \FunctionTok{read.csv}\NormalTok{(}\StringTok{"data/survey.csv"}\NormalTok{, }\AttributeTok{na.strings =} \FunctionTok{c}\NormalTok{(}\StringTok{"NA"}\NormalTok{, }\StringTok{""}\NormalTok{, }\StringTok{"{-}99"}\NormalTok{, }\StringTok{"{-}999"}\NormalTok{))}
\end{Highlighting}
\end{Shaded}

Or fix it afterward, which is clearer when only some columns are
affected.

\begin{Shaded}
\begin{Highlighting}[]
\NormalTok{dat}\SpecialCharTok{$}\NormalTok{income[dat}\SpecialCharTok{$}\NormalTok{income }\SpecialCharTok{==} \SpecialCharTok{{-}}\DecValTok{99}\NormalTok{] }\OtherTok{\textless{}{-}} \ConstantTok{NA}
\end{Highlighting}
\end{Shaded}

Either way, do it before anything else. A mean computed over a column
containing \texttt{-99} is not slightly wrong, it is meaningless, and
nothing in the output will tell you.

\section*{When the commas are not
commas}\label{when-the-commas-are-not-commas}
\addcontentsline{toc}{section}{When the commas are not commas}

\markright{When the commas are not commas}

Files exported from Excel on a European system are often separated by
semicolons rather than commas, and they use commas as decimal points.
Opening one with \texttt{read.csv()} produces a single column of text.

\begin{Shaded}
\begin{Highlighting}[]
\NormalTok{dat }\OtherTok{\textless{}{-}} \FunctionTok{read.csv2}\NormalTok{(}\StringTok{"data/dutch\_export.csv"}\NormalTok{)   }\CommentTok{\# semicolons, comma decimals}
\NormalTok{dat }\OtherTok{\textless{}{-}} \FunctionTok{read\_csv2}\NormalTok{(}\StringTok{"data/dutch\_export.csv"}\NormalTok{)   }\CommentTok{\# tidyverse equivalent}
\end{Highlighting}
\end{Shaded}

Worth knowing here rather than elsewhere, since Dutch-locale Excel does
this by default and several of you will hit it with data you export
yourselves.

If the separator is something else entirely, such as a tab, say so:

\begin{Shaded}
\begin{Highlighting}[]
\NormalTok{dat }\OtherTok{\textless{}{-}} \FunctionTok{read.delim}\NormalTok{(}\StringTok{"data/file.txt"}\NormalTok{)                    }\CommentTok{\# tab separated}
\NormalTok{dat }\OtherTok{\textless{}{-}} \FunctionTok{read.csv}\NormalTok{(}\StringTok{"data/file.txt"}\NormalTok{, }\AttributeTok{sep =} \StringTok{"|"}\NormalTok{)           }\CommentTok{\# anything else}
\end{Highlighting}
\end{Shaded}

The fastest way to diagnose any of this is to open the file in a plain
text editor and look at the first two lines. Two minutes of looking
beats an hour of guessing.

\section*{Accented characters}\label{accented-characters}
\addcontentsline{toc}{section}{Accented characters}

\markright{Accented characters}

Country names, party names, and any text data will contain characters
outside plain English. If they arrive mangled, as \texttt{é} where
\texttt{é} belongs, the file's encoding was misread.

\begin{Shaded}
\begin{Highlighting}[]
\NormalTok{dat }\OtherTok{\textless{}{-}} \FunctionTok{read.csv}\NormalTok{(}\StringTok{"data/parties.csv"}\NormalTok{, }\AttributeTok{encoding =} \StringTok{"UTF{-}8"}\NormalTok{)}
\end{Highlighting}
\end{Shaded}

UTF-8 is correct for the large majority of files you will meet. This
will matter in Week 4, when text becomes data.

\section*{Excel files}\label{excel-files}
\addcontentsline{toc}{section}{Excel files}

\markright{Excel files}

Excel files are not text and cannot be read by \texttt{read.csv()}. Use
\texttt{readxl}.

\begin{Shaded}
\begin{Highlighting}[]
\CommentTok{\# install.packages("readxl")}
\FunctionTok{library}\NormalTok{(readxl)}

\NormalTok{dat }\OtherTok{\textless{}{-}} \FunctionTok{read\_excel}\NormalTok{(}\StringTok{"data/data.xlsx"}\NormalTok{)}
\NormalTok{dat }\OtherTok{\textless{}{-}} \FunctionTok{read\_excel}\NormalTok{(}\StringTok{"data/data.xlsx"}\NormalTok{, }\AttributeTok{sheet =} \StringTok{"Coding"}\NormalTok{)   }\CommentTok{\# a named sheet}
\NormalTok{dat }\OtherTok{\textless{}{-}} \FunctionTok{read\_excel}\NormalTok{(}\StringTok{"data/data.xlsx"}\NormalTok{, }\AttributeTok{skip =} \DecValTok{3}\NormalTok{)           }\CommentTok{\# skip title rows}
\end{Highlighting}
\end{Shaded}

\texttt{skip} earns its place. Spreadsheets kept by humans very often
carry a title, a blank line, and a note before the actual header row,
and without \texttt{skip} all of that is read as data.

A note on Excel more generally. It is a reasonable place to receive data
and a poor place to edit it, because changes made by hand leave no
record. If a value must change, change it in R, where the change is
visible in your code and reversible.

\section*{Stata and SPSS}\label{stata-and-spss}
\addcontentsline{toc}{section}{Stata and SPSS}

\markright{Stata and SPSS}

Political science data arrive as \texttt{.dta} and \texttt{.sav} more
often than anyone would like. The \texttt{haven} package reads both.

\begin{Shaded}
\begin{Highlighting}[]
\CommentTok{\# install.packages("haven")}
\FunctionTok{library}\NormalTok{(haven)}

\NormalTok{dat }\OtherTok{\textless{}{-}} \FunctionTok{read\_dta}\NormalTok{(}\StringTok{"data/anes.dta"}\NormalTok{)   }\CommentTok{\# Stata}
\NormalTok{dat }\OtherTok{\textless{}{-}} \FunctionTok{read\_sav}\NormalTok{(}\StringTok{"data/ess.sav"}\NormalTok{)    }\CommentTok{\# SPSS}
\end{Highlighting}
\end{Shaded}

Older tutorials recommend \texttt{foreign} and \texttt{Hmisc} for this.
Both still work, and \texttt{haven} handles recent file versions and
preserves variable labels better, so prefer it.

Labelled columns from these formats sometimes behave unexpectedly in
models. If one gives you trouble, strip the labels:

\begin{Shaded}
\begin{Highlighting}[]
\NormalTok{dat}\SpecialCharTok{$}\NormalTok{party }\OtherTok{\textless{}{-}}\NormalTok{ haven}\SpecialCharTok{::}\FunctionTok{zap\_labels}\NormalTok{(dat}\SpecialCharTok{$}\NormalTok{party)}
\end{Highlighting}
\end{Shaded}

\section*{Saving}\label{saving}
\addcontentsline{toc}{section}{Saving}

\markright{Saving}

To hand data to someone who may not use R, write a CSV.

\begin{Shaded}
\begin{Highlighting}[]
\FunctionTok{write.csv}\NormalTok{(dat, }\StringTok{"data/processed/cleaned.csv"}\NormalTok{, }\AttributeTok{row.names =} \ConstantTok{FALSE}\NormalTok{)}
\end{Highlighting}
\end{Shaded}

\texttt{row.names\ =\ FALSE} matters. Without it, R adds a column of row
numbers, and a file that is written and read a few times accumulates
several of them.

To save an R object for your own later use, including model objects and
anything else a CSV cannot hold, use \texttt{saveRDS()}.

\begin{Shaded}
\begin{Highlighting}[]
\FunctionTok{saveRDS}\NormalTok{(model, }\StringTok{"output/model.rds"}\NormalTok{)}
\NormalTok{model }\OtherTok{\textless{}{-}} \FunctionTok{readRDS}\NormalTok{(}\StringTok{"output/model.rds"}\NormalTok{)}
\end{Highlighting}
\end{Shaded}

This is worth knowing for the final project, or algorithms that take a
longer time to run, such as Random Forest. A model that takes several
minutes to fit can be fitted once, saved, and read back, so that
re-rendering the document does not refit it.

\section*{Exercise}\label{exercise-3}
\addcontentsline{toc}{section}{Exercise}

\markright{Exercise}

About twenty minutes. Work in the \texttt{ai-practice} project from the
previous chapter.

\begin{enumerate}
\def\labelenumi{\arabic{enumi}.}
\item
  Load the Fearon and Laitin data from the raw URL and confirm its
  dimensions.
\item
  Run \texttt{str()} on it. List two columns that arrived as text. For
  each, say whether that is correct or whether it should have been
  numeric.
\item
  Run \texttt{summary(fl\$polity2)}. Look at the minimum. Polity scores
  are defined on a range of minus ten to plus ten. Is anything outside
  it, and if so, what do you think those values mean?
\item
  Make a deliberately awkward file. In the console:

\begin{Shaded}
\begin{Highlighting}[]
\NormalTok{messy }\OtherTok{\textless{}{-}} \FunctionTok{data.frame}\NormalTok{(}
  \AttributeTok{country =} \FunctionTok{c}\NormalTok{(}\StringTok{"França"}\NormalTok{, }\StringTok{"Türkiye"}\NormalTok{, }\StringTok{"España"}\NormalTok{),}
  \AttributeTok{score   =} \FunctionTok{c}\NormalTok{(}\FloatTok{1.5}\NormalTok{, }\SpecialCharTok{{-}}\DecValTok{99}\NormalTok{, }\FloatTok{3.2}\NormalTok{)}
\NormalTok{)}
\FunctionTok{write.csv}\NormalTok{(messy, }\StringTok{"data/messy.csv"}\NormalTok{, }\AttributeTok{row.names =} \ConstantTok{FALSE}\NormalTok{)}
\end{Highlighting}
\end{Shaded}

  Read it back with \texttt{read.csv("data/messy.csv")}, then compute
  \texttt{mean(score)}. Now read it again treating \texttt{-99} as
  missing, and compute the mean a second time. Report both numbers and
  say which one you would put in a paper.
\item
  Check whether the accented characters survived the round trip. If they
  did not, read the file again with \texttt{encoding\ =\ "UTF-8"}.
\end{enumerate}

Step 4 is the point. Two defensible-looking numbers come out of the same
file, and only one of them means anything. Nothing in R's output
distinguishes them for you.

\part{Weekly Handouts}

\chapter{Prediction, Causation, and the Vocabulary of
AI}\label{prediction-causation-and-the-vocabulary-of-ai}

\emph{Week 1, 2 September}

This first session has two purposes. First, we need to distinguish
between correlation, causation, prediction, and forecasting. These are
related, but they are not the same research goal, and confusing them can
produce bad research. Second, we need a working vocabulary for talking
about artificial intelligence. Terms such as machine learning,
generative AI, reasoning model, and agent are often used as though they
refer to the same thing. They do not.

There is no code in this handout. The work this week is conceptual. That
may sound less exciting than fitting a model, but it matters. A model
can be technically impressive and still answer the wrong question. The
same problem appears in public debates about AI, where a loose use of
one term can hide who designed a system, what it actually does, and who
is responsible when it fails.

Developing a theoretical argument about a political phenomenon is work
you do across the programme. This course adds three things to it. We ask
how evidence can be used to evaluate an argument. We ask how statistical
and machine-learning models can be used to anticipate outcomes. And we
treat AI systems themselves as objects of political analysis, because
they are increasingly part of how states, parties, platforms, and
administrations act.

\section{What kind of question are we
asking?}\label{what-kind-of-question-are-we-asking}

Political scientists use qualitative, quantitative, and mixed-method
research designs. All three can be empirical, and they share an
underlying scientific principle: theoretical claims should be exposed to
evidence that could count against them.

For Popper, a theory is scientific when it produces claims that are
empirically testable and open to falsification. This does not mean that
one study can prove a theory. Evidence can support a hypothesis without
proving that it is true. That support remains provisional and depends on
the quality of the research design, the measurement, and the
alternatives that have been considered.

Lakatos asks us to look beyond the fate of one hypothesis and consider
the development of a research programme (Lakatos 1970). Does it generate
new expectations and improve our understanding, or does it survive only
because its supporters keep adding qualifications after every failure?
We will return to this question when we compare models: improving
performance is useful, but adding complexity without learning anything
substantive is not necessarily scientific progress.

The same care is needed when discussing null hypotheses. Rejecting a
null hypothesis can provide support for a substantive alternative.
Failing to reject a null hypothesis is also evidence from the study and
can be informative, particularly when the design is strong and the test
has adequate power. It does not prove the null hypothesis. More
generally, when observations remain consistent with the implications of
a theory, we may describe this as empirical support, not as final
confirmation.

Publishing research is therefore an invitation to other scholars to
attack the argument scientifically: to examine the assumptions, repeat
the analysis, use other data, and try to show that the finding does not
hold.

The four questions below organize the first part of the course.

{\def\LTcaptype{none} % do not increment counter
\begin{longtable}[]{@{}
  >{\raggedright\arraybackslash}p{(\linewidth - 4\tabcolsep) * \real{0.3333}}
  >{\raggedright\arraybackslash}p{(\linewidth - 4\tabcolsep) * \real{0.3333}}
  >{\raggedright\arraybackslash}p{(\linewidth - 4\tabcolsep) * \real{0.3333}}@{}}
\toprule\noalign{}
\begin{minipage}[b]{\linewidth}\raggedright
Research goal
\end{minipage} & \begin{minipage}[b]{\linewidth}\raggedright
Main question
\end{minipage} & \begin{minipage}[b]{\linewidth}\raggedright
What the evidence can establish
\end{minipage} \\
\midrule\noalign{}
\endhead
\bottomrule\noalign{}
\endlastfoot
Correlation & Do two or more variables vary together? & The presence,
direction, and strength of an association \\
Causation & Would changing \(X\) change \(Y\), all else being equal? & A
causal effect, if the identification strategy is credible \\
Prediction & How well can we estimate an outcome for cases the model has
not seen? & Whether the available information generalizes to unseen
cases \\
Forecasting & How well can we anticipate outcomes that have not yet
occurred? & Whether a model can provide useful information about the
future \\
\end{longtable}
}

Forecasting is a form of prediction in a technical sense, but we keep it
as a separate category in this course. The temporal dimension creates
additional problems that deserve separate attention.

\subsection{Correlation}\label{correlation}

Sometimes our research is descriptive. We want to know whether two
variables are associated and how strong that association is. Descriptive
and correlational research has sometimes been downplayed following the
causal-inference revolution in economics and political science. That is
unfortunate. Description is often the first step in bringing a new
political phenomenon into view, and a reliable pattern can be important
even before its cause is established.

Consider the long-standing debate about economic development and
democracy. There is a well-known association between them: wealthier
countries tend, on average, to be more democratic. But that association
does not tell us which way the causal arrow runs. Economic development
may encourage democratization. Democracy may contribute to economic
development. A third factor may affect both, or the relationship may run
in both directions.

We can present the distinction schematically:

\textbf{Correlation}

\[Economic~development \overset{?; +/-}{\longleftrightarrow} Political~development\]

\textbf{Possible causal relationships}

\[Economic~development \overset{+/-}{\longrightarrow} Political~development\]

\[Economic~development \overset{+/-}{\longleftarrow} Political~development\]

The correlation is a finding. The arrows are additional claims that
require additional evidence.

\subsection{Causation}\label{causation}

A causal study needs an identification strategy: an argument about why
the estimated relationship can be interpreted as the effect of \(X\) on
\(Y\), rather than as selection, reverse causation, measurement error,
or the influence of an omitted factor. Calling \(X\) an independent
variable does not automatically make it exogenous. Exogeneity is an
assumption that must be defended through the design.

Bias is also worth defining carefully. In statistical terms, an
estimator is biased when its expected value differs systematically from
the parameter it is intended to estimate. This is not the same as
finding that one estimate is statistically significant. Omitted
variables, selection, and measurement problems are among the reasons an
estimate may fail to recover the causal quantity we care about.

\subsubsection{The experimental ideal}\label{the-experimental-ideal}

Randomized experiments provide a particularly clear identification
strategy. In one of my research projects, I examined whether information
about the principles of nonviolent resistance changes attitudes toward
this form of political action. Participants were randomly assigned to
two groups. One group watched a clip about International Relations. The
other watched a clip about nonviolent resistance. I then compared their
average attitudes after the intervention.

Random assignment means that treatment status is independent of
participants' pre-treatment characteristics by design. Across repeated
random assignments, the groups are comparable on average. In one
realized experiment, however, the groups do not have to be perfectly
identical on every characteristic. Covariate summaries can help us
describe the sample and identify surprising imbalances, but they do not
make the assignment random; the assignment procedure does that.

The comparison between the two groups gives us an estimate of the
average effect of receiving the information in that study. Notice how
much work the design is doing. The difference in outcomes becomes
interpretable because the intervention was assigned, rather than
selected by participants.

\subsubsection{Observational data}\label{observational-data}

Political scientists often cannot conduct experiments. Think about
conflict research. Can we randomly assign some countries to civil war
and others to peace in order to study the economic consequences? Of
course not.

We therefore rely on observational data produced by political and social
processes that we do not control. If we want to estimate the effect of
democracy on economic development, countries are not assigned randomly
to different levels of democracy. Democratic countries may also have
stronger state institutions, more secure property rights, different
colonial histories, or other characteristics that affect economic
performance.

Adding observed confounders to a model may reduce some of these
problems, but it does not automatically solve them. Matching,
instrumental variables, regression discontinuity, panel designs, and
other methods can support an identification strategy under particular
assumptions. Developing those strategies is outside the scope of this
course. Recognizing when a claim requires one is not.

\subsection{Prediction}\label{prediction}

Prediction changes the question. Given what we know about a case, how
accurately can we estimate its outcome?

Suppose we want to predict civil-war onset. Income, terrain, state
capacity, political institutions, conflict history, and regional
conditions may all carry useful information. A predictive model combines
that information and is evaluated by how well it performs on cases it
did not use for fitting.

This is different from estimating the causal effect of one variable. A
model may predict well even when some of its inputs are only proxies.
Imagine a bank using postal code to predict loan default. Postal code
may carry information about income, employment, housing, and ethnicity.
It can therefore be useful for prediction without providing a causal
explanation of why a person defaults. Intervening on the postal code
would not necessarily change the outcome.

Shmueli distinguishes explanatory modeling from predictive modeling
(Shmueli 2010). Cranmer and Desmarais argue that political science has
paid too little attention to predictive performance (Cranmer and
Desmarais 2017). Hofman, Sharma, and Watts make a related case for
taking prediction more seriously when evaluating explanations of social
phenomena (Hofman et al. 2017). None of this means that accuracy should
replace theory or causal identification.

Predictive performance can discipline an explanation. If an account
repeatedly fails to anticipate the phenomenon it claims to explain, we
have reason to examine it more closely. But predictive performance is
not the same thing as causal evidence, and strong prediction is not
sufficient for a good explanation. Social outcomes can also contain
substantial noise, which means that a credible causal explanation will
not always generate highly accurate individual predictions.

The purpose of the model therefore matters. If the task is prediction,
performance on unseen cases is central. Flexible functional forms may be
useful even when no single coefficient provides a simple substantive
interpretation. If the task is causal inference, good test-set accuracy
does not rescue a weak identification strategy.

\subsection{Forecasting}\label{forecasting}

Forecasting asks about outcomes that have not happened yet.

Suppose we have housing-price data from 1990 to 2026. We fit a model
using observations from 1990 to 2019 and test it on sales from 2020 to
2026. When the model estimates the price of a house sold in 2023, it is
making a prediction for a case whose outcome was hidden from the model
but is already known to us. We can calculate the error immediately.

If we use the same model to estimate prices in 2027 or 2028, we are
forecasting. The outcomes do not yet exist, so the forecast cannot be
fully evaluated until time passes.

When forecasting is the goal, we normally divide training and test data
by time rather than at random. A random split could allow a model
trained on 2025 observations to predict an outcome from 2010. That may
be acceptable for some prediction tasks, but it does not simulate the
information problem faced by someone trying to anticipate the future.

A temporal test set is still not a genuine forecast. It is a simulation
of forecasting using outcomes that are already known to the researcher.
It tells us how the model would have performed in a previous period. It
cannot guarantee that relationships will remain stable when the model is
used later. Early-warning systems such as ViEWS ({Hegre et al.} 2019)
produce forecasts in real time and can be fully judged only after the
forecast period has passed.

\section{Why does confusing these goals produce bad
research?}\label{why-does-confusing-these-goals-produce-bad-research}

There are several common mistakes.

\textbf{Causal language after a predictive model.} High accuracy shows
that the predictors contain information about the outcome. It does not
show that intervening on one of them would change the outcome. A random
forest does not turn an association into a causal effect.

\textbf{Predictive claims from a causal design.} A well-identified
estimate of one causal effect may be a poor predictor of the outcome as
a whole. The design was built to isolate an effect, not necessarily to
anticipate every case.

\textbf{Treating test-set performance as demonstrated forecast
performance.} Good performance on held-out historical years is useful
evidence. It is not proof that the model will work next year,
particularly when political conditions change.

\textbf{Applying the wrong evaluation standard.} Predictive models are
sometimes dismissed because they do not have easily interpretable
coefficients. Causal estimates are sometimes dismissed because they have
a low \(R^2\). Neither criticism is automatically relevant. We first
need to know what the model was built to do.

\section{What is under the hood of
AI?}\label{what-is-under-the-hood-of-ai}

For many contemporary AI systems, machine learning is the motor under
the hood. Credit scoring, content moderation, fraud detection,
political-text classification, recommendation systems, and large
language models all rely heavily on models fitted to data.

AI is nevertheless broader than machine learning. Some systems rely on
written rules, search, optimization, or combinations of these with
fitted models. It is therefore safer to say that machine learning powers
many of the AI systems we study, not that every AI system is machine
learning.

Before moving to specific algorithms, we need several basic terms.

\textbf{Features} are the inputs available to the model. \textbf{Labels}
are known outcomes attached to cases in the training data.
\textbf{Training} or \textbf{fitting} is the process of adjusting a
model's internal quantities so that its outputs correspond more closely
to the examples. A \textbf{test set} contains cases that were not used
for fitting and allows us to examine whether the model generalizes.

\subsection{Supervised, unsupervised, and self-supervised
learning}\label{supervised-unsupervised-and-self-supervised-learning}

\textbf{Supervised learning} uses labeled examples. If country-years in
the training data are labeled as war or peace, or parliamentary motions
are labeled by topic, the model learns a mapping from the features to
those known labels. The classifiers in Weeks 2 through 4 are
supervised-learning methods.

\textbf{Unsupervised learning} starts without an outcome label. It looks
for structure in the data, such as clusters of similar cases or
dimensions along which cases vary. Because there is no known correct
label, evaluation is less straightforward. We will discuss the
distinction even though the applied models in the first half of this
course are supervised.

\textbf{Self-supervised learning} is often treated as a branch of
unsupervised learning, although it creates a supervised-style prediction
problem from unlabeled data. The training data do not need labels
supplied by a human, because the data provide the training signal. In
language-model pretraining, parts of the text provide targets for
predicting other parts. Next-token prediction is the main example: the
preceding text is the input and the next token becomes the target.

Large language models are therefore commonly described as being
pretrained through self-supervised learning. They may later undergo
supervised fine-tuning, preference optimization, reinforcement learning,
or other forms of post-training.

Using an LLM to code or label political data is a separate matter. At
that point, we are using an already trained generative model to perform
a classification task. With examples in the prompt, it may act as a
few-shot classifier; without examples, it may act as a zero-shot
classifier. The task is classification, but that does not retroactively
change how the underlying model was pretrained.

\subsection{Generative, reasoning, and agentic
systems}\label{generative-reasoning-and-agentic-systems}

\textbf{Generative models} produce new content rather than only
assigning a predefined label. Large language models generate text by
estimating plausible continuations given the preceding context. Their
ability to produce fluent answers comes from large-scale pretraining and
later post-training, not from looking up a single stored response.

\textbf{Reasoning models} are language models optimized to perform
better on tasks that require multiple inferential steps, often through
additional training and more computation at inference time. Some produce
visible intermediate text before giving an answer. We should not assume
that this text is a literal transcript of an internal reasoning process.
It is model output and must be evaluated as such.

\textbf{Agentic systems} place a model inside a process in which it can
choose actions, use tools, inspect results, and continue working toward
a goal. The model may search, run code, retrieve documents, or call
other services. This expands what can be automated. It also makes
verification harder: if a system searched, selected, summarized, and
coded evidence for us, we must still decide what parts of that process
we can defend as researchers.

\section{Why is this vocabulary
political?}\label{why-is-this-vocabulary-political}

These distinctions are not only technical. They affect regulation,
responsibility, and the distribution of power.

A supervised classifier used to identify potentially fraudulent benefit
claims raises different problems from a generative model drafting
letters to claimants. The systems may differ in their training data,
output, error patterns, and opportunities for appeal. Regulation that
refers only to ``AI'' can miss those differences.

The source of a problem also affects who can address it. Unequal
treatment may arise from the training data, the objective used to fit
the model, post-training decisions, instructions written by the
deploying organization, or the inputs supplied in a particular case.
Responsibility may therefore be distributed across model developers,
vendors, public organizations, and individual users. Saying only that
``the algorithm is biased'' does not tell us which part of the system
produced the result.

The phrase ``the algorithm decided'' creates a related problem. An
organization selected the system, defined the purpose, chose the data
and threshold, and decided how the output would be used. Describing the
final outcome as the algorithm's decision can make those organizational
choices disappear.

\section{An older analogy: artificial political
entities}\label{an-older-analogy-artificial-political-entities}

Political theory and public-administration scholarship give us older
ways of thinking about artificial and collective decision-making. These
were not theories of artificial intelligence in the modern technical
sense. They are useful analogies because they direct our attention to
rules, information, authority, and responsibility.

On the first page of \emph{Leviathan}, Hobbes describes the commonwealth
as an ``Artificiall Man'' (Hobbes 1651). The sovereign is an artificial
person created through political agreement. The famous frontispiece
shows a body composed of many individual people. Hobbes is describing an
artificial political entity, not a machine-learning system. Still, the
image raises a familiar question: how can a constructed entity perceive,
decide, and act in ways that cannot be reduced to one individual inside
it?

Weberian bureaucracy offers a second analogy (Weber 1978). Bureaucracy
is designed to make decisions impersonally through offices, written
rules, files, and established procedures. Its political attraction is
predictability: similar cases should receive similar treatment
regardless of who happens to sit at the desk. In that limited sense, we
can think of bureaucracy as an algorithm implemented through people and
documents. But this remains an analogy. A written administrative rule
and a fitted statistical model are not the same thing. Or are they?

Herbert Simon provides the closest bridge (Simon 1947, 1969). He treated
organizations as information-processing and decision-making systems
while also contributing to the development of artificial intelligence.
Because individuals have limited information, attention, and
computational capacity, organizations use routines and decision
procedures to reach acceptable rather than perfectly optimal outcomes.

These analogies help us ask political questions about contemporary
systems. Who designed the procedure? What information can it use? What
is invisible to it? Who has the authority to override it? And who
answers for the result when it is wrong? The technology is new, but
several of the institutional questions are not.

\section{The plan for the block}\label{the-plan-for-the-block}

There are seven teaching weeks, each organized around one main idea.

Weeks 2 through 4 develop supervised classifiers, beginning with a
method simple enough to draw on the board and ending with political text
as data. Week 5 examines large language models, including
self-supervised pretraining, post-training, prompts, decoding, and the
use of LLMs for structured political analysis. Week 6 turns to
evaluation: confusion matrices, precision, recall, F1, calibration,
reliability, and the problems created by rare outcomes. Week 7 moves
from AI as a research method to AI as a political and administrative
object.

Basic evaluation begins in Week 2. From the first classifier onward, we
will use a training/test split, compare the model with a simple
baseline, and examine a confusion matrix. Week 6 returns to these ideas
in a more systematic way and adds calibration and the evaluation of
generative systems.

The central thread is the distinction introduced in this handout. You
will learn to build models, evaluate whether they work on unseen cases,
and state honestly what their results do and do not establish. You will
then use that technical understanding to examine political arguments
about AI.

\section{Reading}\label{reading}

Cranmer, S. J., and Desmarais, B. A. 2017. ``What Can We Learn from
Predictive Modeling?'' \emph{Political Analysis} 25(2): 145--166.
(Cranmer and Desmarais 2017)

Read the article for its argument rather than its technique. The authors
are responding to a discipline that has usually given explanation and
causal inference priority over prediction. Pay attention to what they
think predictive performance adds to political science, but also to what
they do not claim it can replace.

\section{Questions for the session}\label{questions-for-the-session}

Come with a view on at least two of these questions. They will structure
the first part of the session.

\begin{enumerate}
\def\labelenumi{\arabic{enumi}.}
\item
  What should count as empirical support for a theoretical claim? At
  what point, if any, does repeated support become confirmation?
\item
  Cranmer and Desmarais argue that political science has neglected
  prediction. What is the strongest objection to their position, and can
  it be answered?
\item
  Take a claim you believe about politics. Is it correlational, causal,
  predictive, or a forecast? What evidence would make you reconsider it?
\item
  A model predicts civil-war onset accurately using income, terrain, and
  state capacity. What has the model taught us about civil war, and what
  has it not taught us?
\item
  If a bureaucracy following written rules is an artificial decision
  system, what changes when the written rule is replaced by a fitted
  model? Try to identify a specific institutional or political
  difference rather than answering only that the model is more complex.
\item
  An LLM assigns political texts to ideological categories. Is this
  supervised, unsupervised, self-supervised, or something else? Does
  your answer refer to how the model was trained or to how it is now
  being used?
\end{enumerate}

\chapter{K-Nearest Neighbors}\label{k-nearest-neighbors}

\emph{Week 2, 9 September}

Last week we separated four research goals and left prediction as a
promise. This week we keep it. The method we start with, \(K\)-nearest
neighbors, was chosen for one reason above all others: it can be drawn.
There is no probability theory to follow, no optimization procedure
running out of sight, and no coefficients to interpret. Everything the
algorithm does can be seen in a picture, which means that when it goes
wrong, you can see that too.

That is worth more at this stage than a method with better performance.
You are going to spend the rest of the block working with models whose
internals are progressively harder to inspect. It helps to have one case
where nothing is hidden.

There is code in this handout, and you should run it rather than read
it. The workshop assumes you have.

\chapter{Fitting a function without assuming its
shape}\label{fitting-a-function-without-assuming-its-shape}

We usually want to know how an outcome varies with something else: how
\(Y\) behaves conditional on \(X\), written \(P(Y=y \mid X=x)\).

The most familiar strategy is to assume a shape for that relationship
and then estimate the parameters of the assumed shape. Ordinary least
squares assumes a line. Logit and probit assume particular shapes for
binary outcomes. These are \textbf{parametric} models, and the
assumption is doing a great deal of work.

Assuming a shape buys you real things. You get coefficients you can
interpret substantively, standard errors, a compact summary you can put
in a table, and a direct connection between a theoretical claim and a
quantity you can estimate. Most of what you have done in your other
courses relies on this.

It also costs something. If the true relationship is not the shape you
assumed, the model is wrong in a way that more data will not fix.
Collecting a thousand more observations does not rescue a straight line
drawn through a curve.

When the assumption looks doubtful, there are two responses. The first
is to stay inside the parametric framework and repair the model:
transform a variable, add an interaction, take logarithms. The second is
to use a \textbf{non-parametric} method, which assumes much less about
the shape and lets the data determine it.

\(K\)-nearest neighbors is the simplest non-parametric method there is.
Here is the entire algorithm:

\begin{quote}
To predict the outcome for a case, find the \(k\) cases in your data
most similar to it, and average their outcomes. If the outcome is a
category rather than a number, take the majority vote instead of the
average.
\end{quote}

There is nothing else. No estimation step, no parameters, no fitting in
the usual sense. The model is the data.

\chapter{Drawing it: development and infant
mortality}\label{drawing-it-development-and-infant-mortality}

Last week's running example was the association between economic
development and political development, and we stopped at the point where
the correlation was clear and the direction of the arrow was not. Set
the arrow aside for the moment. We are asking a prediction question now,
and prediction does not require us to settle it.

We will use national development indicators for 199 countries, which
ship with the \texttt{carData} package.

\begin{Shaded}
\begin{Highlighting}[]
\FunctionTok{install.packages}\NormalTok{(}\FunctionTok{c}\NormalTok{(}\StringTok{"carData"}\NormalTok{, }\StringTok{"kknn"}\NormalTok{, }\StringTok{"caret"}\NormalTok{, }\StringTok{"scatterplot3d"}\NormalTok{))}
\end{Highlighting}
\end{Shaded}

\begin{Shaded}
\begin{Highlighting}[]
\FunctionTok{library}\NormalTok{(carData)}
\FunctionTok{library}\NormalTok{(kknn)}

\FunctionTok{options}\NormalTok{(}\AttributeTok{scipen =} \DecValTok{999}\NormalTok{)   }\CommentTok{\# print numbers in full rather than as 1e+05}

\FunctionTok{data}\NormalTok{(UN, }\AttributeTok{package =} \StringTok{"carData"}\NormalTok{)}

\NormalTok{countries }\OtherTok{\textless{}{-}} \FunctionTok{na.omit}\NormalTok{(UN[, }\FunctionTok{c}\NormalTok{(}\StringTok{"ppgdp"}\NormalTok{, }\StringTok{"infantMortality"}\NormalTok{, }\StringTok{"fertility"}\NormalTok{,}
                            \StringTok{"pctUrban"}\NormalTok{, }\StringTok{"lifeExpF"}\NormalTok{, }\StringTok{"group"}\NormalTok{)])}

\DocumentationTok{\#\# Express per capita GDP in thousands of dollars}
\NormalTok{countries}\SpecialCharTok{$}\NormalTok{ppgdp }\OtherTok{\textless{}{-}}\NormalTok{ countries}\SpecialCharTok{$}\NormalTok{ppgdp }\SpecialCharTok{/} \DecValTok{1000}

\FunctionTok{nrow}\NormalTok{(countries)}
\end{Highlighting}
\end{Shaded}

\begin{verbatim}
[1] 193
\end{verbatim}

\begin{Shaded}
\begin{Highlighting}[]
\FunctionTok{head}\NormalTok{(countries)}
\end{Highlighting}
\end{Shaded}

\begin{verbatim}
             ppgdp infantMortality fertility pctUrban lifeExpF  group
Afghanistan 0.4990         124.535     5.968       23    49.49  other
Albania     3.6772          16.561     1.525       53    80.40  other
Algeria     4.4730          21.458     2.142       67    75.00 africa
Angola      4.3219          96.191     5.135       59    53.17 africa
Argentina   9.1621          12.337     2.172       93    79.89  other
Armenia     3.0307          24.272     1.735       64    77.33  other
\end{verbatim}

\texttt{ppgdp} is per capita gross domestic product and
\texttt{infantMortality} is deaths per thousand live births.

Two small things in that chunk. \texttt{options(scipen\ =\ 999)} tells R
to print numbers in full rather than in scientific notation, so that a
hundred thousand appears as \texttt{100000} and not as \texttt{1e+05}.
And we divide per capita GDP by a thousand, so it is measured in
thousands of dollars. That second step is cosmetic as far as the model
is concerned, since dividing every value by the same constant cannot
change which countries are nearest to which, but it keeps every axis in
this handout readable. Nobody reports per capita GDP to the nearest
dollar anyway.

Start by looking at the data.

\begin{Shaded}
\begin{Highlighting}[]
\FunctionTok{plot}\NormalTok{(countries}\SpecialCharTok{$}\NormalTok{ppgdp, countries}\SpecialCharTok{$}\NormalTok{infantMortality,}
     \AttributeTok{col =} \StringTok{"gray45"}\NormalTok{, }\AttributeTok{pch =} \DecValTok{19}\NormalTok{, }\AttributeTok{cex =} \FloatTok{0.7}\NormalTok{,}
     \AttributeTok{xlab =} \StringTok{"Per capita GDP (thousands of US$)"}\NormalTok{,}
     \AttributeTok{ylab =} \StringTok{"Infant mortality (per 1,000 live births)"}\NormalTok{)}
\end{Highlighting}
\end{Shaded}

\pandocbounded{\includegraphics[keepaspectratio]{handouts/w02-knn_files/figure-pdf/unnamed-chunk-3-1.pdf}}

The shape is the point. Infant mortality falls very steeply across the
poorest range and then flattens out almost completely. The difference
between a country at \$500 and a country at \$2,000 is enormous. The
difference between \$30,000 and \$40,000 is close to nothing.

Now put a straight line through it.

\begin{Shaded}
\begin{Highlighting}[]
\FunctionTok{plot}\NormalTok{(countries}\SpecialCharTok{$}\NormalTok{ppgdp, countries}\SpecialCharTok{$}\NormalTok{infantMortality,}
     \AttributeTok{col =} \StringTok{"gray45"}\NormalTok{, }\AttributeTok{pch =} \DecValTok{19}\NormalTok{, }\AttributeTok{cex =} \FloatTok{0.7}\NormalTok{,}
     \AttributeTok{xlab =} \StringTok{"Per capita GDP (thousands of US$)"}\NormalTok{,}
     \AttributeTok{ylab =} \StringTok{"Infant mortality (per 1,000 live births)"}\NormalTok{)}

\NormalTok{fit\_ols }\OtherTok{\textless{}{-}} \FunctionTok{lm}\NormalTok{(infantMortality }\SpecialCharTok{\textasciitilde{}}\NormalTok{ ppgdp, }\AttributeTok{data =}\NormalTok{ countries)}
\FunctionTok{abline}\NormalTok{(fit\_ols, }\AttributeTok{col =} \StringTok{"navy"}\NormalTok{, }\AttributeTok{lwd =} \FloatTok{2.5}\NormalTok{)}
\end{Highlighting}
\end{Shaded}

\pandocbounded{\includegraphics[keepaspectratio]{handouts/w02-knn_files/figure-pdf/unnamed-chunk-4-1.pdf}}

The line is not a small approximation to the truth here. It is
systematically wrong in a patterned way: too low for the poorest
countries, too high through the middle, and it eventually predicts
negative infant mortality, which is not a quantity that exists.

You could repair this without leaving the parametric framework, and in a
paper you probably would. Taking the logarithm of per capita GDP would
straighten most of the curve, and that transformation has a defensible
substantive reading: what matters for infant mortality is proportional
differences in income, not absolute ones. That is a real modeling
decision and a good one.

The point of this week is that \(K\)-nearest neighbors requires no such
decision. It does not need to be told the shape.

\begin{Shaded}
\begin{Highlighting}[]
\NormalTok{grid }\OtherTok{\textless{}{-}} \FunctionTok{data.frame}\NormalTok{(}\AttributeTok{ppgdp =} \FunctionTok{sort}\NormalTok{(countries}\SpecialCharTok{$}\NormalTok{ppgdp))}

\NormalTok{fit\_15 }\OtherTok{\textless{}{-}} \FunctionTok{kknn}\NormalTok{(infantMortality }\SpecialCharTok{\textasciitilde{}}\NormalTok{ ppgdp,}
               \AttributeTok{train =}\NormalTok{ countries, }\AttributeTok{test =}\NormalTok{ grid,}
               \AttributeTok{k =} \DecValTok{15}\NormalTok{, }\AttributeTok{kernel =} \StringTok{"rectangular"}\NormalTok{)}

\FunctionTok{plot}\NormalTok{(countries}\SpecialCharTok{$}\NormalTok{ppgdp, countries}\SpecialCharTok{$}\NormalTok{infantMortality,}
     \AttributeTok{col =} \StringTok{"gray45"}\NormalTok{, }\AttributeTok{pch =} \DecValTok{19}\NormalTok{, }\AttributeTok{cex =} \FloatTok{0.7}\NormalTok{,}
     \AttributeTok{xlab =} \StringTok{"Per capita GDP (thousands of US$)"}\NormalTok{,}
     \AttributeTok{ylab =} \StringTok{"Infant mortality (per 1,000 live births)"}\NormalTok{,}
     \AttributeTok{main =} \StringTok{"k = 15"}\NormalTok{)}

\FunctionTok{lines}\NormalTok{(grid}\SpecialCharTok{$}\NormalTok{ppgdp, fit\_15}\SpecialCharTok{$}\NormalTok{fitted.values, }\AttributeTok{col =} \StringTok{"maroon"}\NormalTok{, }\AttributeTok{lwd =} \DecValTok{3}\NormalTok{)}
\end{Highlighting}
\end{Shaded}

\pandocbounded{\includegraphics[keepaspectratio]{handouts/w02-knn_files/figure-pdf/unnamed-chunk-5-1.pdf}}

The argument \texttt{kernel\ =\ "rectangular"} means that all \(k\)
neighbors count equally. Other choices give closer neighbors more
weight. Rectangular is the plain version of the algorithm, and it is the
one described above.

\section{What the algorithm actually
did}\label{what-the-algorithm-actually-did}

Take one point on that curve. To produce a prediction for a country with
a per capita GDP of \$2,000, the algorithm found the fifteen countries
in the data closest to \$2,000 on that axis and averaged their infant
mortality. To produce the prediction at \$3,000, it did the same thing
with a neighborhood that had shifted slightly to the right, dropping
some countries and picking up others.

The curve is the trace left behind by that moving neighborhood. It is
not a formula, and it cannot be written as one. It is a record of what
the neighbors did.

This is also why the curve wobbles where the data are sparse. On the
right-hand side there are few countries, so the fifteen nearest
neighbors of a rich country may be spread across a very wide range of
incomes. The algorithm still returns an answer. It does not warn you
that the answer came from a neighborhood stretching across \$40,000.

\chapter{The choice of k}\label{the-choice-of-k}

Nothing so far told us to use fifteen. Try several values.

\begin{Shaded}
\begin{Highlighting}[]
\FunctionTok{plot}\NormalTok{(countries}\SpecialCharTok{$}\NormalTok{ppgdp, countries}\SpecialCharTok{$}\NormalTok{infantMortality,}
     \AttributeTok{col =} \StringTok{"gray75"}\NormalTok{, }\AttributeTok{pch =} \DecValTok{19}\NormalTok{, }\AttributeTok{cex =} \FloatTok{0.7}\NormalTok{,}
     \AttributeTok{xlab =} \StringTok{"Per capita GDP (thousands of US$)"}\NormalTok{,}
     \AttributeTok{ylab =} \StringTok{"Infant mortality (per 1,000 live births)"}\NormalTok{,}
     \AttributeTok{main =} \StringTok{"The same data, four values of k"}\NormalTok{)}

\NormalTok{kvals  }\OtherTok{\textless{}{-}} \FunctionTok{c}\NormalTok{(}\DecValTok{3}\NormalTok{, }\DecValTok{15}\NormalTok{, }\DecValTok{50}\NormalTok{, }\DecValTok{150}\NormalTok{)}
\NormalTok{kcols  }\OtherTok{\textless{}{-}} \FunctionTok{c}\NormalTok{(}\StringTok{"darkred"}\NormalTok{, }\StringTok{"maroon"}\NormalTok{, }\StringTok{"navy"}\NormalTok{, }\StringTok{"darkgreen"}\NormalTok{)}

\ControlFlowTok{for}\NormalTok{ (i }\ControlFlowTok{in} \FunctionTok{seq\_along}\NormalTok{(kvals)) \{}
\NormalTok{  f }\OtherTok{\textless{}{-}} \FunctionTok{kknn}\NormalTok{(infantMortality }\SpecialCharTok{\textasciitilde{}}\NormalTok{ ppgdp,}
            \AttributeTok{train =}\NormalTok{ countries, }\AttributeTok{test =}\NormalTok{ grid,}
            \AttributeTok{k =}\NormalTok{ kvals[i], }\AttributeTok{kernel =} \StringTok{"rectangular"}\NormalTok{)}
  \FunctionTok{lines}\NormalTok{(grid}\SpecialCharTok{$}\NormalTok{ppgdp, f}\SpecialCharTok{$}\NormalTok{fitted.values, }\AttributeTok{col =}\NormalTok{ kcols[i], }\AttributeTok{lwd =} \FloatTok{2.5}\NormalTok{)}
\NormalTok{\}}

\FunctionTok{legend}\NormalTok{(}\StringTok{"topright"}\NormalTok{, }\AttributeTok{legend =} \FunctionTok{paste}\NormalTok{(}\StringTok{"k ="}\NormalTok{, kvals),}
       \AttributeTok{col =}\NormalTok{ kcols, }\AttributeTok{lwd =} \FloatTok{2.5}\NormalTok{, }\AttributeTok{box.lty =} \DecValTok{0}\NormalTok{, }\AttributeTok{cex =} \FloatTok{0.8}\NormalTok{)}
\end{Highlighting}
\end{Shaded}

\pandocbounded{\includegraphics[keepaspectratio]{handouts/w02-knn_files/figure-pdf/unnamed-chunk-6-1.pdf}}

At \(k=3\) the curve is jagged. It follows every local fluctuation,
including fluctuations produced by the fact that we happen to have these
199 countries rather than some other set. At \(k=150\) the curve is
nearly flat, close to the overall average, and it has erased a
difference between rich and poor countries that is obviously real.

Both failures are informative, and they are opposite failures.

A small \(k\) takes every wiggle in the data seriously. Some of those
wiggles are the relationship we want. Others are noise. A model flexible
enough to reproduce the training data exactly has reproduced the noise
along with the signal, and the noise will not be there next time.

We can state this in one line. Write the outcome as

\[Y = \underbrace{f(X)}_{\text{signal}} + \underbrace{\varepsilon}_{\text{noise}}\]

What we want is \(f\). The term \(\varepsilon\) is whatever is left
over: measurement error, idiosyncrasy, the particular history of the
particular countries in this dataset. It is not a stable feature of the
world, and a model that has learned it has learned something that will
not recur.

So here is the sentence to take away from this week. \textbf{A model
that fits your data perfectly has told you nothing, because all it has
done is recite back the data you already had.} This is the opposite of
the instinct most people arrive with, which is that a closer fit is a
better model.

A large \(k\) fails in the other direction. It averages across countries
that are not comparable, and it produces a model too rigid to represent
a relationship that genuinely exists.

\chapter{Why a perfect fit should worry
you}\label{why-a-perfect-fit-should-worry-you}

To choose between these we need to stop grading the model on the data it
was built from. Split the data in two: fit on one part, evaluate on the
other.

\begin{Shaded}
\begin{Highlighting}[]
\FunctionTok{set.seed}\NormalTok{(}\DecValTok{7}\NormalTok{)}

\NormalTok{n        }\OtherTok{\textless{}{-}} \FunctionTok{nrow}\NormalTok{(countries)}
\NormalTok{train\_id }\OtherTok{\textless{}{-}} \FunctionTok{sample}\NormalTok{(}\DecValTok{1}\SpecialCharTok{:}\NormalTok{n, }\FunctionTok{floor}\NormalTok{(}\FloatTok{0.7} \SpecialCharTok{*}\NormalTok{ n))}

\NormalTok{Train }\OtherTok{\textless{}{-}}\NormalTok{ countries[train\_id, ]}
\NormalTok{Test  }\OtherTok{\textless{}{-}}\NormalTok{ countries[}\SpecialCharTok{{-}}\NormalTok{train\_id, ]}

\FunctionTok{nrow}\NormalTok{(Train)}
\end{Highlighting}
\end{Shaded}

\begin{verbatim}
[1] 135
\end{verbatim}

\begin{Shaded}
\begin{Highlighting}[]
\FunctionTok{nrow}\NormalTok{(Test)}
\end{Highlighting}
\end{Shaded}

\begin{verbatim}
[1] 58
\end{verbatim}

\texttt{set.seed()} fixes the random draw so that you and I get the same
split. Use it every time you do anything random. Without it your results
change on every run and neither you nor anyone else can reproduce them.

Now measure the error. For a continuous outcome the usual summary is the
root mean squared error,

\[\text{RMSE} = \sqrt{\frac{1}{n}\sum_{i=1}^{n}\left[y_i - \hat{f}(x_i)\right]^2},\]

which is roughly the size of a typical prediction error, in the units of
the outcome. Here that means deaths per thousand live births, so an RMSE
of 15 means the model is typically off by about fifteen deaths per
thousand.

\begin{Shaded}
\begin{Highlighting}[]
\NormalTok{rmse }\OtherTok{\textless{}{-}} \ControlFlowTok{function}\NormalTok{(y, yhat) }\FunctionTok{sqrt}\NormalTok{(}\FunctionTok{mean}\NormalTok{((y }\SpecialCharTok{{-}}\NormalTok{ yhat)}\SpecialCharTok{\^{}}\DecValTok{2}\NormalTok{))}

\NormalTok{kseq     }\OtherTok{\textless{}{-}} \DecValTok{2}\SpecialCharTok{:}\DecValTok{60}
\NormalTok{err\_test }\OtherTok{\textless{}{-}} \FunctionTok{rep}\NormalTok{(}\ConstantTok{NA}\NormalTok{, }\FunctionTok{length}\NormalTok{(kseq))}
\NormalTok{err\_train }\OtherTok{\textless{}{-}} \FunctionTok{rep}\NormalTok{(}\ConstantTok{NA}\NormalTok{, }\FunctionTok{length}\NormalTok{(kseq))}

\ControlFlowTok{for}\NormalTok{ (i }\ControlFlowTok{in} \FunctionTok{seq\_along}\NormalTok{(kseq)) \{}
\NormalTok{  f\_test }\OtherTok{\textless{}{-}} \FunctionTok{kknn}\NormalTok{(infantMortality }\SpecialCharTok{\textasciitilde{}}\NormalTok{ ppgdp, Train, Test,}
                 \AttributeTok{k =}\NormalTok{ kseq[i], }\AttributeTok{kernel =} \StringTok{"rectangular"}\NormalTok{)}
\NormalTok{  err\_test[i] }\OtherTok{\textless{}{-}} \FunctionTok{rmse}\NormalTok{(Test}\SpecialCharTok{$}\NormalTok{infantMortality, f\_test}\SpecialCharTok{$}\NormalTok{fitted.values)}

\NormalTok{  f\_train }\OtherTok{\textless{}{-}} \FunctionTok{kknn}\NormalTok{(infantMortality }\SpecialCharTok{\textasciitilde{}}\NormalTok{ ppgdp, Train, Train,}
                  \AttributeTok{k =}\NormalTok{ kseq[i], }\AttributeTok{kernel =} \StringTok{"rectangular"}\NormalTok{)}
\NormalTok{  err\_train[i] }\OtherTok{\textless{}{-}} \FunctionTok{rmse}\NormalTok{(Train}\SpecialCharTok{$}\NormalTok{infantMortality, f\_train}\SpecialCharTok{$}\NormalTok{fitted.values)}
\NormalTok{\}}

\FunctionTok{plot}\NormalTok{(kseq, err\_test, }\AttributeTok{type =} \StringTok{"l"}\NormalTok{, }\AttributeTok{col =} \StringTok{"maroon"}\NormalTok{, }\AttributeTok{lwd =} \FloatTok{2.5}\NormalTok{,}
     \AttributeTok{ylim =} \FunctionTok{range}\NormalTok{(}\FunctionTok{c}\NormalTok{(err\_test, err\_train)),}
     \AttributeTok{xlab =} \StringTok{"k"}\NormalTok{, }\AttributeTok{ylab =} \StringTok{"RMSE"}\NormalTok{,}
     \AttributeTok{main =} \StringTok{"Error on data the model has seen, and data it has not"}\NormalTok{)}

\FunctionTok{lines}\NormalTok{(kseq, err\_train, }\AttributeTok{col =} \StringTok{"navy"}\NormalTok{, }\AttributeTok{lwd =} \FloatTok{2.5}\NormalTok{)}

\FunctionTok{legend}\NormalTok{(}\StringTok{"bottomright"}\NormalTok{, }\AttributeTok{legend =} \FunctionTok{c}\NormalTok{(}\StringTok{"Test (unseen)"}\NormalTok{, }\StringTok{"Training (seen)"}\NormalTok{),}
       \AttributeTok{col =} \FunctionTok{c}\NormalTok{(}\StringTok{"maroon"}\NormalTok{, }\StringTok{"navy"}\NormalTok{), }\AttributeTok{lwd =} \FloatTok{2.5}\NormalTok{, }\AttributeTok{box.lty =} \DecValTok{0}\NormalTok{, }\AttributeTok{cex =} \FloatTok{0.8}\NormalTok{)}
\end{Highlighting}
\end{Shaded}

\pandocbounded{\includegraphics[keepaspectratio]{handouts/w02-knn_files/figure-pdf/unnamed-chunk-8-1.pdf}}

\begin{Shaded}
\begin{Highlighting}[]
\NormalTok{kseq[}\FunctionTok{which.min}\NormalTok{(err\_test)]}
\end{Highlighting}
\end{Shaded}

\begin{verbatim}
[1] 17
\end{verbatim}

Read the two curves against each other, because the comparison is the
lesson.

The training error keeps falling as \(k\) gets smaller. It is lowest at
the far left, where the model is most flexible. If you had only this
curve, you would conclude that the best model is the most complicated
one available, and you would conclude that every single time, on every
dataset, regardless of the truth.

The test error does something different. It falls, reaches a minimum,
and then rises again. Past that minimum the model is no longer learning
the relationship. It is learning the training countries.

One caution before you go looking for the best \(k\). Change the seed
and rerun the code. The U shape will still be there; the exact location
of the minimum will move. The shape is a property of the problem. The
precise minimum is a property of one arbitrary split. Do not report
``the optimal \(k\) is 11'' as though you had discovered something about
infant mortality. We will fix this properly in a moment.

\section{The same picture, drawn the conventional
way}\label{the-same-picture-drawn-the-conventional-way}

That plot has an awkward feature: its horizontal axis runs backwards. A
large \(k\) is a \emph{simple} model, so complexity increases as you
move to the \emph{left}, which is the opposite of how you will see this
drawn everywhere else.

The usual fix is to plot against \(\log(1/k)\), which increases as the
model gets more flexible.

\begin{Shaded}
\begin{Highlighting}[]
\FunctionTok{plot}\NormalTok{(}\FunctionTok{log}\NormalTok{(}\DecValTok{1} \SpecialCharTok{/}\NormalTok{ kseq), err\_test, }\AttributeTok{type =} \StringTok{"l"}\NormalTok{, }\AttributeTok{col =} \StringTok{"maroon"}\NormalTok{, }\AttributeTok{lwd =} \FloatTok{2.5}\NormalTok{,}
     \AttributeTok{ylim =} \FunctionTok{range}\NormalTok{(}\FunctionTok{c}\NormalTok{(err\_test, err\_train)),}
     \AttributeTok{xlab =} \StringTok{"Model complexity, log(1/k)"}\NormalTok{, }\AttributeTok{ylab =} \StringTok{"RMSE"}\NormalTok{,}
     \AttributeTok{main =} \StringTok{"Complexity increases to the right"}\NormalTok{)}

\FunctionTok{lines}\NormalTok{(}\FunctionTok{log}\NormalTok{(}\DecValTok{1} \SpecialCharTok{/}\NormalTok{ kseq), err\_train, }\AttributeTok{col =} \StringTok{"navy"}\NormalTok{, }\AttributeTok{lwd =} \FloatTok{2.5}\NormalTok{)}

\FunctionTok{text}\NormalTok{(}\FunctionTok{log}\NormalTok{(}\DecValTok{1} \SpecialCharTok{/} \DecValTok{60}\NormalTok{), err\_test[}\FunctionTok{length}\NormalTok{(kseq)], }\StringTok{"k = 60"}\NormalTok{, }\AttributeTok{pos =} \DecValTok{4}\NormalTok{, }\AttributeTok{cex =} \FloatTok{0.8}\NormalTok{, }\AttributeTok{col =} \StringTok{"gray30"}\NormalTok{)}
\FunctionTok{text}\NormalTok{(}\FunctionTok{log}\NormalTok{(}\DecValTok{1} \SpecialCharTok{/} \DecValTok{2}\NormalTok{),  err\_test[}\DecValTok{1}\NormalTok{],            }\StringTok{"k = 2"}\NormalTok{,  }\AttributeTok{pos =} \DecValTok{2}\NormalTok{, }\AttributeTok{cex =} \FloatTok{0.8}\NormalTok{, }\AttributeTok{col =} \StringTok{"gray30"}\NormalTok{)}

\FunctionTok{legend}\NormalTok{(}\StringTok{"top"}\NormalTok{, }\AttributeTok{legend =} \FunctionTok{c}\NormalTok{(}\StringTok{"Test (unseen)"}\NormalTok{, }\StringTok{"Training (seen)"}\NormalTok{),}
       \AttributeTok{col =} \FunctionTok{c}\NormalTok{(}\StringTok{"maroon"}\NormalTok{, }\StringTok{"navy"}\NormalTok{), }\AttributeTok{lwd =} \FloatTok{2.5}\NormalTok{, }\AttributeTok{box.lty =} \DecValTok{0}\NormalTok{, }\AttributeTok{cex =} \FloatTok{0.8}\NormalTok{)}
\end{Highlighting}
\end{Shaded}

\pandocbounded{\includegraphics[keepaspectratio]{handouts/w02-knn_files/figure-pdf/unnamed-chunk-9-1.pdf}}

Identical numbers, identical conclusion, drawn so that ``further right''
means ``more flexible model.'' Every version of this figure you meet
from now on, in any textbook and for any method, has this orientation,
and the two curves always behave this way: the training curve slides
downward forever, and the test curve turns around.

\chapter{Bias and variance}\label{bias-and-variance}

We now have the picture. The vocabulary for it is worth learning,
because it is how everyone else in this literature talks, and because
the two failures it names are genuinely different failures.

Return to the decomposition from earlier:

\[Y = \underbrace{f(X)}_{\text{signal}} + \underbrace{\varepsilon}_{\text{noise}}\]

A model can be wrong about \(f\) in two ways.

\textbf{Bias} is systematic error. The model is too rigid to represent
the shape of \(f\), so it is wrong in the same direction every time, and
it would still be wrong if you handed it a million more countries. The
straight line we drew through the infant mortality data has high bias.
So does \(K\)-nearest neighbors at \(k=150\), which flattens a real
curve into something close to a horizontal line.

\textbf{Variance} is instability. The model is flexible enough to
represent \(f\), but it responds so eagerly to the particular cases in
front of it that a different sample would have produced a noticeably
different answer. It is not systematically wrong. It is just not
reliable, and the specific answer you got is partly an accident of the
data you happen to hold.

The trade-off is that reducing one usually increases the other, which is
what produces the U.

\begin{figure}[H]

{\centering \pandocbounded{\includegraphics[keepaspectratio]{handouts/../images/bias-variance-tradeoff.png}}

}

\caption{Test and training error as a function of model complexity.
Reproduced from Hastie, T., Tibshirani, R., and Friedman, J., \emph{The
Elements of Statistical Learning}, 2nd edition, Springer, 2009.}

\end{figure}%

This is a schematic, not a fit to anything. It is the shape you just
produced from real data, and the shape you will produce from almost any
dataset with almost any flexible method.

Read the two arrows at the top. Moving left, the model gets simpler:
variance falls, bias rises, and eventually the model is too rigid to
represent what is there. Moving right, the model gets more flexible:
bias falls, variance rises, and eventually the model is chasing noise.
The training curve slides downward the whole way and never turns, which
is why it cannot be used to choose anything. The test curve has a
minimum, and finding it is the entire problem.

\section{Seeing variance directly}\label{seeing-variance-directly}

Bias and variance are hard to separate with real data, for a simple
reason: you would need to know the true \(f\) in order to say how far
the model sits from it, and if you knew \(f\) you would not be modeling.

So we simulate. We invent a function, add noise to it, draw many samples
from it, and fit a model to each. Because we wrote down the truth
ourselves, we can measure exactly how far off each fit lands.

\begin{Shaded}
\begin{Highlighting}[]
\FunctionTok{set.seed}\NormalTok{(}\DecValTok{1364}\NormalTok{)}

\NormalTok{N }\OtherTok{\textless{}{-}} \DecValTok{1000}
\NormalTok{x }\OtherTok{\textless{}{-}} \FunctionTok{rnorm}\NormalTok{(N, }\DecValTok{2}\NormalTok{, }\DecValTok{1}\NormalTok{)}
\NormalTok{f\_true }\OtherTok{\textless{}{-}} \ControlFlowTok{function}\NormalTok{(x) }\FunctionTok{cospi}\NormalTok{(x)}
\NormalTok{y }\OtherTok{\textless{}{-}} \FunctionTok{f\_true}\NormalTok{(x) }\SpecialCharTok{+} \FunctionTok{rnorm}\NormalTok{(N, }\DecValTok{0}\NormalTok{, }\FloatTok{0.5}\NormalTok{)}

\NormalTok{sim   }\OtherTok{\textless{}{-}} \FunctionTok{data.frame}\NormalTok{(x, y)}
\NormalTok{grid2 }\OtherTok{\textless{}{-}} \FunctionTok{data.frame}\NormalTok{(}\AttributeTok{x =} \FunctionTok{sort}\NormalTok{(x))}

\NormalTok{kvec   }\OtherTok{\textless{}{-}} \FunctionTok{c}\NormalTok{(}\DecValTok{5}\NormalTok{, }\DecValTok{40}\NormalTok{, }\DecValTok{150}\NormalTok{)}
\NormalTok{nsim   }\OtherTok{\textless{}{-}} \DecValTok{30}
\NormalTok{ntrain }\OtherTok{\textless{}{-}} \DecValTok{200}
\NormalTok{spot   }\OtherTok{\textless{}{-}} \DecValTok{400}                       \CommentTok{\# one x value we will watch closely}
\NormalTok{x\_spot }\OtherTok{\textless{}{-}}\NormalTok{ grid2}\SpecialCharTok{$}\NormalTok{x[spot]}

\FunctionTok{par}\NormalTok{(}\AttributeTok{mfrow =} \FunctionTok{c}\NormalTok{(}\DecValTok{2}\NormalTok{, }\DecValTok{3}\NormalTok{))}

\NormalTok{preds }\OtherTok{\textless{}{-}} \FunctionTok{matrix}\NormalTok{(}\ConstantTok{NA}\NormalTok{, nsim, }\FunctionTok{length}\NormalTok{(kvec))}

\ControlFlowTok{for}\NormalTok{ (j }\ControlFlowTok{in} \FunctionTok{seq\_along}\NormalTok{(kvec)) \{}

  \FunctionTok{plot}\NormalTok{(x, y, }\AttributeTok{col =} \StringTok{"gray70"}\NormalTok{, }\AttributeTok{pch =} \DecValTok{1}\NormalTok{, }\AttributeTok{cex =} \FloatTok{0.5}\NormalTok{,}
       \AttributeTok{main =} \FunctionTok{paste}\NormalTok{(}\StringTok{"k ="}\NormalTok{, kvec[j]), }\AttributeTok{font.main =} \DecValTok{1}\NormalTok{)}

  \ControlFlowTok{for}\NormalTok{ (i }\ControlFlowTok{in} \DecValTok{1}\SpecialCharTok{:}\NormalTok{nsim) \{}
\NormalTok{    draw }\OtherTok{\textless{}{-}}\NormalTok{ sim[}\FunctionTok{sample}\NormalTok{(}\DecValTok{1}\SpecialCharTok{:}\NormalTok{N, ntrain), ]}
\NormalTok{    f    }\OtherTok{\textless{}{-}} \FunctionTok{kknn}\NormalTok{(y }\SpecialCharTok{\textasciitilde{}}\NormalTok{ x, draw, grid2, }\AttributeTok{k =}\NormalTok{ kvec[j], }\AttributeTok{kernel =} \StringTok{"rectangular"}\NormalTok{)}
    \FunctionTok{lines}\NormalTok{(grid2}\SpecialCharTok{$}\NormalTok{x, f}\SpecialCharTok{$}\NormalTok{fitted.values, }\AttributeTok{col =} \StringTok{"royalblue2"}\NormalTok{, }\AttributeTok{lwd =} \FloatTok{0.4}\NormalTok{)}
\NormalTok{    preds[i, j] }\OtherTok{\textless{}{-}}\NormalTok{ f}\SpecialCharTok{$}\NormalTok{fitted.values[spot]}
\NormalTok{  \}}

  \FunctionTok{abline}\NormalTok{(}\AttributeTok{v =}\NormalTok{ x\_spot, }\AttributeTok{col =} \StringTok{"gray40"}\NormalTok{, }\AttributeTok{lty =} \DecValTok{2}\NormalTok{)}
\NormalTok{\}}

\ControlFlowTok{for}\NormalTok{ (j }\ControlFlowTok{in} \FunctionTok{seq\_along}\NormalTok{(kvec)) \{}
  \FunctionTok{boxplot}\NormalTok{(preds[, j], }\AttributeTok{ylim =} \FunctionTok{range}\NormalTok{(preds), }\AttributeTok{col =} \StringTok{"gray90"}\NormalTok{,}
          \AttributeTok{main =} \FunctionTok{paste}\NormalTok{(}\StringTok{"Predictions at one point, k ="}\NormalTok{, kvec[j]),}
          \AttributeTok{font.main =} \DecValTok{1}\NormalTok{, }\AttributeTok{cex.main =} \FloatTok{0.9}\NormalTok{)}
  \FunctionTok{abline}\NormalTok{(}\AttributeTok{h =} \FunctionTok{f\_true}\NormalTok{(x\_spot), }\AttributeTok{col =} \StringTok{"red"}\NormalTok{, }\AttributeTok{lwd =} \DecValTok{2}\NormalTok{)}
\NormalTok{\}}
\end{Highlighting}
\end{Shaded}

\pandocbounded{\includegraphics[keepaspectratio]{handouts/w02-knn_files/figure-pdf/unnamed-chunk-10-1.pdf}}

Thirty samples, three values of \(k\), and the same underlying truth
throughout. The red line in the lower row marks the correct answer at
the point we are watching.

At \(k=5\) the fitted curves scatter all over the plot. Each one is
chasing its own two hundred observations. The predictions center roughly
on the truth, so the bias is small, but any single one of them could be
well off. That is variance.

At \(k=150\) the curves lie almost on top of one another. Run the study
again and you would get essentially the same answer, which sounds like a
virtue until you notice where that answer sits: the box is nowhere near
the red line. The model is stable and consistently wrong. That is bias.

At \(k=40\) the box is both reasonably tight and reasonably close. That
is the whole idea.

Notice what this means for how you read a result. A model with high
variance can be badly wrong without giving any sign of it, because the
fit to the data in front of you looks excellent. Nothing in the output
announces that a different sample would have said something else.

\chapter{Cross-validation}\label{cross-validation}

Now back to the problem we left hanging. We chose \(k\) using one train
and test split, and the choice moved when the seed moved. Two complaints
about the single split, both fair:

\begin{enumerate}
\def\labelenumi{\arabic{enumi}.}
\tightlist
\item
  The answer depends on which split we happened to draw, which is
  exactly the instability we just spent a section learning to distrust.
\item
  We set aside thirty percent of the countries and never trained on
  them. With 199 cases that is a real loss of information.
\end{enumerate}

\textbf{Cross-validation} addresses both by rotating the role of the
held-out set.

\section{Leave-one-out}\label{leave-one-out}

Set aside a single observation, train on the remaining \(n-1\), predict
the one you left out, and record the error. Do it again with the next
observation. Repeat \(n\) times and average.

\[\text{RMSE}_{\text{LOO}} = \sqrt{\frac{1}{n}\sum_{i=1}^{n}\left(y_i - \hat{y}_i\right)^2}\]

Every case gets used for training almost every time and for testing
exactly once. Nothing is wasted, and there is no arbitrary split to be
at the mercy of. The cost is that you have to fit the model \(n\) times.

\section{k-fold}\label{k-fold}

The practical compromise is to hold out a group instead of a single
case. Split the data at random into a number of roughly equal folds.
Hold out the first fold, train on the rest, and record the error on the
fold you held out. Then hold out the second fold, and so on, until every
fold has served once as the test set. Average the results.

\pandocbounded{\includegraphics[keepaspectratio]{handouts/w02-knn_files/figure-pdf/unnamed-chunk-11-1.pdf}}

\[\text{CV RMSE} = \sqrt{\frac{1}{\text{nfold}}\sum_{f=1}^{\text{nfold}} \text{MSE}_f}\]

\textbf{One warning about names before the code.} The \(k\) in
``\(k\)-fold'' has nothing whatever to do with the \(k\) in
``\(K\)-nearest neighbors.'' They are two unrelated numbers that
unfortunately share a letter, and we are about to choose one of them
using the other. To keep them apart, the code below calls the number of
folds \texttt{nfold} and never \texttt{k}.

\begin{Shaded}
\begin{Highlighting}[]
\NormalTok{cv\_knn }\OtherTok{\textless{}{-}} \ControlFlowTok{function}\NormalTok{(data, kvals, }\AttributeTok{nfold =} \DecValTok{10}\NormalTok{, }\AttributeTok{seed =} \DecValTok{7}\NormalTok{) \{}

  \FunctionTok{set.seed}\NormalTok{(seed)}
\NormalTok{  fold }\OtherTok{\textless{}{-}} \FunctionTok{sample}\NormalTok{(}\FunctionTok{rep}\NormalTok{(}\DecValTok{1}\SpecialCharTok{:}\NormalTok{nfold, }\AttributeTok{length.out =} \FunctionTok{nrow}\NormalTok{(data)))}
\NormalTok{  out  }\OtherTok{\textless{}{-}} \FunctionTok{rep}\NormalTok{(}\ConstantTok{NA}\NormalTok{, }\FunctionTok{length}\NormalTok{(kvals))}

  \ControlFlowTok{for}\NormalTok{ (j }\ControlFlowTok{in} \FunctionTok{seq\_along}\NormalTok{(kvals)) \{}

\NormalTok{    fold\_mse }\OtherTok{\textless{}{-}} \FunctionTok{rep}\NormalTok{(}\ConstantTok{NA}\NormalTok{, nfold)}

    \ControlFlowTok{for}\NormalTok{ (f }\ControlFlowTok{in} \DecValTok{1}\SpecialCharTok{:}\NormalTok{nfold) \{}
\NormalTok{      tr  }\OtherTok{\textless{}{-}}\NormalTok{ data[fold }\SpecialCharTok{!=}\NormalTok{ f, ]}
\NormalTok{      te  }\OtherTok{\textless{}{-}}\NormalTok{ data[fold }\SpecialCharTok{==}\NormalTok{ f, ]}
\NormalTok{      fit }\OtherTok{\textless{}{-}} \FunctionTok{kknn}\NormalTok{(infantMortality }\SpecialCharTok{\textasciitilde{}}\NormalTok{ ppgdp, tr, te,}
                  \AttributeTok{k =}\NormalTok{ kvals[j], }\AttributeTok{kernel =} \StringTok{"rectangular"}\NormalTok{)}
\NormalTok{      fold\_mse[f] }\OtherTok{\textless{}{-}} \FunctionTok{mean}\NormalTok{((te}\SpecialCharTok{$}\NormalTok{infantMortality }\SpecialCharTok{{-}}\NormalTok{ fit}\SpecialCharTok{$}\NormalTok{fitted.values)}\SpecialCharTok{\^{}}\DecValTok{2}\NormalTok{)}
\NormalTok{    \}}

\NormalTok{    out[j] }\OtherTok{\textless{}{-}} \FunctionTok{sqrt}\NormalTok{(}\FunctionTok{mean}\NormalTok{(fold\_mse))}
\NormalTok{  \}}

\NormalTok{  out}
\NormalTok{\}}
\end{Highlighting}
\end{Shaded}

You are not expected to be able to write that from scratch. You are
expected to be able to read it and say in words what it does, which is
the standard this course applies to code generally.

\begin{Shaded}
\begin{Highlighting}[]
\NormalTok{kseq2  }\OtherTok{\textless{}{-}} \DecValTok{2}\SpecialCharTok{:}\DecValTok{60}

\NormalTok{cv\_5  }\OtherTok{\textless{}{-}} \FunctionTok{cv\_knn}\NormalTok{(countries, kseq2, }\AttributeTok{nfold =} \DecValTok{5}\NormalTok{)}
\NormalTok{cv\_10 }\OtherTok{\textless{}{-}} \FunctionTok{cv\_knn}\NormalTok{(countries, kseq2, }\AttributeTok{nfold =} \DecValTok{10}\NormalTok{)}

\FunctionTok{plot}\NormalTok{(kseq2, cv\_10, }\AttributeTok{type =} \StringTok{"l"}\NormalTok{, }\AttributeTok{col =} \StringTok{"maroon"}\NormalTok{, }\AttributeTok{lwd =} \FloatTok{2.5}\NormalTok{,}
     \AttributeTok{ylim =} \FunctionTok{range}\NormalTok{(}\FunctionTok{c}\NormalTok{(cv\_5, cv\_10)),}
     \AttributeTok{xlab =} \StringTok{"k"}\NormalTok{, }\AttributeTok{ylab =} \StringTok{"Cross{-}validated RMSE"}\NormalTok{,}
     \AttributeTok{main =} \StringTok{"Choosing k by cross{-}validation"}\NormalTok{)}

\FunctionTok{lines}\NormalTok{(kseq2, cv\_5, }\AttributeTok{col =} \StringTok{"navyblue"}\NormalTok{, }\AttributeTok{lwd =} \FloatTok{2.5}\NormalTok{)}

\FunctionTok{abline}\NormalTok{(}\AttributeTok{v =}\NormalTok{ kseq2[}\FunctionTok{which.min}\NormalTok{(cv\_10)], }\AttributeTok{lty =} \DecValTok{2}\NormalTok{, }\AttributeTok{col =} \StringTok{"maroon"}\NormalTok{)}

\FunctionTok{legend}\NormalTok{(}\StringTok{"topleft"}\NormalTok{, }\AttributeTok{legend =} \FunctionTok{c}\NormalTok{(}\StringTok{"10{-}fold"}\NormalTok{, }\StringTok{"5{-}fold"}\NormalTok{),}
       \AttributeTok{col =} \FunctionTok{c}\NormalTok{(}\StringTok{"maroon"}\NormalTok{, }\StringTok{"navyblue"}\NormalTok{), }\AttributeTok{lwd =} \FloatTok{2.5}\NormalTok{, }\AttributeTok{box.lty =} \DecValTok{0}\NormalTok{, }\AttributeTok{cex =} \FloatTok{0.8}\NormalTok{)}
\end{Highlighting}
\end{Shaded}

\pandocbounded{\includegraphics[keepaspectratio]{handouts/w02-knn_files/figure-pdf/unnamed-chunk-13-1.pdf}}

\begin{Shaded}
\begin{Highlighting}[]
\NormalTok{kseq2[}\FunctionTok{which.min}\NormalTok{(cv\_5)]}
\end{Highlighting}
\end{Shaded}

\begin{verbatim}
[1] 24
\end{verbatim}

\begin{Shaded}
\begin{Highlighting}[]
\NormalTok{kseq2[}\FunctionTok{which.min}\NormalTok{(cv\_10)]}
\end{Highlighting}
\end{Shaded}

\begin{verbatim}
[1] 27
\end{verbatim}

Compare this curve with the single-split version. It is smoother,
because each point averages several held-out sets rather than resting on
one, and the two fold counts broadly agree with each other. That
agreement is the reassurance we were missing before.

It does not make the choice of \(k\) objective. Cross-validation still
leaves you with a range of values that perform about equally well, and
picking the exact minimum of a flat curve is reading precision into
noise. Report the range, say what you chose, and say why.

\section{The same thing, with caret}\label{the-same-thing-with-caret}

Nobody writes that loop in practice. Now that you know what
cross-validation does, here is the tool you will actually use for it,
and use for the rest of this course.

\texttt{caret} is the standard R package for fitting and comparing
predictive models through a single interface. It is worth installing
now, because Week 6 uses it again for evaluation.

\begin{Shaded}
\begin{Highlighting}[]
\FunctionTok{install.packages}\NormalTok{(}\StringTok{"caret"}\NormalTok{)}
\end{Highlighting}
\end{Shaded}

\begin{Shaded}
\begin{Highlighting}[]
\FunctionTok{library}\NormalTok{(caret)}

\FunctionTok{set.seed}\NormalTok{(}\DecValTok{7}\NormalTok{)}

\NormalTok{cv\_caret }\OtherTok{\textless{}{-}} \FunctionTok{train}\NormalTok{(}
\NormalTok{  infantMortality }\SpecialCharTok{\textasciitilde{}}\NormalTok{ ppgdp,}
  \AttributeTok{data       =}\NormalTok{ countries,}
  \AttributeTok{method     =} \StringTok{"knn"}\NormalTok{,}
  \AttributeTok{preProcess =} \FunctionTok{c}\NormalTok{(}\StringTok{"range"}\NormalTok{),}
  \AttributeTok{trControl  =} \FunctionTok{trainControl}\NormalTok{(}\AttributeTok{method =} \StringTok{"cv"}\NormalTok{, }\AttributeTok{number =} \DecValTok{10}\NormalTok{),}
  \AttributeTok{tuneGrid   =} \FunctionTok{data.frame}\NormalTok{(}\AttributeTok{k =}\NormalTok{ kseq2)}
\NormalTok{)}

\NormalTok{cv\_caret}\SpecialCharTok{$}\NormalTok{bestTune}
\end{Highlighting}
\end{Shaded}

\begin{verbatim}
    k
26 27
\end{verbatim}

\begin{Shaded}
\begin{Highlighting}[]
\FunctionTok{plot}\NormalTok{(cv\_caret)}
\end{Highlighting}
\end{Shaded}

\pandocbounded{\includegraphics[keepaspectratio]{handouts/w02-knn_files/figure-pdf/unnamed-chunk-15-1.pdf}}

Four arguments carry the work, and each corresponds to something we did
by hand.

\texttt{method\ =\ "knn"} picks the algorithm. \texttt{caret} supports
well over a hundred, and swapping this one string is how you will fit a
random forest next week.

\texttt{trControl\ =\ trainControl(method\ =\ "cv",\ number\ =\ 10)} is
the cross-validation. \texttt{number} is the fold count, our
\texttt{nfold}. Set \texttt{method\ =\ "LOOCV"} here and you get
leave-one-out instead.

\texttt{tuneGrid\ =\ data.frame(k\ =\ kseq2)} gives it the values of
\(k\) to try, the same sequence we looped over. Leave this out and
\texttt{caret} quietly tries three values of its own choosing, which is
almost never what you want.

\texttt{preProcess\ =\ c("range")} rescales the features to lie between
zero and one, which is the first of the two formulas from the previous
section. \textbf{This is not the default.} Omit it and \texttt{caret}
uses the raw units.

One honest caveat about that last point here. With a single predictor,
rescaling makes no difference at all: stretching one axis cannot change
which countries are nearest to which. It is in the argument above
because it becomes essential the moment we add a second feature, and
because a habit formed on the easy case is one you still have on the
hard one.

The answer should land close to what the loop produced, though not
necessarily identical, since \texttt{caret} draws its own folds. If the
two disagree by more than a few units of \(k\) on a curve this flat,
that is the flatness talking rather than a real disagreement.

Week 6 returns to evaluation and takes it further, into the metrics that
matter when the outcome is rare and accuracy stops being informative.

\chapter{More than one feature, and the problem of
units}\label{more-than-one-feature-and-the-problem-of-units}

So far similarity has meant closeness on a single axis, which is easy.
With several features we need a general definition of distance, and the
usual one is Euclidean distance. For a case \(f\) with features
\((x_{f1}, \dots, x_{fp})\) and a case \(i\) in the training data,

\[d_{if} = \sqrt{\sum_{j=1}^{p}\left(x_{fj} - x_{ij}\right)^2}.\]

This is the straight-line distance you already know, extended to more
than two dimensions. And it introduces a problem that is easy to miss
and consequential when you do.

Our features are measured on very different scales. Per capita GDP runs
from well under one up to about a hundred, in thousands of dollars.
Percent urban runs from zero to a hundred. Fertility runs from about one
to about seven, and female life expectancy from the high forties to the
high eighties.

\begin{Shaded}
\begin{Highlighting}[]
\FunctionTok{summary}\NormalTok{(countries[, }\FunctionTok{c}\NormalTok{(}\StringTok{"ppgdp"}\NormalTok{, }\StringTok{"pctUrban"}\NormalTok{, }\StringTok{"fertility"}\NormalTok{, }\StringTok{"lifeExpF"}\NormalTok{)])}
\end{Highlighting}
\end{Shaded}

\begin{verbatim}
     ppgdp             pctUrban       fertility        lifeExpF    
 Min.   :  0.1148   Min.   : 11.0   Min.   :1.134   Min.   :48.11  
 1st Qu.:  1.2398   1st Qu.: 39.0   1st Qu.:1.750   1st Qu.:65.10  
 Median :  4.4958   Median : 59.0   Median :2.264   Median :75.57  
 Mean   : 12.2911   Mean   : 57.1   Mean   :2.780   Mean   :72.08  
 3rd Qu.: 14.4973   3rd Qu.: 75.0   3rd Qu.:3.700   3rd Qu.:79.07  
 Max.   :105.0954   Max.   :100.0   Max.   :6.925   Max.   :87.12  
\end{verbatim}

Put those into the distance formula unchanged and the two wide-ranging
features decide everything. A difference of fifty between two countries
on per capita GDP contributes two and a half thousand to the sum under
the square root. The entire possible range of fertility, from one child
per woman to seven, contributes about thirty-five. Fertility is not
being weighted lightly here. For practical purposes it is not in the
model at all, and neither is anything else measured on a small scale.

The standard fix is to put all features on a common scale first, either

\[\tilde{x} = \frac{x - \min(x)}{\max(x) - \min(x)} \qquad \text{or} \qquad \tilde{x} = \frac{x - \bar{x}}{\sigma_x}.\]

\begin{Shaded}
\begin{Highlighting}[]
\NormalTok{rescale }\OtherTok{\textless{}{-}} \ControlFlowTok{function}\NormalTok{(x) (x }\SpecialCharTok{{-}} \FunctionTok{min}\NormalTok{(x)) }\SpecialCharTok{/}\NormalTok{ (}\FunctionTok{max}\NormalTok{(x) }\SpecialCharTok{{-}} \FunctionTok{min}\NormalTok{(x))}

\NormalTok{feats }\OtherTok{\textless{}{-}} \FunctionTok{c}\NormalTok{(}\StringTok{"ppgdp"}\NormalTok{, }\StringTok{"pctUrban"}\NormalTok{, }\StringTok{"fertility"}\NormalTok{, }\StringTok{"lifeExpF"}\NormalTok{)}

\NormalTok{countries\_s }\OtherTok{\textless{}{-}}\NormalTok{ countries}
\NormalTok{countries\_s[feats] }\OtherTok{\textless{}{-}} \FunctionTok{lapply}\NormalTok{(countries\_s[feats], rescale)}

\FunctionTok{summary}\NormalTok{(countries\_s[, feats])}
\end{Highlighting}
\end{Shaded}

\begin{verbatim}
     ppgdp            pctUrban        fertility         lifeExpF     
 Min.   :0.00000   Min.   :0.0000   Min.   :0.0000   Min.   :0.0000  
 1st Qu.:0.01072   1st Qu.:0.3146   1st Qu.:0.1064   1st Qu.:0.4355  
 Median :0.04173   Median :0.5393   Median :0.1951   Median :0.7039  
 Mean   :0.11599   Mean   :0.5180   Mean   :0.2843   Mean   :0.6144  
 3rd Qu.:0.13700   3rd Qu.:0.7191   3rd Qu.:0.4431   3rd Qu.:0.7936  
 Max.   :1.00000   Max.   :1.0000   Max.   :1.0000   Max.   :1.0000  
\end{verbatim}

The first thing to establish is that this is a harmless operation.
Rescaling does not change the relationship between any feature and the
outcome. Look at two features before and after.

\begin{Shaded}
\begin{Highlighting}[]
\FunctionTok{par}\NormalTok{(}\AttributeTok{mfrow =} \FunctionTok{c}\NormalTok{(}\DecValTok{2}\NormalTok{, }\DecValTok{2}\NormalTok{))}

\FunctionTok{plot}\NormalTok{(countries}\SpecialCharTok{$}\NormalTok{ppgdp, countries}\SpecialCharTok{$}\NormalTok{infantMortality, }\AttributeTok{col =} \StringTok{"maroon"}\NormalTok{, }\AttributeTok{cex =} \FloatTok{0.7}\NormalTok{,}
     \AttributeTok{xlab =} \StringTok{"Per capita GDP (thousands of US$)"}\NormalTok{, }\AttributeTok{ylab =} \StringTok{"Infant mortality"}\NormalTok{,}
     \AttributeTok{main =} \StringTok{"Raw units"}\NormalTok{, }\AttributeTok{font.main =} \DecValTok{1}\NormalTok{)}
\FunctionTok{plot}\NormalTok{(countries\_s}\SpecialCharTok{$}\NormalTok{ppgdp, countries\_s}\SpecialCharTok{$}\NormalTok{infantMortality, }\AttributeTok{col =} \StringTok{"maroon"}\NormalTok{, }\AttributeTok{cex =} \FloatTok{0.7}\NormalTok{,}
     \AttributeTok{xlab =} \StringTok{"Per capita GDP (rescaled 0 to 1)"}\NormalTok{, }\AttributeTok{ylab =} \StringTok{"Infant mortality"}\NormalTok{,}
     \AttributeTok{main =} \StringTok{"Rescaled"}\NormalTok{, }\AttributeTok{font.main =} \DecValTok{1}\NormalTok{)}

\FunctionTok{plot}\NormalTok{(countries}\SpecialCharTok{$}\NormalTok{pctUrban, countries}\SpecialCharTok{$}\NormalTok{infantMortality, }\AttributeTok{col =} \StringTok{"maroon"}\NormalTok{, }\AttributeTok{cex =} \FloatTok{0.7}\NormalTok{,}
     \AttributeTok{xlab =} \StringTok{"Percent urban"}\NormalTok{, }\AttributeTok{ylab =} \StringTok{"Infant mortality"}\NormalTok{,}
     \AttributeTok{main =} \StringTok{"Raw units"}\NormalTok{, }\AttributeTok{font.main =} \DecValTok{1}\NormalTok{)}

\FunctionTok{plot}\NormalTok{(countries\_s}\SpecialCharTok{$}\NormalTok{pctUrban, countries\_s}\SpecialCharTok{$}\NormalTok{infantMortality, }\AttributeTok{col =} \StringTok{"maroon"}\NormalTok{, }\AttributeTok{cex =} \FloatTok{0.7}\NormalTok{,}
     \AttributeTok{xlab =} \StringTok{"Percent urban (rescaled 0 to 1)"}\NormalTok{, }\AttributeTok{ylab =} \StringTok{"Infant mortality"}\NormalTok{,}
     \AttributeTok{main =} \StringTok{"Rescaled"}\NormalTok{, }\AttributeTok{font.main =} \DecValTok{1}\NormalTok{)}
\end{Highlighting}
\end{Shaded}

\pandocbounded{\includegraphics[keepaspectratio]{handouts/w02-knn_files/figure-pdf/unnamed-chunk-18-1.pdf}}

The clouds of points are identical in shape. Only the numbers on the
horizontal axis have moved. Nothing about the association between any
feature and the outcome has been altered.

What has changed is invisible in these plots, because it is not about
any single feature. It is about the space the features jointly define.

\section{What the algorithm now sees}\label{what-the-algorithm-now-sees}

With two features rather than one, cases no longer sit along a line.
They sit in a plane, and the outcome rises above it.

\begin{Shaded}
\begin{Highlighting}[]
\FunctionTok{library}\NormalTok{(scatterplot3d)}

\FunctionTok{scatterplot3d}\NormalTok{(countries\_s}\SpecialCharTok{$}\NormalTok{ppgdp, countries\_s}\SpecialCharTok{$}\NormalTok{pctUrban, countries\_s}\SpecialCharTok{$}\NormalTok{infantMortality,}
              \AttributeTok{color =} \StringTok{"maroon"}\NormalTok{, }\AttributeTok{pch =} \DecValTok{1}\NormalTok{, }\AttributeTok{cex.symbols =} \FloatTok{0.7}\NormalTok{,}
              \AttributeTok{xlab =} \StringTok{"Per capita GDP (rescaled)"}\NormalTok{,}
              \AttributeTok{ylab =} \StringTok{"Percent urban (rescaled)"}\NormalTok{,}
              \AttributeTok{zlab =} \StringTok{"Infant mortality"}\NormalTok{)}
\end{Highlighting}
\end{Shaded}

\pandocbounded{\includegraphics[keepaspectratio]{handouts/w02-knn_files/figure-pdf/unnamed-chunk-19-1.pdf}}

Add a third feature and the picture cannot be drawn at all, which is the
last point in this course at which you can see the space a model is
working in. From Week 3 onward you will have to reason about it rather
than look at it. The arithmetic does not care: the distance formula
extends to any number of features without complaint, and it is only our
eyes that stop at three.

\section{What the neighborhood looks
like}\label{what-the-neighborhood-looks-like}

Now the thing worth seeing. Pick a country, and ask which twenty-five
countries the algorithm considers its neighbors. We do it twice, once in
raw units and once rescaled, using the same country both times.

\begin{Shaded}
\begin{Highlighting}[]
\NormalTok{target\_raw }\OtherTok{\textless{}{-}} \FunctionTok{c}\NormalTok{(}\FunctionTok{median}\NormalTok{(countries}\SpecialCharTok{$}\NormalTok{ppgdp),   }\FunctionTok{median}\NormalTok{(countries}\SpecialCharTok{$}\NormalTok{fertility))}
\NormalTok{target\_s   }\OtherTok{\textless{}{-}} \FunctionTok{c}\NormalTok{(}\FunctionTok{median}\NormalTok{(countries\_s}\SpecialCharTok{$}\NormalTok{ppgdp), }\FunctionTok{median}\NormalTok{(countries\_s}\SpecialCharTok{$}\NormalTok{fertility))}

\DocumentationTok{\#\# Euclidean distance, exactly the formula above, computed by hand}
\NormalTok{dist\_raw }\OtherTok{\textless{}{-}} \FunctionTok{sqrt}\NormalTok{((countries}\SpecialCharTok{$}\NormalTok{ppgdp     }\SpecialCharTok{{-}}\NormalTok{ target\_raw[}\DecValTok{1}\NormalTok{])}\SpecialCharTok{\^{}}\DecValTok{2} \SpecialCharTok{+}
\NormalTok{                 (countries}\SpecialCharTok{$}\NormalTok{fertility }\SpecialCharTok{{-}}\NormalTok{ target\_raw[}\DecValTok{2}\NormalTok{])}\SpecialCharTok{\^{}}\DecValTok{2}\NormalTok{)}

\NormalTok{dist\_s   }\OtherTok{\textless{}{-}} \FunctionTok{sqrt}\NormalTok{((countries\_s}\SpecialCharTok{$}\NormalTok{ppgdp     }\SpecialCharTok{{-}}\NormalTok{ target\_s[}\DecValTok{1}\NormalTok{])}\SpecialCharTok{\^{}}\DecValTok{2} \SpecialCharTok{+}
\NormalTok{                 (countries\_s}\SpecialCharTok{$}\NormalTok{fertility }\SpecialCharTok{{-}}\NormalTok{ target\_s[}\DecValTok{2}\NormalTok{])}\SpecialCharTok{\^{}}\DecValTok{2}\NormalTok{)}

\NormalTok{nb\_raw }\OtherTok{\textless{}{-}} \FunctionTok{order}\NormalTok{(dist\_raw)[}\DecValTok{1}\SpecialCharTok{:}\DecValTok{25}\NormalTok{]}
\NormalTok{nb\_s   }\OtherTok{\textless{}{-}} \FunctionTok{order}\NormalTok{(dist\_s)[}\DecValTok{1}\SpecialCharTok{:}\DecValTok{25}\NormalTok{]}

\FunctionTok{par}\NormalTok{(}\AttributeTok{mfrow =} \FunctionTok{c}\NormalTok{(}\DecValTok{1}\NormalTok{, }\DecValTok{2}\NormalTok{))}

\FunctionTok{plot}\NormalTok{(countries}\SpecialCharTok{$}\NormalTok{ppgdp, countries}\SpecialCharTok{$}\NormalTok{fertility, }\AttributeTok{col =} \StringTok{"gray55"}\NormalTok{, }\AttributeTok{pch =} \DecValTok{19}\NormalTok{, }\AttributeTok{cex =} \FloatTok{0.7}\NormalTok{,}
     \AttributeTok{xlab =} \StringTok{"Per capita GDP (thousands of US$)"}\NormalTok{, }\AttributeTok{ylab =} \StringTok{"Fertility (births per woman)"}\NormalTok{,}
     \AttributeTok{main =} \StringTok{"Raw units, k = 25"}\NormalTok{, }\AttributeTok{font.main =} \DecValTok{1}\NormalTok{)}
\FunctionTok{points}\NormalTok{(countries}\SpecialCharTok{$}\NormalTok{ppgdp[nb\_raw], countries}\SpecialCharTok{$}\NormalTok{fertility[nb\_raw], }\AttributeTok{col =} \StringTok{"red"}\NormalTok{, }\AttributeTok{cex =} \FloatTok{1.6}\NormalTok{)}
\FunctionTok{points}\NormalTok{(target\_raw[}\DecValTok{1}\NormalTok{], target\_raw[}\DecValTok{2}\NormalTok{], }\AttributeTok{col =} \StringTok{"red"}\NormalTok{, }\AttributeTok{pch =} \DecValTok{19}\NormalTok{, }\AttributeTok{cex =} \FloatTok{1.7}\NormalTok{)}
\FunctionTok{abline}\NormalTok{(}\AttributeTok{v =}\NormalTok{ target\_raw[}\DecValTok{1}\NormalTok{], }\AttributeTok{h =}\NormalTok{ target\_raw[}\DecValTok{2}\NormalTok{], }\AttributeTok{col =} \StringTok{"darkgreen"}\NormalTok{, }\AttributeTok{lty =} \DecValTok{2}\NormalTok{)}

\FunctionTok{plot}\NormalTok{(countries\_s}\SpecialCharTok{$}\NormalTok{ppgdp, countries\_s}\SpecialCharTok{$}\NormalTok{fertility, }\AttributeTok{col =} \StringTok{"gray55"}\NormalTok{, }\AttributeTok{pch =} \DecValTok{19}\NormalTok{, }\AttributeTok{cex =} \FloatTok{0.7}\NormalTok{,}
     \AttributeTok{xlab =} \StringTok{"Per capita GDP (rescaled)"}\NormalTok{, }\AttributeTok{ylab =} \StringTok{"Fertility (rescaled)"}\NormalTok{,}
     \AttributeTok{main =} \StringTok{"Rescaled, k = 25"}\NormalTok{, }\AttributeTok{font.main =} \DecValTok{1}\NormalTok{)}
\FunctionTok{points}\NormalTok{(countries\_s}\SpecialCharTok{$}\NormalTok{ppgdp[nb\_s], countries\_s}\SpecialCharTok{$}\NormalTok{fertility[nb\_s], }\AttributeTok{col =} \StringTok{"red"}\NormalTok{, }\AttributeTok{cex =} \FloatTok{1.6}\NormalTok{)}
\FunctionTok{points}\NormalTok{(target\_s[}\DecValTok{1}\NormalTok{], target\_s[}\DecValTok{2}\NormalTok{], }\AttributeTok{col =} \StringTok{"red"}\NormalTok{, }\AttributeTok{pch =} \DecValTok{19}\NormalTok{, }\AttributeTok{cex =} \FloatTok{1.7}\NormalTok{)}
\FunctionTok{abline}\NormalTok{(}\AttributeTok{v =}\NormalTok{ target\_s[}\DecValTok{1}\NormalTok{], }\AttributeTok{h =}\NormalTok{ target\_s[}\DecValTok{2}\NormalTok{], }\AttributeTok{col =} \StringTok{"darkgreen"}\NormalTok{, }\AttributeTok{lty =} \DecValTok{2}\NormalTok{)}
\end{Highlighting}
\end{Shaded}

\pandocbounded{\includegraphics[keepaspectratio]{handouts/w02-knn_files/figure-pdf/unnamed-chunk-20-1.pdf}}

The filled red point is the country we are predicting for. The circled
points are the twenty-five the algorithm chose as most similar to it.

On the right, with both features on a common scale, the neighborhood is
roughly a disc around the target. Both features contributed to picking
it, which is what we asked for.

On the left it is a narrow vertical band. The selected countries share
almost exactly the target's income and are spread across nearly the
whole range of fertility, from around one child per woman to well above
the average. Per capita GDP spans about a hundred units on this scale
and fertility spans about six, so a difference of a few units of income
outweighs the entire fertility axis. The algorithm did not ignore
fertility because anybody judged it unimportant. It ignored it because
of the units it happens to be measured in.

Note also what the model would have reported. Both versions return
predictions. Neither prints a warning. The left-hand model would have
gone into a paper as a two-feature model, and the fact that it was a
one-feature model in disguise would appear nowhere in the output.

Stay with that, because it is not a technical footnote. The algorithm's
entire operation rests on its definition of ``similar,'' and that
definition was just set by two choices: which features we included, and
how we scaled them. Neither choice required anyone to make a claim about
the world. Had we left the units alone, the model would have decided
that countries are similar when their incomes are similar, and it would
have reported that decision as a finding rather than as an artifact of
the fact that GDP happens to be counted in dollars.

\chapter{What this model cannot tell
you}\label{what-this-model-cannot-tell-you}

The model has no coefficients. There is no number to report saying how
much per capita GDP matters, no standard error, no statement of the form
``a one-unit increase in \(X\) is associated with a change in \(Y\).''
The model cannot be written down apart from the data it holds.

More importantly, it has said nothing about \emph{why} infant mortality
is lower in richer countries. It has found that countries with similar
incomes tend to have similar infant mortality, and it can use that to
predict. That is a prediction, in exactly the sense we distinguished
last week, and it does not become an explanation by performing well.

There is a second limit, and the neighborhood figures above showed it. A
model of this kind answers every question put to it, at the same
apparent level of confidence, whether the case sits in a crowded part of
the data or alone at the edge of it. The richest country in this dataset
gets a prediction assembled from neighbors spread across tens of
thousands of dollars of income. That prediction is printed in the same
format, to the same number of decimal places, as one built from
twenty-five genuinely similar cases. Nothing in the output distinguishes
them.

\chapter{The politics of similarity}\label{the-politics-of-similarity}

The rule at the center of this week is that cases similar to past cases
receive what those past cases received.

That is also a fair description of precedent in law, of case-based
reasoning in administration, and of the Weberian bureaucratic ideal we
discussed last week, in which similar cases are treated alike regardless
of who happens to be at the desk. Consistency of this kind is not a
defect. It is one of the oldest arguments for rule-bound government, and
part of what makes an administration predictable enough to be held to
account.

What changes when the rule is a fitted model is who defines similarity,
and how visibly. In \(K\)-nearest neighbors that definition came from
feature selection and rescaling, decisions made early, made quietly, and
almost never revisited or documented once the system is running. Nobody
writes them down as a policy, because they do not look like a policy.

And a model of this kind can only reproduce the pattern in the cases it
was given. If the treatment recorded in those cases was unjust, then
similarity to past cases carries the injustice forward, and it does so
with the appearance of perfect consistency. The system is not
malfunctioning when this happens. It is doing precisely what it was
built to do.

\chapter{Reading}\label{reading-1}

James, G., Witten, D., Hastie, T., and Tibshirani, R. \emph{An
Introduction to Statistical Learning}. Sections 2.1 and 2.2, Section
3.5, and Section 5.1. Freely available at
\href{https://www.statlearning.com/}{statlearning.com} (James et al.
2021).

Read Section 2.2 for the figures rather than the derivations. The
pictures of training error against test error, and of the bias-variance
decomposition, are the same pictures you produced above, and seeing them
drawn by someone else is a useful check on whether you understood what
you made. Section 5.1 covers cross-validation and is short.

\chapter{Before the workshop}\label{before-the-workshop}

Install \texttt{carData}, \texttt{kknn}, \texttt{caret}, and
\texttt{scatterplot3d}, run every chunk in this handout once, and bring
a charged laptop. The workshop begins from the assumption that the code
above ran on your machine, and the first part of a three-hour session is
not the place to discover an installation problem.

If something failed, say so at the start rather than at the end.

\chapter{Questions for the session}\label{questions-for-the-session-1}

Come with a view on at least two of these. They will structure the first
part of the session.

\begin{enumerate}
\def\labelenumi{\arabic{enumi}.}
\item
  \(K\)-nearest neighbors has no coefficients and makes no claim about
  mechanism. Is it therefore useless for political science, or is that
  the wrong test to apply to it? Last week's reading is relevant here.
\item
  The choice of \(k\) has no correct answer that can be read off the
  data alone. What does it mean that a central decision in a model is
  made by the researcher rather than discovered?
\item
  We rescaled the features so that income would not dominate the
  distance calculation. Suppose someone had not. Would you describe the
  resulting model as wrong, or as encoding a particular view of what
  makes two countries similar? Does the distinction matter for how such
  a system should be reviewed?
\item
  In the neighborhood figures, the left-hand model was given two
  features and used essentially one. It returned predictions and printed
  no warning. What follows for how a system you did not build should be
  reviewed before anyone relies on its outputs?
\item
  The model returns a prediction for the richest country in the data,
  assembled from neighbors that are nothing like it, and reports it
  exactly as it reports any other. Where in political or administrative
  use would that distinction matter most, and what would a responsible
  system do instead?
\item
  Bureaucratic rule-following and nearest-neighbor prediction both treat
  like cases alike. Identify one political difference between them that
  is not simply that the model is more complicated.
\end{enumerate}

\chapter{Trees and Random Forests}\label{trees-and-random-forests}

\emph{Week 3, 16 September}

Last week's method concealed nothing, in the sense that every prediction
\(K\)-nearest neighbors produced could be traced back to a small set of
countries that could be located on a scatterplot. This week begins from
the same premise, with a model that can be drawn on a whiteboard and
read aloud as a sequence of rules, and then abandons that transparency
deliberately, in exchange for a considerably better model.

The session therefore stages a tension that runs through much of
contemporary machine learning. On the one hand, a model whose reasoning
can be inspected is easier to defend, to audit, and to explain to
somebody who has to act on its output. On the other hand, the models
that predict best are frequently the ones whose reasoning cannot be
inspected at all. Trees and forests are unusually well suited to
examining this tension, because the movement from the first situation to
the second occurs within a single family of methods and can be observed
at each step.

\section{What trees are good at, and what they are
not}\label{what-trees-are-good-at-and-what-they-are-not}

Tree models in statistical learning bear an obvious kinship to decision
trees as we ordinarily understand them, and they have been used widely
for prediction. It is worth setting out at the start what they offer and
what they cost.

\textbf{Benefits.}

\begin{itemize}
\tightlist
\item
  They capture interactions between features and non-linear associations
  without being instructed to look for them.
\item
  They handle categorical variables well, and the results remain
  interpretable.
\item
  They are relatively fast compared with most other machine-learning
  algorithms.
\item
  \textbf{They do not require the features to be rescaled.}
\end{itemize}

The final point deserves a moment, since it undoes an entire section of
last week. \(K\)-nearest neighbors was at the mercy of the units in
which its features happened to be measured, and we devoted considerable
effort to placing everything on a common scale so that income would not
swamp fertility. A tree, however, never compares two features to each
other. Instead, it asks one question at a time, about a single variable,
and on that variable's own scale, so that multiplying a feature by a
thousand changes nothing whatever and the problem simply does not arise.

\textbf{The disadvantage,} and it is a serious one, is that a single
tree does not predict well out of sample.

That tension is the shape of the week, in that everything after the
first third is an attempt to retain the benefits while addressing the
disadvantage, by growing a forest of trees and allowing them to vote.

\section{Two datasets this week}\label{two-datasets-this-week}

Part I returns to the country data from last week, on the grounds that
the fastest way to see what is distinctive about a tree is to apply it
to material you have already modeled by other means.

Part II moves to Fearon and Laitin's civil war data (Fearon and Laitin
2003), where the outcome is binary and rare, and where the more
difficult problems arise.

\begin{Shaded}
\begin{Highlighting}[]
\FunctionTok{install.packages}\NormalTok{(}\FunctionTok{c}\NormalTok{(}\StringTok{"carData"}\NormalTok{, }\StringTok{"rpart"}\NormalTok{, }\StringTok{"rpart.plot"}\NormalTok{, }\StringTok{"tree"}\NormalTok{,}
                   \StringTok{"randomForest"}\NormalTok{, }\StringTok{"caret"}\NormalTok{, }\StringTok{"ggplot2"}\NormalTok{))}
\end{Highlighting}
\end{Shaded}

\begin{Shaded}
\begin{Highlighting}[]
\FunctionTok{library}\NormalTok{(carData)}
\FunctionTok{library}\NormalTok{(rpart)}
\FunctionTok{library}\NormalTok{(rpart.plot)}
\FunctionTok{library}\NormalTok{(tree)}
\FunctionTok{library}\NormalTok{(randomForest)}
\FunctionTok{library}\NormalTok{(caret)}
\FunctionTok{library}\NormalTok{(ggplot2)}

\FunctionTok{options}\NormalTok{(}\AttributeTok{scipen =} \DecValTok{999}\NormalTok{)}

\FunctionTok{data}\NormalTok{(UN, }\AttributeTok{package =} \StringTok{"carData"}\NormalTok{)}

\NormalTok{countries }\OtherTok{\textless{}{-}} \FunctionTok{na.omit}\NormalTok{(UN[, }\FunctionTok{c}\NormalTok{(}\StringTok{"ppgdp"}\NormalTok{, }\StringTok{"infantMortality"}\NormalTok{, }\StringTok{"fertility"}\NormalTok{,}
                            \StringTok{"pctUrban"}\NormalTok{, }\StringTok{"lifeExpF"}\NormalTok{, }\StringTok{"group"}\NormalTok{)])}

\NormalTok{countries}\SpecialCharTok{$}\NormalTok{ppgdp }\OtherTok{\textless{}{-}}\NormalTok{ countries}\SpecialCharTok{$}\NormalTok{ppgdp }\SpecialCharTok{/} \DecValTok{1000}   \CommentTok{\# thousands, as last week}

\FunctionTok{nrow}\NormalTok{(countries)}
\end{Highlighting}
\end{Shaded}

\begin{verbatim}
[1] 193
\end{verbatim}

One row per country, which is worth noting now, because Part II will not
have that luxury.

\section{Part I. A regression tree}\label{part-i.-a-regression-tree}

Last week we fitted a \(K\)-nearest neighbors curve to the relationship
between per capita GDP and infant mortality, and gave some attention to
the fact that the relationship bends: mortality falls very steeply
across the poorest range and then flattens almost completely. A straight
line was therefore a poor description of the relationship, and
\(K\)-nearest neighbors accommodated the curvature instead by averaging
over a moving window of neighbors.

A tree accommodates it differently, and placing the two side by side is
the most efficient way to establish what a tree is.

\subsection{One feature}\label{one-feature}

\begin{Shaded}
\begin{Highlighting}[]
\NormalTok{m\_tree }\OtherTok{\textless{}{-}} \FunctionTok{rpart}\NormalTok{(infantMortality }\SpecialCharTok{\textasciitilde{}}\NormalTok{ ppgdp, }\AttributeTok{data =}\NormalTok{ countries)}

\FunctionTok{rpart.plot}\NormalTok{(m\_tree, }\AttributeTok{box.palette =} \StringTok{"GnBu"}\NormalTok{)}
\end{Highlighting}
\end{Shaded}

\pandocbounded{\includegraphics[keepaspectratio]{handouts/w03-trees-forests_files/figure-pdf/unnamed-chunk-3-1.pdf}}

At each interior node there is a decision rule of the form \(x > c\).
Follow the branch that matches, and continue until you reach a bottom
node, which is called a leaf. The number in the leaf is the prediction,
namely the average infant mortality of all the countries that arrived
there.

Reading one path aloud amounts to stating a substantive claim about
development and child survival in ordinary language, without recourse to
a formula, and very few of the models encountered in this course can be
communicated in that way.

We can now examine what the tree is as a function.

\begin{Shaded}
\begin{Highlighting}[]
\NormalTok{oo  }\OtherTok{\textless{}{-}} \FunctionTok{order}\NormalTok{(countries}\SpecialCharTok{$}\NormalTok{ppgdp)}
\NormalTok{ddf }\OtherTok{\textless{}{-}}\NormalTok{ countries[oo, ]}

\FunctionTok{plot}\NormalTok{(ddf}\SpecialCharTok{$}\NormalTok{ppgdp, ddf}\SpecialCharTok{$}\NormalTok{infantMortality,}
     \AttributeTok{col =} \StringTok{"gray55"}\NormalTok{, }\AttributeTok{pch =} \DecValTok{19}\NormalTok{, }\AttributeTok{cex =} \FloatTok{0.7}\NormalTok{,}
     \AttributeTok{xlab =} \StringTok{"Per capita GDP (thousands of US$)"}\NormalTok{,}
     \AttributeTok{ylab =} \StringTok{"Infant mortality (per 1,000 live births)"}\NormalTok{)}

\FunctionTok{lines}\NormalTok{(ddf}\SpecialCharTok{$}\NormalTok{ppgdp, }\FunctionTok{predict}\NormalTok{(m\_tree, ddf), }\AttributeTok{col =} \StringTok{"maroon"}\NormalTok{, }\AttributeTok{lwd =} \DecValTok{3}\NormalTok{)}

\ControlFlowTok{for}\NormalTok{ (i }\ControlFlowTok{in}\NormalTok{ m\_tree}\SpecialCharTok{$}\NormalTok{splits[, }\DecValTok{4}\NormalTok{]) \{}
  \FunctionTok{abline}\NormalTok{(}\AttributeTok{v =}\NormalTok{ i, }\AttributeTok{col =} \StringTok{"orange"}\NormalTok{, }\AttributeTok{lty =} \DecValTok{2}\NormalTok{)}
\NormalTok{\}}
\end{Highlighting}
\end{Shaded}

\pandocbounded{\includegraphics[keepaspectratio]{handouts/w03-trees-forests_files/figure-pdf/unnamed-chunk-4-1.pdf}}

The set of leaves provides a partition of the predictor space into
disjoint regions. With a single feature that partition is simply a set
of intervals, marked by the dashed vertical lines. Within each interval
we take the average outcome of the countries that fall there, which
produces the step function drawn in red. That step function is
\(\hat{f}\), and to predict for a new country we locate its interval and
read off the corresponding step.

The placement of the boundaries is itself informative, in that they are
tightly spaced at the left, where mortality changes rapidly, and widely
spaced at the right, where it does not. Notably, nothing in the
procedure instructed the algorithm to distribute them in this way.

\subsection{The same data, two
algorithms}\label{the-same-data-two-algorithms}

Last week's model can usefully be placed on the same figure, so that the
two procedures may be compared directly.

\begin{Shaded}
\begin{Highlighting}[]
\FunctionTok{library}\NormalTok{(kknn)}

\NormalTok{grid }\OtherTok{\textless{}{-}} \FunctionTok{data.frame}\NormalTok{(}\AttributeTok{ppgdp =} \FunctionTok{sort}\NormalTok{(countries}\SpecialCharTok{$}\NormalTok{ppgdp))}

\NormalTok{knn\_fit }\OtherTok{\textless{}{-}} \FunctionTok{kknn}\NormalTok{(infantMortality }\SpecialCharTok{\textasciitilde{}}\NormalTok{ ppgdp, }\AttributeTok{train =}\NormalTok{ countries, }\AttributeTok{test =}\NormalTok{ grid,}
                \AttributeTok{k =} \DecValTok{10}\NormalTok{, }\AttributeTok{kernel =} \StringTok{"rectangular"}\NormalTok{)}

\FunctionTok{plot}\NormalTok{(ddf}\SpecialCharTok{$}\NormalTok{ppgdp, ddf}\SpecialCharTok{$}\NormalTok{infantMortality,}
     \AttributeTok{col =} \StringTok{"gray75"}\NormalTok{, }\AttributeTok{pch =} \DecValTok{19}\NormalTok{, }\AttributeTok{cex =} \FloatTok{0.7}\NormalTok{,}
     \AttributeTok{xlab =} \StringTok{"Per capita GDP (thousands of US$)"}\NormalTok{,}
     \AttributeTok{ylab =} \StringTok{"Infant mortality (per 1,000 live births)"}\NormalTok{,}
     \AttributeTok{main =} \StringTok{"Tree and KNN on the same data"}\NormalTok{)}

\FunctionTok{lines}\NormalTok{(grid}\SpecialCharTok{$}\NormalTok{ppgdp, knn\_fit}\SpecialCharTok{$}\NormalTok{fitted.values, }\AttributeTok{col =} \StringTok{"navy"}\NormalTok{, }\AttributeTok{lwd =} \FloatTok{2.5}\NormalTok{)}
\FunctionTok{lines}\NormalTok{(ddf}\SpecialCharTok{$}\NormalTok{ppgdp, }\FunctionTok{predict}\NormalTok{(m\_tree, ddf),  }\AttributeTok{col =} \StringTok{"maroon"}\NormalTok{, }\AttributeTok{lwd =} \DecValTok{3}\NormalTok{)}

\FunctionTok{legend}\NormalTok{(}\StringTok{"topright"}\NormalTok{, }\AttributeTok{legend =} \FunctionTok{c}\NormalTok{(}\StringTok{"Regression tree"}\NormalTok{, }\StringTok{"KNN, k = 10"}\NormalTok{),}
       \AttributeTok{col =} \FunctionTok{c}\NormalTok{(}\StringTok{"maroon"}\NormalTok{, }\StringTok{"navy"}\NormalTok{), }\AttributeTok{lwd =} \FunctionTok{c}\NormalTok{(}\DecValTok{3}\NormalTok{, }\FloatTok{2.5}\NormalTok{), }\AttributeTok{box.lty =} \DecValTok{0}\NormalTok{, }\AttributeTok{cex =} \FloatTok{0.85}\NormalTok{)}
\end{Highlighting}
\end{Shaded}

\pandocbounded{\includegraphics[keepaspectratio]{handouts/w03-trees-forests_files/figure-pdf/unnamed-chunk-5-1.pdf}}

Both are step functions constructed from local averages, and neither was
told the shape of the relationship in advance. The difference between
them lies entirely in the manner in which the neighborhoods are
determined.

\(K\)-nearest neighbors uses a window of fixed size that slides along
the axis. Every prediction averages exactly ten countries, whether those
ten lie within a few hundred dollars of one another or are distributed
across forty thousand. This is why the curve becomes unstable on the
right, where the data thin out: the window preserves its count while
losing its substantive meaning.

The tree, by contrast, uses regions of fixed content, having selected
where to cut and therefore how wide each region should be, narrowly
where the outcome changes quickly and broadly where it does not. The
cost of proceeding this way is that the steps are coarse. The gain,
however, is that the width of a neighborhood becomes a decision the
algorithm reached about the data, rather than a number imposed upon it
in advance.

\subsection{What the algorithm is
minimizing}\label{what-the-algorithm-is-minimizing}

We are splitting the feature space \(X = \{x_1, x_2, \dots, x_k\}\) into
smaller rectangles, or boxes. If we can identify a pattern of splits
with low prediction error on the training data, we can then predict for
a new case by determining which box it belongs to.

Roughly, a tree may be thought of as a \(K\)-nearest neighbors algorithm
in which \(k\) is dynamic: the neighborhood is a box rather than a
window, and the boxes are wide where the data are sparse and narrow
where they are dense, rather than fixed in advance. The figure above is
that proposition drawn, and it is worth returning to whenever the
relationship between the two methods becomes unclear.

Formally, for a continuous outcome:

\begin{align}
\min_{R_1, R_2, \dots, R_J} \sum_{j=1}^{J} \sum_{i \in R_j} (y_i - \hat{y}_{R_j})^2
\end{align}

The boxes \(R_j\) are chosen so as to minimize the total squared
deviation of observations from the mean response within each box. In
principle there are infinitely many ways of defining the boxes, and
evaluating all of them is computationally impossible. A
\textbf{recursive partitioning} procedure is therefore used instead:

\begin{itemize}
\tightlist
\item
  Scan every feature to identify the one whose split yields the best
  improvement, and split there.
\item
  Repeat within each resulting box, and continue until splitting becomes
  too costly or a pre-set size is reached.
\end{itemize}

This is a \textbf{greedy} approach, in that at each step it takes the
split that appears best at that moment, without considering the
consequences for the splits beneath it. The resulting tree is therefore
a good tree arrived at through a sequence of locally sensible decisions.
It is not the best one, and the best one is not known.

\subsection{Two features, and the partition of the
space}\label{two-features-and-the-partition-of-the-space}

Before fitting anything, it is worth looking at the two features
together, with the outcome carried by the color of each point.

\begin{Shaded}
\begin{Highlighting}[]
\FunctionTok{ggplot}\NormalTok{(countries, }\FunctionTok{aes}\NormalTok{(ppgdp, pctUrban, }\AttributeTok{colour =}\NormalTok{ infantMortality)) }\SpecialCharTok{+}
  \FunctionTok{geom\_point}\NormalTok{(}\AttributeTok{size =} \DecValTok{2}\NormalTok{)}
\end{Highlighting}
\end{Shaded}

\pandocbounded{\includegraphics[keepaspectratio]{handouts/w03-trees-forests_files/figure-pdf/unnamed-chunk-6-1.pdf}}

Each country is a point located by its income and its urbanization, and
its shade reports the outcome we are trying to predict. The dark points,
where mortality is high, gather in one corner of the plot, and the pale
ones occupy the rest. A tree fitted to these two features is an attempt
to draw rectangles around those regions.

With two features the boxes do become rectangles, and they can still be
drawn.

\begin{Shaded}
\begin{Highlighting}[]
\NormalTok{t2 }\OtherTok{\textless{}{-}} \FunctionTok{tree}\NormalTok{(infantMortality }\SpecialCharTok{\textasciitilde{}}\NormalTok{ ppgdp }\SpecialCharTok{+}\NormalTok{ pctUrban, }\AttributeTok{data =}\NormalTok{ countries,}
           \AttributeTok{control =} \FunctionTok{tree.control}\NormalTok{(}\AttributeTok{nobs =} \FunctionTok{nrow}\NormalTok{(countries), }\AttributeTok{mindev =} \FloatTok{0.002}\NormalTok{))}

\FunctionTok{partition.tree}\NormalTok{(t2, }\AttributeTok{main =} \StringTok{"The partition of the feature space"}\NormalTok{)}
\end{Highlighting}
\end{Shaded}

\pandocbounded{\includegraphics[keepaspectratio]{handouts/w03-trees-forests_files/figure-pdf/unnamed-chunk-7-1.pdf}}

Every rectangle is a leaf, and the number printed within it is the
predicted infant mortality for a country falling there. This is the
image to retain for the remainder of the block, namely that a tree
divides the space of features into boxes and then assigns a single
number to each box.

The \texttt{mindev} argument lets the tree grow a little further than it
otherwise would, since \texttt{tree()} stops splitting quite early by
default and income alone accounts for most of the variation in infant
mortality.

One feature of that figure deserves emphasis, given how much of last
week was devoted to the opposite situation. Nothing in the partition
required the two features to be placed on a common scale. Per capita GDP
runs to roughly a hundred and percent urban to a hundred, and had income
been left in dollars the boundaries would sit at different numbers on
the axis while \textbf{the partition would be identical}. The difficulty
that the rescaling section of last week was designed to address
therefore does not arise here at all.

\subsection{A categorical feature, with no dummy
coding}\label{a-categorical-feature-with-no-dummy-coding}

The UN data include \texttt{group}, which sorts countries into OECD,
African, and other categories. It is a factor rather than a number, and
a tree accepts it as it stands.

\begin{Shaded}
\begin{Highlighting}[]
\NormalTok{m\_tree3 }\OtherTok{\textless{}{-}} \FunctionTok{rpart}\NormalTok{(infantMortality }\SpecialCharTok{\textasciitilde{}}\NormalTok{ ppgdp }\SpecialCharTok{+}\NormalTok{ pctUrban }\SpecialCharTok{+}\NormalTok{ group, }\AttributeTok{data =}\NormalTok{ countries)}

\FunctionTok{rpart.plot}\NormalTok{(m\_tree3, }\AttributeTok{box.palette =} \StringTok{"GnBu"}\NormalTok{)}
\end{Highlighting}
\end{Shaded}

\pandocbounded{\includegraphics[keepaspectratio]{handouts/w03-trees-forests_files/figure-pdf/unnamed-chunk-8-1.pdf}}

There are no dummy variables, no reference category, and no interaction
terms. The tree simply asks which group a country belongs to and directs
it down the appropriate branch, and the split is then printed as the
names of the categories rather than as a coefficient on an indicator.

Consider whether \texttt{group} appears at all, and, if it does, what
happens to income on either side of it. If the tree asks about income at
a different threshold within one group than within another, it has
identified an interaction, which is to say that the relationship between
wealth and child mortality is not the same everywhere. Nothing in the
specification requested this, whereas a linear model recovers such an
interaction only if the researcher thinks in advance to look for it.

Those are the first two items on the list at the head of this handout,
and together they are the principal reason trees are worth knowing.

\subsection{Cost complexity, and the parameter that plays the role of
k}\label{cost-complexity-and-the-parameter-that-plays-the-role-of-k}

Left to itself, the procedure continues splitting until every leaf is
small and homogeneous, which fits the training data very well indeed.
Last week, however, taught us to regard that result with suspicion
rather than satisfaction.

The standard remedy rewrites the problem as

\begin{align}
C(T, y) = L(T, y) + \alpha |T|
\end{align}

where \(L(T, y)\) is the loss from fitting outcome \(y\) with tree
\(T\), and \(|T|\) is the number of nodes. Small loss is always
preferred. Nevertheless, we do not want a tree so complex that it
captures every fluctuation in the training data and consequently
predicts poorly outside it, and the size of the tree is therefore
penalized, with \(\alpha\) determining how heavily that penalty is
applied.

A large \(\alpha\) makes complexity expensive and returns a small tree,
which may in consequence fit poorly, whereas a small \(\alpha\) permits
the tree to grow large. \textbf{Here} \(\alpha\) is analogous to \(k\)
in \(K\)-nearest neighbors, and, as with \(k\), no value is correct in
advance.

The question of how to choose it has the same answer as last week, and
as most weeks: cross-validation.

\subsection{Pruning}\label{pruning}

The practical procedure is to allow the tree to grow first, without any
concern for its size, and then to cut it back to the complexity that
cross-validation prefers. Since \texttt{rpart} cross-validates as it
proceeds, the relevant table is already available.

The operative word is \textbf{grow}, since a tree that is stopped early
never becomes sufficiently complex to overfit, and in that case there is
no trade-off to observe at all. We therefore set \texttt{cp} to
essentially zero and allow the procedure to run until the leaves contain
a handful of countries each.

\begin{Shaded}
\begin{Highlighting}[]
\FunctionTok{set.seed}\NormalTok{(}\DecValTok{7}\NormalTok{)}

\NormalTok{big\_tree }\OtherTok{\textless{}{-}} \FunctionTok{rpart}\NormalTok{(infantMortality }\SpecialCharTok{\textasciitilde{}}\NormalTok{ ppgdp }\SpecialCharTok{+}\NormalTok{ pctUrban }\SpecialCharTok{+}\NormalTok{ group,}
                  \AttributeTok{data =}\NormalTok{ countries,}
                  \AttributeTok{control =} \FunctionTok{rpart.control}\NormalTok{(}\AttributeTok{cp =} \FloatTok{0.000001}\NormalTok{, }\AttributeTok{minsplit =} \DecValTok{10}\NormalTok{,}
                                          \AttributeTok{minbucket =} \DecValTok{5}\NormalTok{, }\AttributeTok{xval =} \DecValTok{10}\NormalTok{))}

\NormalTok{cpt }\OtherTok{\textless{}{-}}\NormalTok{ big\_tree}\SpecialCharTok{$}\NormalTok{cptable}

\FunctionTok{cat}\NormalTok{(}\StringTok{"Leaves in the grown tree:"}\NormalTok{, }\FunctionTok{max}\NormalTok{(cpt[, }\StringTok{"nsplit"}\NormalTok{]) }\SpecialCharTok{+} \DecValTok{1}\NormalTok{, }\StringTok{"}\SpecialCharTok{\textbackslash{}n}\StringTok{"}\NormalTok{)}
\end{Highlighting}
\end{Shaded}

\begin{verbatim}
Leaves in the grown tree: 29 
\end{verbatim}

With roughly a hundred and eighty countries and leaves holding five
apiece, that tree has divided the world into a considerably larger
number of categories than the world contains. It will predict its own
training data very accurately and generalize poorly, which is precisely
the condition we need to be able to observe.

\texttt{plotcp()} draws the cross-validation table that \texttt{rpart}
constructed along the way, which records, for every size of tree from
one split upward, the cross-validated error of a tree of that size.

\begin{Shaded}
\begin{Highlighting}[]
\FunctionTok{plotcp}\NormalTok{(big\_tree)}
\end{Highlighting}
\end{Shaded}

\pandocbounded{\includegraphics[keepaspectratio]{handouts/w03-trees-forests_files/figure-pdf/unnamed-chunk-10-1.pdf}}

This figure carries four pieces of information at once, none of them
labeled obviously, and it is worth taking them in turn.

\textbf{The bottom axis is the complexity parameter,} on a log scale,
running from large on the left to small on the right. A large \(\alpha\)
implies an expensive penalty and therefore a small tree, so that
complexity increases as one moves to the right, which is the same
orientation as last week's \(\log(1/k)\) figure.

\textbf{The top axis is the size of the tree,} measured in number of
leaves. It is the same axis expressed in more intuitive units, and it is
generally the one to attend to.

\textbf{The vertical axis is cross-validated error,} scaled so that a
tree with no splits at all sits at 1. A value of 0.3 indicates that the
tree has removed seventy per cent of the error that would result from
predicting the same number for every country.

\textbf{The vertical bars are one standard error} on either side of each
estimate, and the dotted horizontal line sits one standard error above
the lowest point on the curve. That line supports an alternative rule
for choosing a tree, which we set aside here in favor of the simpler one
below.

The figure should be read in the same way as last week's test-error
curve. The error falls while additional splits are still recovering real
structure, flattens, and then rises once the tree begins fitting
individual countries rather than the pattern among them.

\subsubsection{Choosing where to cut}\label{choosing-where-to-cut}

We take the tree with the lowest cross-validated error, which is the
value the table reports directly.

\begin{Shaded}
\begin{Highlighting}[]
\NormalTok{i\_min   }\OtherTok{\textless{}{-}} \FunctionTok{which.min}\NormalTok{(cpt[, }\StringTok{"xerror"}\NormalTok{])}
\NormalTok{cp\_best }\OtherTok{\textless{}{-}}\NormalTok{ cpt[i\_min, }\StringTok{"CP"}\NormalTok{]}

\FunctionTok{cat}\NormalTok{(}\StringTok{"Lowest CV error at"}\NormalTok{, cpt[i\_min, }\StringTok{"nsplit"}\NormalTok{] }\SpecialCharTok{+} \DecValTok{1}\NormalTok{, }\StringTok{"leaves, cp ="}\NormalTok{,}
\NormalTok{    cp\_best, }\StringTok{"}\SpecialCharTok{\textbackslash{}n}\StringTok{"}\NormalTok{)}
\end{Highlighting}
\end{Shaded}

\begin{verbatim}
Lowest CV error at 5 leaves, cp = 0.01229076 
\end{verbatim}

One caution before pruning. The curve is typically rather flat near its
bottom, so the difference between the best tree and its immediate
neighbors is often smaller than the uncertainty in the estimate. The
minimum is therefore a reasonable choice rather than a precise one, and
a tree one or two splits away would usually serve about as well.

\begin{Shaded}
\begin{Highlighting}[]
\NormalTok{best\_tree }\OtherTok{\textless{}{-}} \FunctionTok{prune}\NormalTok{(big\_tree, }\AttributeTok{cp =}\NormalTok{ cp\_best)}

\FunctionTok{rpart.plot}\NormalTok{(best\_tree, }\AttributeTok{box.palette =} \StringTok{"GnBu"}\NormalTok{, }\AttributeTok{main =} \StringTok{"After pruning"}\NormalTok{)}
\end{Highlighting}
\end{Shaded}

\pandocbounded{\includegraphics[keepaspectratio]{handouts/w03-trees-forests_files/figure-pdf/unnamed-chunk-12-1.pdf}}

That is the whole of the method. Everything established so far, however,
was established with a continuous outcome, one row per country, and a
dataset already familiar from last week, and the second half of this
handout changes all three of those conditions at once.

\section{Part II. The same idea, on a rare binary
outcome}\label{part-ii.-the-same-idea-on-a-rare-binary-outcome}

We now turn to \texttt{onset}, and the algorithm itself barely changes:
instead of averaging the outcome within each box we take the majority
vote, and instead of squared error we use a measure of how mixed each
box is. Everything concerning splitting, greediness, cost complexity,
and pruning carries over unaltered.

Three things do change, however: the data are new, the outcome is
binary, and the outcome is very rare. The third of these occupies most
of what follows, and it shapes nearly every judgment we go on to make
about the model.

\begin{Shaded}
\begin{Highlighting}[]
\NormalTok{url }\OtherTok{\textless{}{-}} \FunctionTok{paste0}\NormalTok{(}\StringTok{"https://raw.githubusercontent.com/babakrezaee/"}\NormalTok{,}
              \StringTok{"AppliedAI{-}PoliSci/refs/heads/main/data/FearonLaitin2003.csv"}\NormalTok{)}

\NormalTok{fl }\OtherTok{\textless{}{-}} \FunctionTok{read.csv}\NormalTok{(url)}

\NormalTok{vars }\OtherTok{\textless{}{-}} \FunctionTok{c}\NormalTok{(}\StringTok{"onset"}\NormalTok{, }\StringTok{"warl"}\NormalTok{, }\StringTok{"gdpenl"}\NormalTok{, }\StringTok{"lpopl1"}\NormalTok{, }\StringTok{"lmtnest"}\NormalTok{, }\StringTok{"ncontig"}\NormalTok{,}
          \StringTok{"Oil"}\NormalTok{, }\StringTok{"nwstate"}\NormalTok{, }\StringTok{"instab"}\NormalTok{, }\StringTok{"polity2l"}\NormalTok{, }\StringTok{"ethfrac"}\NormalTok{, }\StringTok{"relfrac"}\NormalTok{)}

\NormalTok{cw }\OtherTok{\textless{}{-}} \FunctionTok{na.omit}\NormalTok{(fl[, vars])}
\NormalTok{cw }\OtherTok{\textless{}{-}}\NormalTok{ cw[cw}\SpecialCharTok{$}\NormalTok{onset }\SpecialCharTok{\%in\%} \FunctionTok{c}\NormalTok{(}\DecValTok{0}\NormalTok{, }\DecValTok{1}\NormalTok{), ]}

\FunctionTok{nrow}\NormalTok{(cw)}
\end{Highlighting}
\end{Shaded}

\begin{verbatim}
[1] 6326
\end{verbatim}

These are Fearon and Laitin's civil war data, and \texttt{onset} takes
the value 1 if a civil war began in that country in that year.

The covariates are drawn from Fearon and Laitin's own specification,
although the models below use eight of them rather than the full set.
The reduction is made for legibility rather than for substance: a tree
fitted on eleven features produces a figure that nobody can read on a
screen, and the point of Part II is to be able to see what the algorithm
did. The three variables set aside, prior war, non-contiguity, and
political instability, remain in \texttt{cw}, so that restoring them is
a matter of adding three terms to a formula. Doing so, and observing how
much the figures deteriorate, is the first thing worth trying in the
workshop.

Two features of that block are worth noting, since each can consume an
hour when encountered for the first time in one's own project.
\texttt{Oil} is capitalized in this file while nothing else is, and R
will report that the object does not exist rather than inferring what
was intended. \texttt{onset}, meanwhile, carries an occasional code
other than 0 or 1, for wars already under way when a country entered the
data, and we therefore retain only the two values we intend to model.

\textbf{One structural difference from Part I should be acknowledged
here, since we are about to accept it rather than solve it.} Every row
is a country in a single year, from 1945--1999, so that the United
States appears fifty-five times with identical terrain, identical
colonial history, and an income that barely moves between consecutive
rows. Part I had one row per country and did not confront this. Here the
structure cannot be collapsed away, because a civil war begins in a
particular year and cannot be averaged into a country's post-war
summary. The point returns at the end of this handout, where it acquires
a name.

The outcome can now be counted, and the resulting figure is the one on
which the remainder of this handout turns.

\begin{Shaded}
\begin{Highlighting}[]
\FunctionTok{table}\NormalTok{(cw}\SpecialCharTok{$}\NormalTok{onset)}
\end{Highlighting}
\end{Shaded}

\begin{verbatim}

   0    1 
6221  105 
\end{verbatim}

\begin{Shaded}
\begin{Highlighting}[]
\FunctionTok{round}\NormalTok{(}\FunctionTok{mean}\NormalTok{(cw}\SpecialCharTok{$}\NormalTok{onset), }\DecValTok{4}\NormalTok{)}
\end{Highlighting}
\end{Shaded}

\begin{verbatim}
[1] 0.0166
\end{verbatim}

Civil wars begin in roughly two country-years in a hundred, and that
figure is the single most consequential quantity on this page, since
nearly everything that follows is shaped by it.

The immediate implication can be stated directly by considering a model
that consists of one sentence, and that contains no data, no fitting,
and no thought whatever:

\begin{quote}
Predict that no civil war will begin, anywhere, ever.
\end{quote}

That model is correct approximately ninety-eight per cent of the time,
and reported as accuracy it appears excellent. It is nevertheless
worthless, since the only cases of any substantive interest are
precisely the ones it gets wrong, in every instance, by construction.

The model is worth retaining for comparison, since we are about to build
something considerably more sophisticated and will then need to
establish whether it improves upon this.

\begin{Shaded}
\begin{Highlighting}[]
\FunctionTok{set.seed}\NormalTok{(}\DecValTok{7}\NormalTok{)}

\NormalTok{n        }\OtherTok{\textless{}{-}} \FunctionTok{nrow}\NormalTok{(cw)}
\NormalTok{train\_id }\OtherTok{\textless{}{-}} \FunctionTok{sample}\NormalTok{(}\DecValTok{1}\SpecialCharTok{:}\NormalTok{n, }\FunctionTok{floor}\NormalTok{(}\FloatTok{0.7} \SpecialCharTok{*}\NormalTok{ n))}

\NormalTok{Train }\OtherTok{\textless{}{-}}\NormalTok{ cw[train\_id, ]}
\NormalTok{Test  }\OtherTok{\textless{}{-}}\NormalTok{ cw[}\SpecialCharTok{{-}}\NormalTok{train\_id, ]}
\end{Highlighting}
\end{Shaded}

\subsection{One tree on onset}\label{one-tree-on-onset}

\begin{Shaded}
\begin{Highlighting}[]
\NormalTok{m\_stump }\OtherTok{\textless{}{-}} \FunctionTok{rpart}\NormalTok{(}\FunctionTok{factor}\NormalTok{(onset) }\SpecialCharTok{\textasciitilde{}}\NormalTok{ gdpenl }\SpecialCharTok{+}\NormalTok{ lpopl1 }\SpecialCharTok{+}\NormalTok{ lmtnest }\SpecialCharTok{+}\NormalTok{ Oil }\SpecialCharTok{+}
\NormalTok{                   nwstate }\SpecialCharTok{+}\NormalTok{ polity2l }\SpecialCharTok{+}\NormalTok{ ethfrac }\SpecialCharTok{+}\NormalTok{ relfrac,}
                 \AttributeTok{data =}\NormalTok{ Train, }\AttributeTok{method =} \StringTok{"class"}\NormalTok{,}
                 \AttributeTok{control =} \FunctionTok{rpart.control}\NormalTok{(}\AttributeTok{cp =} \FloatTok{0.004}\NormalTok{, }\AttributeTok{minbucket =} \DecValTok{25}\NormalTok{))}

\FunctionTok{rpart.plot}\NormalTok{(m\_stump, }\AttributeTok{box.palette =} \StringTok{"GnBu"}\NormalTok{)}
\end{Highlighting}
\end{Shaded}

\pandocbounded{\includegraphics[keepaspectratio]{handouts/w03-trees-forests_files/figure-pdf/unnamed-chunk-16-1.pdf}}

That figure is not an error, and it is worth pausing over rather than
correcting immediately. The tree has declined to split at all. It
reports a single node containing every case, predicts no onset for all
of them, and attaches a probability of 0.02, which is simply the base
rate we computed a moment ago.

The reason lies in the criterion. A split is retained only if it
improves the tree's classification sufficiently to justify its cost, and
on an outcome occurring in two per cent of cases no split improves
classification whatever. Predicting no war everywhere is already correct
ninety-eight per cent of the time, and no partition of the feature space
available to \texttt{rpart} does better than that. The algorithm
therefore concluded, entirely correctly on its own terms, that the best
available tree is the one that asks no questions at all.

We have thus encountered the central difficulty of this half of the
handout before fitting anything of interest. The one-sentence model set
out above has not merely proved difficult to beat. It has turned out to
be the model that \texttt{rpart} selects when left to its own devices,
and it arrived not through any failure of the software but through the
faithful execution of the objective we gave it.

\subsubsection{Telling the algorithm that the rare class
matters}\label{telling-the-algorithm-that-the-rare-class-matters}

The standard remedy is to instruct the procedure to treat the two
classes as equally important, notwithstanding that they are not equally
common.

\begin{Shaded}
\begin{Highlighting}[]
\NormalTok{m\_class }\OtherTok{\textless{}{-}} \FunctionTok{rpart}\NormalTok{(}\FunctionTok{factor}\NormalTok{(onset) }\SpecialCharTok{\textasciitilde{}}\NormalTok{ gdpenl }\SpecialCharTok{+}\NormalTok{ lpopl1 }\SpecialCharTok{+}\NormalTok{ lmtnest }\SpecialCharTok{+}\NormalTok{ Oil }\SpecialCharTok{+}
\NormalTok{                   nwstate }\SpecialCharTok{+}\NormalTok{ polity2l }\SpecialCharTok{+}\NormalTok{ ethfrac }\SpecialCharTok{+}\NormalTok{ relfrac,}
                 \AttributeTok{data =}\NormalTok{ Train, }\AttributeTok{method =} \StringTok{"class"}\NormalTok{,}
                 \AttributeTok{parms   =} \FunctionTok{list}\NormalTok{(}\AttributeTok{prior =} \FunctionTok{c}\NormalTok{(}\FloatTok{0.5}\NormalTok{, }\FloatTok{0.5}\NormalTok{)),}
                 \AttributeTok{control =} \FunctionTok{rpart.control}\NormalTok{(}\AttributeTok{cp =} \FloatTok{0.002}\NormalTok{, }\AttributeTok{minbucket =} \DecValTok{20}\NormalTok{))}

\FunctionTok{rpart.plot}\NormalTok{(m\_class, }\AttributeTok{box.palette =} \StringTok{"GnBu"}\NormalTok{)}
\end{Highlighting}
\end{Shaded}

\pandocbounded{\includegraphics[keepaspectratio]{handouts/w03-trees-forests_files/figure-pdf/unnamed-chunk-17-1.pdf}}

The \texttt{prior} argument instructs \texttt{rpart} to proceed as
though onsets and non-onsets were equally frequent, which has the effect
of making each onset count for roughly fifty times what it counted for
before. Splits that separate a handful of wars from a great many
non-wars now improve the objective, and the tree consequently grows. The
\texttt{method\ =\ "class"} argument remains the only difference in kind
from Part I, and the resulting figure is read exactly as the earlier
ones were.

Two observations about that argument deserve emphasis, since the second
of them is the more important.

First, the reweighting is a decision rather than a correction. Nothing
in the data licensed it, and no diagnostic recommended it. We supplied
it because we judged that a missed war matters more than a false alarm,
which is a substantive judgment about consequences and not a statistical
one. The tree without it was not broken; it was answering a different
question, one we had posed carelessly.

Second, and more consequentially, \textbf{we have reweighted the tree
and we have not reweighted the forest.} That asymmetry is deliberate,
and its consequences appear shortly.

\subsection{Why one tree is not
enough}\label{why-one-tree-is-not-enough}

We arrive here at the problem that motivates the remainder of the week.
Fitting the same tree to slightly different data does not produce a
slightly different tree; it frequently produces a different tree
altogether.

\begin{Shaded}
\begin{Highlighting}[]
\FunctionTok{par}\NormalTok{(}\AttributeTok{mfrow =} \FunctionTok{c}\NormalTok{(}\DecValTok{1}\NormalTok{, }\DecValTok{3}\NormalTok{))}

\FunctionTok{set.seed}\NormalTok{(}\DecValTok{7}\NormalTok{)}

\ControlFlowTok{for}\NormalTok{ (i }\ControlFlowTok{in} \DecValTok{1}\SpecialCharTok{:}\DecValTok{3}\NormalTok{) \{}
\NormalTok{  draw }\OtherTok{\textless{}{-}}\NormalTok{ Train[}\FunctionTok{sample}\NormalTok{(}\DecValTok{1}\SpecialCharTok{:}\FunctionTok{nrow}\NormalTok{(Train), }\FunctionTok{nrow}\NormalTok{(Train), }\AttributeTok{replace =} \ConstantTok{TRUE}\NormalTok{), ]}

\NormalTok{  t\_i }\OtherTok{\textless{}{-}} \FunctionTok{rpart}\NormalTok{(}\FunctionTok{factor}\NormalTok{(onset) }\SpecialCharTok{\textasciitilde{}}\NormalTok{ gdpenl }\SpecialCharTok{+}\NormalTok{ lpopl1 }\SpecialCharTok{+}\NormalTok{ lmtnest }\SpecialCharTok{+}\NormalTok{ Oil }\SpecialCharTok{+}
\NormalTok{                 nwstate }\SpecialCharTok{+}\NormalTok{ polity2l }\SpecialCharTok{+}\NormalTok{ ethfrac }\SpecialCharTok{+}\NormalTok{ relfrac,}
               \AttributeTok{data =}\NormalTok{ draw, }\AttributeTok{method =} \StringTok{"class"}\NormalTok{,}
               \AttributeTok{parms   =} \FunctionTok{list}\NormalTok{(}\AttributeTok{prior =} \FunctionTok{c}\NormalTok{(}\FloatTok{0.5}\NormalTok{, }\FloatTok{0.5}\NormalTok{)),}
               \AttributeTok{control =} \FunctionTok{rpart.control}\NormalTok{(}\AttributeTok{cp =} \FloatTok{0.004}\NormalTok{, }\AttributeTok{minbucket =} \DecValTok{20}\NormalTok{))}

  \FunctionTok{rpart.plot}\NormalTok{(t\_i, }\AttributeTok{box.palette =} \StringTok{"GnBu"}\NormalTok{, }\AttributeTok{main =} \FunctionTok{paste}\NormalTok{(}\StringTok{"Sample"}\NormalTok{, i))}
\NormalTok{\}}
\end{Highlighting}
\end{Shaded}

\pandocbounded{\includegraphics[keepaspectratio]{handouts/w03-trees-forests_files/figure-pdf/unnamed-chunk-18-1.pdf}}

Three samples drawn from the same data yield three trees, and the
variable sitting at the top of each is the one to watch. If the first
split changes, everything beneath it changes with it, since every
subsequent question is asked only of the cases that the first question
directed that way.

The vocabulary for this was established last week, and the term is
\textbf{variance}, in the sense that the model responds to the
particular cases placed before it rather than to the structure
underlying them. A single tree has low bias, since given sufficient
depth it can represent almost any shape, and high variance, since which
shape it represents depends heavily on the sample. The greedy procedure
aggravates the problem, because a single early decision taken on a
narrow margin propagates through everything below it.

A single tree is thus readable and unreliable at the same time, and the
response, perhaps counterintuitively, is to stop reading it.

\subsection{From bagging to random
forests}\label{from-bagging-to-random-forests}

\textbf{Bagging}, an abbreviation of bootstrap aggregating, is the first
idea. The data are resampled with replacement a sufficient number of
times that trees with good explanatory power are captured, a tree is
fitted to each sample, and their predictions are averaged. The procedure
is computationally intensive, and it works for the reason averaging
generally works: individually noisy estimates, once averaged, are less
noisy, provided their errors are not all the same error.

That proviso, as it happens, is precisely where the method acquires its
name. \textbf{Random forests begin from bagging and add a second
randomization.} Rather than scanning all the features at each split, the
algorithm samples a subset of \(m\) variables and must then select its
split from among those alone.

The rationale is precise. If income is the strongest predictor, for
instance, then every bagged tree splits on income first, and the trees
therefore resemble one another closely and err in the same places.
Requiring most trees to proceed without it at any given split breaks
that correlation. The individual trees thereby become \textbf{worse},
and yet their average becomes \textbf{better}.

This is the proposition to carry out of the week. A forest is not an
attempt to build one excellent model. It is an attempt to build many
mediocre models whose errors point in different directions, and then to
let them vote.

Two parameters govern the procedure: \(B\) is the number of bootstrapped
trees, and \(m\) is the number of variables sampled at each split. A
common choice is \(m = \sqrt{p}\), where \(p\) is the number of
features. Notably, when \(m = p\) the random forest reduces to bagging.

\subsection{Fitting the forest}\label{fitting-the-forest}

Last week's claim was that changing one string changes the algorithm,
and this is that change.

\begin{Shaded}
\begin{Highlighting}[]
\NormalTok{Train}\SpecialCharTok{$}\NormalTok{onset\_f }\OtherTok{\textless{}{-}} \FunctionTok{factor}\NormalTok{(}\FunctionTok{ifelse}\NormalTok{(Train}\SpecialCharTok{$}\NormalTok{onset }\SpecialCharTok{==} \DecValTok{1}\NormalTok{, }\StringTok{"war"}\NormalTok{, }\StringTok{"no\_war"}\NormalTok{),}
                        \AttributeTok{levels =} \FunctionTok{c}\NormalTok{(}\StringTok{"no\_war"}\NormalTok{, }\StringTok{"war"}\NormalTok{))}
\NormalTok{Test}\SpecialCharTok{$}\NormalTok{onset\_f  }\OtherTok{\textless{}{-}} \FunctionTok{factor}\NormalTok{(}\FunctionTok{ifelse}\NormalTok{(Test}\SpecialCharTok{$}\NormalTok{onset }\SpecialCharTok{==} \DecValTok{1}\NormalTok{, }\StringTok{"war"}\NormalTok{, }\StringTok{"no\_war"}\NormalTok{),}
                        \AttributeTok{levels =} \FunctionTok{c}\NormalTok{(}\StringTok{"no\_war"}\NormalTok{, }\StringTok{"war"}\NormalTok{))}

\FunctionTok{set.seed}\NormalTok{(}\DecValTok{7}\NormalTok{)}

\NormalTok{rf }\OtherTok{\textless{}{-}} \FunctionTok{train}\NormalTok{(onset\_f }\SpecialCharTok{\textasciitilde{}}\NormalTok{ gdpenl }\SpecialCharTok{+}\NormalTok{ lpopl1 }\SpecialCharTok{+}\NormalTok{ lmtnest }\SpecialCharTok{+}\NormalTok{ Oil }\SpecialCharTok{+}
\NormalTok{              nwstate }\SpecialCharTok{+}\NormalTok{ polity2l }\SpecialCharTok{+}\NormalTok{ ethfrac }\SpecialCharTok{+}\NormalTok{ relfrac,}
            \AttributeTok{data      =}\NormalTok{ Train,}
            \AttributeTok{method    =} \StringTok{"rf"}\NormalTok{,}
            \AttributeTok{ntree     =} \DecValTok{300}\NormalTok{,}
            \AttributeTok{trControl =} \FunctionTok{trainControl}\NormalTok{(}\AttributeTok{method =} \StringTok{"cv"}\NormalTok{, }\AttributeTok{number =} \DecValTok{5}\NormalTok{),}
            \AttributeTok{tuneGrid  =} \FunctionTok{data.frame}\NormalTok{(}\AttributeTok{mtry =} \FunctionTok{c}\NormalTok{(}\DecValTok{2}\NormalTok{, }\DecValTok{3}\NormalTok{, }\DecValTok{5}\NormalTok{)))}

\NormalTok{rf}
\end{Highlighting}
\end{Shaded}

\begin{verbatim}
Random Forest 

4428 samples
   8 predictor
   2 classes: 'no_war', 'war' 

No pre-processing
Resampling: Cross-Validated (5 fold) 
Summary of sample sizes: 3543, 3542, 3542, 3543, 3542 
Resampling results across tuning parameters:

  mtry  Accuracy   Kappa       
  2     0.9805785  -0.003979084
  3     0.9794498  -0.005356815
  5     0.9789981  -0.005866209

Accuracy was used to select the optimal model using the largest value.
The final value used for the model was mtry = 2.
\end{verbatim}

Everything apart from \texttt{method\ =\ "rf"} and the grid over
\texttt{mtry} is the structure established last week: a formula,
training data, a resampling scheme, and a set of values to try. The
outcome had to be converted to a factor, since that is how
\texttt{caret} is informed that the problem is one of classification
rather than regression. Note also that nothing was rescaled, and that
nothing needed to be.

Note as well what is \textbf{not} in that call. There is no
\texttt{prior}, no class weighting, and no instruction of any kind that
an onset matters more than a non-onset. We supplied precisely such an
instruction to the single tree a few pages ago, and we have deliberately
withheld it here. What follows is the consequence.

This chunk takes a minute or two to run, which is worth knowing before
you begin.

\subsubsection{Predicting for cases the model has never
seen}\label{predicting-for-cases-the-model-has-never-seen}

\texttt{train()} refits the model on the whole of the training data once
tuning is complete, so that the object is a fitted model and prediction
proceeds from it directly.

\begin{Shaded}
\begin{Highlighting}[]
\NormalTok{pred }\OtherTok{\textless{}{-}} \FunctionTok{predict}\NormalTok{(rf, }\AttributeTok{newdata =}\NormalTok{ Test)}

\FunctionTok{head}\NormalTok{(pred)}
\end{Highlighting}
\end{Shaded}

\begin{verbatim}
[1] no_war no_war no_war no_war no_war no_war
Levels: no_war war
\end{verbatim}

Prediction should proceed from \texttt{rf} rather than from
\texttt{rf\$finalModel}, since the final model holds the fitted forest
but retains no record of any preprocessing. Calling it directly may
therefore supply it with data in the wrong form, and it will return an
answer regardless.

\subsection{The accuracy trap}\label{the-accuracy-trap}

\begin{Shaded}
\begin{Highlighting}[]
\NormalTok{cm }\OtherTok{\textless{}{-}} \FunctionTok{table}\NormalTok{(}\AttributeTok{Predicted =}\NormalTok{ pred, }\AttributeTok{Actual =}\NormalTok{ Test}\SpecialCharTok{$}\NormalTok{onset\_f)}
\NormalTok{cm}
\end{Highlighting}
\end{Shaded}

\begin{verbatim}
         Actual
Predicted no_war  war
   no_war   1865   30
   war         3    0
\end{verbatim}

\begin{Shaded}
\begin{Highlighting}[]
\NormalTok{accuracy }\OtherTok{\textless{}{-}} \FunctionTok{sum}\NormalTok{(}\FunctionTok{diag}\NormalTok{(cm)) }\SpecialCharTok{/} \FunctionTok{sum}\NormalTok{(cm)}
\FunctionTok{round}\NormalTok{(accuracy, }\DecValTok{4}\NormalTok{)}
\end{Highlighting}
\end{Shaded}

\begin{verbatim}
[1] 0.9826
\end{verbatim}

The number itself should be considered first, and it is very high, such
that in the abstract of an article it would attract no particular
scrutiny.

The one-sentence model from the start of Part II, the one containing no
data at all, provides the comparison.

\begin{Shaded}
\begin{Highlighting}[]
\NormalTok{baseline }\OtherTok{\textless{}{-}} \FunctionTok{mean}\NormalTok{(Test}\SpecialCharTok{$}\NormalTok{onset\_f }\SpecialCharTok{==} \StringTok{"no\_war"}\NormalTok{)}
\FunctionTok{round}\NormalTok{(baseline, }\DecValTok{4}\NormalTok{)}
\end{Highlighting}
\end{Shaded}

\begin{verbatim}
[1] 0.9842
\end{verbatim}

Whatever the forest achieved must be read against that figure, and the
gap between them is small or absent. We have grown nine hundred trees,
tuned \texttt{mtry} by cross-validation, and arrived at approximately
the proposition that nothing ever happens.

The table itself, rather than the summary, is where the substance lies.
Reading down the column of actual wars, the question is how many of the
onsets in the test set the model predicted, and the answer is generally
close to none.

The forest is not malfunctioning. However, we requested accuracy, and
accuracy on a two per cent outcome is maximized by predicting zero
indefinitely, so that the model did precisely as it was instructed.

This is the stump from the beginning of Part II, returning in a more
elaborate form. The single tree refused to split because no split
improved accuracy; the forest grew nine hundred trees and then predicted
almost nothing, for the same reason. We repaired the tree by telling it
that the rare class mattered, and we declined to repair the forest, and
the two figures should be read together. The disease is identical in
both cases, and it lies in the objective rather than in the algorithm.
\textbf{\texttt{caret} tuned \texttt{mtry} so as to maximize accuracy,
which is to say that it tuned toward the useless model, and reported
success in doing so.}

Nothing in the output announced this. The accuracy appeared excellent,
the cross-validation was correctly specified, and the test set was
genuinely held out. Every methodological requirement was satisfied, and
the result is a model that has not once predicted the outcome it was
built to predict.

This is the argument of the second reading for this week (Muchlinski et
al. 2016), and it is the reason that article exists. Week 6 takes up the
repair in full, covering precision, recall, the F1 score, thresholds,
and class weighting, together with the question of what to do when one
error costs substantially more than the other. For the present, the
diagnosis matters more than the cure.

\subsubsection{The four cells}\label{the-four-cells}

The table printed above is a \textbf{confusion matrix}, and naming its
cells requires first determining which class counts as
\textbf{positive}, which is a decision made by the analyst rather than a
property of the data. Here the positive class is war onset.

{\def\LTcaptype{none} % do not increment counter
\begin{longtable}[]{@{}lll@{}}
\toprule\noalign{}
& Actually war & Actually no war \\
\midrule\noalign{}
\endhead
\bottomrule\noalign{}
\endlastfoot
\textbf{Predicted war} & True positive & False positive \\
\textbf{Predicted no war} & False negative & True negative \\
\end{longtable}
}

A \textbf{false positive} is a country flagged as heading toward civil
war that does not in fact experience one, whereas a \textbf{false
negative} is a civil war that nobody anticipated. These are not the same
error, and no serious application treats them as equivalent. Which of
them one would rather commit is a political and ethical judgment about
consequences, and accuracy collapses the two into a single number
precisely by declining to pose the question.

\subsection{Variable importance, and what it does not
mean}\label{variable-importance-and-what-it-does-not-mean}

We surrendered the readable tree, and variable importance is what is
returned to us in its place.

\begin{Shaded}
\begin{Highlighting}[]
\FunctionTok{varImpPlot}\NormalTok{(rf}\SpecialCharTok{$}\NormalTok{finalModel, }\AttributeTok{main =} \StringTok{"Variable importance"}\NormalTok{)}
\end{Highlighting}
\end{Shaded}

\pandocbounded{\includegraphics[keepaspectratio]{handouts/w03-trees-forests_files/figure-pdf/unnamed-chunk-23-1.pdf}}

The underlying idea is straightforward. For each feature, one measures
the extent to which the forest's performance depends upon it, either by
totalling the improvement that feature produced across all the splits in
which it was used, or by permuting its values and observing how far the
predictions deteriorate. Thus a feature that can be scrambled without
consequence was evidently not contributing much.

Since the forest is bootstrapped, features that genuinely matter will on
average be selected and used more frequently than others, and this is
what makes the ranking meaningful at all.

The limits of the measure, however, matter more than the figure itself.

\textbf{Importance is not effect size.} The plot has no interpretable
units. A feature twice as high as another is not twice as influential in
any substantive sense, and the numbers are not comparable across models
or datasets.

\textbf{Importance carries no sign.} Nothing here indicates whether
higher income raises or lowers the risk of onset. A feature may sit at
the top of the plot with a relationship that is positive, negative, or
U-shaped, and the ranking cannot distinguish among them.

\textbf{Correlated features divide their importance between them.} Where
two features carry substantially the same information, the forest uses
them interchangeably, and each appears half as important as either would
appear alone. A feature may therefore look unimportant because it has a
twin rather than because it is irrelevant. In these data, where income,
population, and several regime measures are entangled with one another,
the possibility is not hypothetical.

\textbf{Impurity-based importance is biased} toward features with many
distinct values, so that continuous variables tend to outrank binary
ones for reasons unrelated to the world. It is worth comparing where
\texttt{Oil} and \texttt{nwstate}, which are binary, sit relative to
\texttt{gdpenl}, which is not.

\textbf{None of it, finally, is causal.} Importance describes what the
model used in order to predict. It says nothing about what would follow
from changing a feature, since nothing in the procedure ever posed that
question.

The last point has a practical edge, in that variable importance is
regularly presented, in policy briefs and in press coverage, as
identifying the drivers of a problem. What it identifies is the set of
predictors that a particular fitted model found useful on a particular
dataset. These are different claims, and the distance between them is
where a good deal of poor AI-informed policy analysis is located.

\section{Prediction, forecasting, and early
warning}\label{prediction-forecasting-and-early-warning}

Week 1 distinguished prediction from forecasting: prediction estimates
an outcome for a case the model has not seen, whereas forecasting
estimates one that has not yet occurred. Everything above is prediction,
since our test set consists of country-years drawn at random from the
same period as the training set.

The third reading for this week concerns the other activity ({Hegre et
al.} 2019). ViEWS is an operating early-warning system that publishes
monthly forecasts of political violence for months that have not yet
occurred, and the difference is not merely presentational.

As soon as the target is the future, a random split becomes
indefensible. If the model trains on 1995 and 1999 and is tested on
1997, it has observed the years on either side of the case it is
predicting, and it will make use of them. Forecasting therefore requires
splitting on time: training up to a date, forecasting beyond it, and
accepting that one cannot look ahead. Nearly every published forecasting
model that appeared too good to be true was too good on account of some
version of this.

A subtler problem remains, which cross-validation does not address, in
that cross-validation assumes cases to be interchangeable while
country-years plainly are not. Sudan in 1996 is not independent of Sudan
in 1995, and a fold that holds one while the other trains the model is
leaking information. This is the structural difference flagged at the
start of Part II, and it has a name: the observations are not
independent, while the standard resampling procedures assume that they
are. Week 6 returns to the problem and to the ways in which it can be
mitigated.

\section{The politics of early
warning}\label{the-politics-of-early-warning}

A forecast is not merely an analysis that happens to concern the future.
It is, rather, an object produced in order to be acted upon, and that
purpose changes what its errors mean.

Suppose a system of this kind flags a country as being at elevated risk.
Consequences follow, in investment decisions, aid allocations, insurance
premiums, travel advisories, diplomatic posture, and on occasion in the
argument for intervention. The forecast thereby enters the politics of
the country it describes.

That observation should be combined with the base rate. Thus, where the
outcome occurs in two per cent of cases, even a model with genuinely
good discrimination will produce considerably more false alarms than
true warnings, simply because there are so many more non-events
available to be wrong about. A country may therefore be flagged, bear
the consequences of having been flagged, and have been a false positive
throughout. It will not be informed of this, and the system has no means
of establishing it either, until the year passes without incident and
the case leaves the window.

Moreover, a forecast can act upon the world it describes. If a warning
of instability withdraws investment, and withdrawn investment raises the
likelihood of instability, then the model has participated in the
outcome it predicted and will appear more accurate for having done so.
This is not a defect in the statistics. It is, instead, what occurs
whenever a prediction about people is published to people who have the
power to respond to it.

The closing question last week concerned who defines similarity, and the
closing question this week concerns who defines the threshold. Somewhere
between the forest's vote share and a red mark on a map, a number is
selected that determines which countries are flagged, and that number
sets the ratio of false alarms to missed wars. It is among the most
consequential decisions in the entire pipeline, it is rarely made by
anybody who bears its consequences, and it almost never appears in the
write-up.

\section{Reading}\label{reading-2}

Fearon, J. D., and Laitin, D. D. 2003. ``Ethnicity, Insurgency, and
Civil War.'' \emph{American Political Science Review} 97(1): 75--90.
(Fearon and Laitin 2003)

These are the data used above and the specification fitted above. The
article should be read for its argument about the causes of civil war,
and it is worth noting throughout that it is an explanatory article. Its
authors are not attempting to predict onset; they are attempting to
explain it, and the difference between those two objectives is the
thread running through this week.

Muchlinski, D., Siroky, D., He, J., and Kocher, M. 2016. ``Comparing
Random Forest with Logistic Regression for Predicting Class-Imbalanced
Civil War Onset Data.'' \emph{Political Analysis} 24(1): 87--103.
(Muchlinski et al. 2016)

This one should be read for its argument about rare outcomes rather than
for its implementation, and, if time permits, alongside the critical
exchange it provoked, since the disagreement is instructive about what a
performance comparison can and cannot settle.

Hegre, H., et al.~2019. ``ViEWS: A Political Violence Early-Warning
System.'' \emph{Journal of Peace Research} 56(2): 155--174. ({Hegre et
al.} 2019)

This should be read as a description of a system that is actually in
use, rather than as a methods article. Particular attention is worth
giving to what its authors say about how the forecasts are intended to
be interpreted, and by whom.

\textbf{Recommended, not required.}

Pinckney, J., and RezaeeDaryakenari, B. 2022. ``When the Levee Breaks: A
Forecasting Model of Violent and Nonviolent Dissent.''
\emph{International Interactions} 48(5): 997--1026. (Pinckney and
RezaeeDaryakenari 2022)

This article forecasts both armed and unarmed conflict at the
country-year level, and it does so by training eight machine-learning
algorithms and combining them into five forecasting ensembles. The
relevance to this week is direct, in that the logic of the random forest
reappears one level higher: where a forest averages many trees whose
errors point in different directions, an ensemble forecast averages many
algorithms on the same reasoning. It is also worth reading for the
contrast with the cases above, since the outcome of interest includes
nonviolent mobilization rather than civil war alone, and unarmed
uprisings are neither as rare nor as well measured.

\section{Before the workshop}\label{before-the-workshop-1}

Install \texttt{rpart}, \texttt{rpart.plot}, \texttt{tree},
\texttt{randomForest}, and \texttt{caret}, and run every chunk in this
handout once. The forest takes a minute or two, which is better
discovered at home than in the room.

\section{Questions for the session}\label{questions-for-the-session-2}

Come with a view on at least two of these.

\begin{enumerate}
\def\labelenumi{\arabic{enumi}.}
\item
  Fearon and Laitin set out to explain civil war, whereas we used their
  variables to predict it. What did we learn that they did not, and what
  did they establish that our model cannot?
\item
  A single tree is readable and unreliable, while a forest is reliable
  and unreadable. For a political scientist, which of these failures is
  the more serious, and does the answer depend on the purpose the model
  is to serve?
\item
  The forest deliberately makes each tree worse, by withholding features
  from it, in order that the ensemble should be better. Is there
  anything analogous in the way political institutions aggregate
  judgments? Committees, juries, federal structures, and multiple veto
  players are all admissible.
\item
  Our model achieved very high accuracy and predicted almost no wars,
  and every technical step was performed correctly. What does this
  suggest about the relationship between methodological correctness and
  usefulness?
\item
  Variable importance is regularly reported in policy documents as
  identifying the drivers of a problem. Write the one sentence you would
  add to such a document in order to make it honest, and state what
  would be lost in the process.
\item
  Who should select the threshold at which a forecast becomes a warning,
  and to whom should that selection be explained? Note that those
  bearing the cost of a false alarm and those bearing the cost of a
  missed war are generally not the same people.
\end{enumerate}

\chapter{Text as Data}\label{text-as-data}

\emph{In preparation.}

\chapter{Large Language Models}\label{large-language-models}

\emph{In preparation.}

\chapter{Evaluating AI Systems}\label{evaluating-ai-systems}

\emph{In preparation.}

\chapter{Final Project}\label{final-project}

\emph{In preparation.}

\bookmarksetup{startatroot}

\chapter*{References}\label{references}
\addcontentsline{toc}{chapter}{References}

\markboth{References}{References}

\protect\phantomsection\label{refs}
\begin{CSLReferences}{1}{1}
\bibitem[\citeproctext]{ref-cranmer2017}
Cranmer, Skyler J., and Bruce A. Desmarais. 2017. {``What Can We Learn
from Predictive Modeling?''} \emph{Political Analysis} 25 (2): 145--66.
\url{https://doi.org/10.1017/pan.2017.3}.

\bibitem[\citeproctext]{ref-fearon2003}
Fearon, James D., and David D. Laitin. 2003. {``Ethnicity, Insurgency,
and Civil War.''} \emph{American Political Science Review} 97 (1):
75--90. \url{https://doi.org/10.1017/S0003055403000534}.

\bibitem[\citeproctext]{ref-hegre2019}
{Hegre, Håvard et al.} 2019. {``{ViEWS}: A Political Violence
Early-Warning System.''} \emph{Journal of Peace Research} 56 (2):
155--74. \url{https://doi.org/10.1177/0022343319823860}.

\bibitem[\citeproctext]{ref-hobbes1651}
Hobbes, Thomas. 1651. \emph{Leviathan}. Andrew Crooke.

\bibitem[\citeproctext]{ref-hofman2017}
Hofman, Jake M., Amit Sharma, and Duncan J. Watts. 2017. {``Prediction
and Explanation in Social Systems.''} \emph{Science} 355 (6324):
486--88. \url{https://doi.org/10.1126/science.aal3856}.

\bibitem[\citeproctext]{ref-james2021}
James, Gareth, Daniela Witten, Trevor Hastie, and Robert Tibshirani.
2021. \emph{An Introduction to Statistical Learning: With Applications
in {R}}. 2nd ed. Springer.

\bibitem[\citeproctext]{ref-lakatos1970}
Lakatos, Imre. 1970. {``Falsification and the Methodology of Scientific
Research Programmes.''} In \emph{Criticism and the Growth of Knowledge},
edited by Imre Lakatos and Alan Musgrave. Cambridge University Press.

\bibitem[\citeproctext]{ref-muchlinski2016}
Muchlinski, David, David Siroky, Jingrui He, and Matthew Kocher. 2016.
{``Comparing Random Forest with Logistic Regression for Predicting
Class-Imbalanced Civil War Onset Data.''} \emph{Political Analysis} 24
(1): 87--103. \url{https://doi.org/10.1093/pan/mpv024}.

\bibitem[\citeproctext]{ref-pinckney2022}
Pinckney, Jonathan, and Babak RezaeeDaryakenari. 2022. {``When the Levee
Breaks: A Forecasting Model of Violent and Nonviolent Dissent.''}
\emph{International Interactions} 48 (5): 997--1026.
\url{https://doi.org/10.1080/03050629.2022.2090933}.

\bibitem[\citeproctext]{ref-shmueli2010}
Shmueli, Galit. 2010. {``To Explain or to Predict?''} \emph{Statistical
Science} 25 (3): 289--310. \url{https://doi.org/10.1214/10-STS330}.

\bibitem[\citeproctext]{ref-simon1947}
Simon, Herbert A. 1947. \emph{Administrative Behavior: A Study of
Decision-Making Processes in Administrative Organization}. Macmillan.

\bibitem[\citeproctext]{ref-simon1969}
Simon, Herbert A. 1969. \emph{The Sciences of the Artificial}. MIT
Press.

\bibitem[\citeproctext]{ref-weber1978}
Weber, Max. 1978. \emph{Economy and Society: An Outline of Interpretive
Sociology}. Edited by Guenther Roth and Claus Wittich. University of
California Press.

\end{CSLReferences}


\backmatter


\end{document}
